\section{Loop integration}
The calculation proceeds from loop parameters to elastic angles, then to
soft matching and real-emission energies. All boundary constants are
fixed by the original integrals and their causal boundary values.

\subsection{Analytic integration of the radial loop parameter}
\label{sec:comp-supp-integrating}

We now integrate the preceding amplitudes independently of the
total cross section of Lee, Schwartz and Zhang
\cite{qed-processes-lsz}. In this section all momenta and
photon masses are divided by the electron mass, so $m=1$.
The final correction is restored as
$\sigma_1=\alpha_{\mathrm T}^3 C_{\mathrm{total}}(s)/m^2$, with the argument
$s$ here denoting the dimensionless ratio $s_{\mathrm{phys}}/m^2$.
Likewise $t$, $u$, and $\lambda$ in this chapter stand for
$t_{\mathrm{phys}}/m^2$, $u_{\mathrm{phys}}/m^2$, and
$\lambda_{\mathrm{phys}}/m$. Thus $s+t+u=2$ in elastic kinematics.
The incident-energy variables $X$ and $H$ used below both equal $s-1$.
Integration coordinates with a lowercase letter are defined locally
and must not be identified with these external invariants.
Remove the common factor $\alpha_{\mathrm T}/(4\pi)$ from the loop
kernel and denote the resulting matrix by $\widehat R_1$.
The notation $\mathcal C$ with a sector subscript denotes that sector's
contribution to $C_{\mathrm{total}}$. The complete coefficient satisfies
$C_{\mathrm{total}}(s)=C_1(s-1)$.
Throughout this chapter the interference trace is
$\mathcal T=\mathcal T(\widehat R_1,R_0)$, using
\eqref{eq:comp-supp-trace-bilinear} with $m=1$.

For the spanning graph write the three electron parameters as
\begin{equation}
 a_1=\rho(1-b)z,\qquad a_2=\rho b,\qquad
 a_3=\rho(1-b)(1-z),\qquad a_0=1-\rho,
 \quad 0\leq\rho,b,z\leq1.
\end{equation}
The simplex measure becomes $\rho^2(1-b)\dd\rho\dd b\dd z$.
Let $Q$ be the internal electron momentum and define
\begin{align}
 D&=Q^2-1,\qquad \overline K=bQ+(1-b)[zp+(1-z)p'],\\
 M_z&=1-tz(1-z),\qquad C=1-t(1-b)^2z(1-z),\\
 A&=\overline K^2=C+bD,\qquad B=-bD,\qquad
 \Delta=\rho(A\rho+B)+(1-\rho)\lambda^2.
\end{align}
Here $t=(p-p')^2\leq0$, and $C=A+B>0$.
This last inequality makes the analytic continuation particularly
simple. The internal $s$ ordering has $D>0$, while the $u$ ordering
has $D<0$.

Expand the numerator and its second derivative in $\rho$,
\begin{equation}
 N(\rho\overline K)=\sum_{n=0}^3\rho^n N_n,
 \qquad (\Box_\ell N)(\rho\overline K)=L_0+\rho L_1.
\end{equation}
These coefficients are fixed without any integration by
\begin{equation}
 N_n=\frac{(\overline K\cdot\partial_\ell)^n}{n!}N(\ell)\big|_{\ell=0},
 \qquad
 L_n=\frac{(\overline K\cdot\partial_\ell)^n}{n!}
       \Box_\ell N(\ell)\big|_{\ell=0}\quad(n=0,1).
\end{equation}
Only the $N_0$ term is infrared divergent. All the other terms may
be integrated with $\lambda=0$. Their radial integrals are
\begin{align}
 F_n(A,B)&=\int_0^1\frac{\rho^n\dd\rho}{A\rho+B-\ii0},\qquad
 R_n(A,B)=\int_0^1\frac{\rho^n\dd\rho}{(A\rho+B-\ii0)^2},\\
 F_0&=\frac{\ln C-\ln(B-\ii0)}A,\qquad
 F_n=\frac1{nA}-\frac BA F_{n-1}\quad(n\geq1),\\
 R_0&=\frac1{BC},\qquad
 R_n=\frac{F_{n-1}-BR_{n-1}}A\quad(n\geq1).
 \label{eq:comp-supp-radial-moments}
\end{align}
Polynomial division gives the recurrence for $F_n$, and the identity
$\rho=[(A\rho+B)-B]/A$ gives that for $R_n$.
For $B<0$, use $\ln(B-\ii0)=\ln|B|-\ii\pi$.
Keeping this term is essential for the physical cut.
For $|A/B|<1$, expanding the denominator as a geometric series gives
the following convergent representations. They also avoid cancellation
when the recurrence divides a small difference by $A$,
\begin{equation}
 F_n=\frac1B\sum_{j=0}^\infty\frac{(-A/B)^j}{n+j+1},\qquad
 R_n=\frac1{B^2}\sum_{j=0}^\infty
                    \frac{(j+1)(-A/B)^j}{n+j+1}.
\end{equation}

\subsection{Evaluation of the infrared scalar box}
The scalar coefficient of $N_0$ is
\begin{equation}
 D_0=\int_0^1\dd z\int_0^1\dd b\int_0^1\dd\rho\,
 \frac{\rho^2(1-b)}{[\rho(A\rho+B)+(1-\rho)\lambda^2-\ii0]^2}.
\end{equation}
Its infrared endpoint is $b=\rho=0$. The Feynman boundary value
means that $-\ii0$ is first replaced by $-\ii\epsilon$, with
$\epsilon\downarrow0$ at each fixed $\lambda>0$. We next take
$\lambda\to0$ at fixed nonzero $D$ and fixed $t$. A direct estimate
on the physical side is useful because an asymptotic expansion in the
region $D<0$ need not remain uniform on a cut.

For $D>0$, deform the $b$ contour to $\Gamma$, parametrized by
\begin{equation}
 b(v)=v+\ii\eta v(1-v),\qquad 0\leq v\leq1,
 \qquad 0<\eta\ll1.
\end{equation}
Here $\eta$ is fixed independently of $\lambda$. Along this contour,
with $h=M_z-1\geq0$, the exact denominator obeys
\begin{equation}
 \operatorname{Im}\Delta
 =-\eta v(1-v)\rho
       [D(1-\rho)+2h\rho(1-v)]\leq0.
\end{equation}
Thus the deformation crosses no displaced pole. At $\rho=1$ and
$h=0$ the denominator is $1$, so the vanishing imaginary part there
causes no singularity. The undeformed endpoint $b=1$ is understood
by the same Feynman limit. For $D<0$ the real contour suffices, and
we set $\eta=0$. Choose a temporary boundary $v=\delta\ll1$ and
write $b_*=b(\delta)$. The two contour segments define $D_0^{<}$
and $D_0^{>}$.

On the segment $v\geq\delta$, the radial integral has the uniform
small-mass limit $R_0=1/(BC)$. To see this directly for $D>0$,
factor its quadratic denominator. Its roots are
\begin{equation}
 \rho_-=\frac{\lambda^2}{Db}+O_\delta(\lambda^4),\qquad
 \rho_+=\frac{Db}{A}+O_\delta(\lambda^2).
\end{equation}
The roots remain separated, and $b$, $A$, and $C=A-Db$ stay away
from zero for sufficiently small fixed $\eta$. In particular,
$\rho_+$ stays away from both radial endpoints, uniformly in $z$
and $v\geq\delta$. In the partial fractions, the double-pole and
simple-pole coefficients at $\rho_-$ are respectively $O_\delta(\lambda^4)$
and $O_\delta(\lambda^2)$. Their integrals are therefore
$O_\delta(\lambda^2)$ and $O_\delta(\lambda^2|\ln\lambda|)$.
The coefficients at $\rho_+$ change by $O_\delta(\lambda^2)$, and
their endpoint primitives remain bounded on the chosen branch.
These statements also hold as $v\to1$, where a displaced pole approaches
the real radial interval but never either of its endpoints.
They are uniform at $h=0$, including $z=0,1$ and $t=0$.
Thus the radial error is $O_\delta(\lambda^2|\ln\lambda|)$.
At zero mass the primitive is
$-1/[A(A\rho+B-\ii0)]$, which gives $1/(BC)$ also when $B<0$.
For $D<0$ no radial pole occurs, and the same limit follows directly.
Consequently,
\begin{equation}
 D_0^{>}=-\frac1D\int_0^1\dd z\int_{\Gamma^{>}}\dd b\,
 \frac{1-b}{b[1+(M_z-1)(1-b)^2]}+o_\lambda(1).
\end{equation}

On $0\leq v\leq\delta$, put
$G=\rho^2+v\rho+\lambda^2$ and
$F=M_z\rho^2-Db\rho+\lambda^2$.
The fixed contour angle gives $|F|\geq cG$, with $c>0$ uniform in
$z$, $v$, $\rho$, and small $\lambda$. For $D<0$ positivity gives
the same bound. Moreover,
\begin{equation}
 \Delta-F=O(v\rho^2+\rho\lambda^2),\qquad
 |\Delta|\geq c'G.
\end{equation}
The second inequality follows for small $\delta$ and $\lambda$,
since $|\Delta-F|/G=O(v+\lambda)$.
Replacing $1-b$ by $1$ and $\Delta$ by $F$ consequently changes
this segment by at most a constant times
\begin{equation}
 \int_0^\delta\dd v\int_0^1\dd\rho
 \left[\frac{v\rho^2}{G^2}
       +\frac{\lambda^2\rho^3}{G^3}\right]
 =O(\delta)+O(\lambda).
\end{equation}
For the first term use $G\geq\rho(\rho+v)$ and integrate $\rho$.
For the second integrate $v$ first, leaving a bound proportional
to $\lambda^2\int_0^1\rho^2(\rho^2+\lambda^2)^{-2}\dd\rho=O(\lambda)$.
This controls the endpoint approximation on the physical branch,
including its subsequent $z$ integral.

The model denominator is linear in $b$, so its contour integral gives
\begin{equation}
 D_0^{<}=\frac1D\int_0^1\dd z\int_0^1\rho\,\dd\rho
 \left[\frac1{M_z\rho^2-Db_*\rho+\lambda^2}
       -\frac1{M_z\rho^2+\lambda^2}\right]
       +O(\delta)+o_\lambda(1).
\end{equation}
At fixed $b_*$ the first radial integral tends to
\begin{equation}
 \frac1{M_z}\bigl[\ln(M_z-Db_*)-\ln(-Db_*)\bigr],
\end{equation}
whereas the second is exactly
$(2M_z)^{-1}\ln[(M_z+\lambda^2)/\lambda^2]$.
For $D>0$, $-Db_*$ lies below the negative real axis and
$\ln(-Db_*)=\ln(-D-\ii0)+\ln b_*$.
This fixes the physical branch directly, without continuing a remainder
from $D<0$. Expanding only after the fixed-$\delta$ mass limit gives
\begin{equation}
 D_0^{<}=\frac1D\int_0^1\frac{\dd z}{M_z}
 \left[\ln\lambda+\frac12\ln M_z-\ln b_*
                       -\ln(-D-\ii0)\right]
       +O(\delta)+o_\lambda(1).
\end{equation}
The regular part of the outer contour integral can be deformed back
to the real interval. Its singular part is $-\ln b_*/M_z$,
so the complex cutoff cancels just as a real cutoff would.
The finite part of the remaining $b$ integral is elementary.
With $h=M_z-1\geq0$ and $r=1-b$,
\begin{align}
 \int_0^1\dd b\left[\frac{1-b}{b[1+h(1-b)^2]}
                       -\frac1{bM_z}\right]
 &=\frac1{M_z}\int_0^1\frac{hr-1}{1+hr^2}\dd r\\
 &=\frac1{M_z}\left[\frac12\ln M_z
                  -\frac{\arctan\sqrt h}{\sqrt h}\right].
\end{align}
Adding the two regions cancels $b_*$. After the mass limit the
remaining error is bounded by $O(\delta)$, and therefore vanishes
as $\delta\to0$. We obtain
\begin{equation}
 \boxed{D_0=\frac1D\left[g_0(t)\ln\frac\lambda{-D-\ii0}
                              +g_a(t)\right]+o(1),}
 \label{eq:comp-supp-box-integrated}
\end{equation}
where the remainder vanishes as $\lambda\to0$ at fixed nonzero
$D$ and fixed $t$, and
\begin{align}
 g_0(t)&=\int_0^1\frac{\dd z}{M_z}
 =\frac{4\operatorname{arsinh}(\sqrt{-t}/2)}{\sqrt{(-t)(4-t)}},\\
 g_a(t)&=\int_0^1\frac{\dd z}{M_z}
             \frac{\arctan\sqrt{M_z-1}}{\sqrt{M_z-1}}.
 \label{eq:comp-supp-box-one-dimensional}
\end{align}
Both functions have value $1$ at $t=0$. The remaining integral is
smooth and real, and can also be performed analytically. Define
\begin{equation}
 a=\frac{\sqrt{-t}}2,\qquad h_a=\sqrt{1+a^2},\qquad
 \kappa=\frac{a}{h_a},\qquad 0\leq\kappa<1.
\end{equation}
Introduce an auxiliary angle $\vartheta$, set $z=(1+\cos\vartheta)/2$, and use reflection symmetry
about $\vartheta=\pi/2$. Next set $\tan\vartheta=\tan\varphi/h_a$.
These substitutions give
\begin{align}
 g_a(t)&=\frac1a\int_0^{\pi/2}
       \frac{\arctan(a\sin\vartheta)}{1+a^2\sin^2\vartheta}\dd\vartheta
       =\frac{J(\kappa)}{a h_a},\\
 J(\kappa)&=\int_0^{\pi/2}\arcsin(\kappa\sin\varphi)\dd\varphi.
\end{align}
Differentiating the last expression gives an elementary integral,
\begin{equation}
 J'(\kappa)=\int_0^1\frac{\dd y}
                  {\sqrt{1-\kappa^2+\kappa^2y^2}}
             =\frac{\operatorname{artanh}\kappa}{\kappa}.
\end{equation}
The boundary condition $J(0)=0$ then fixes the answer in terms
of the dilogarithm $\operatorname{Li}_2$,
\begin{equation}
 \boxed{g_a(t)=\frac{\operatorname{Li}_2(\kappa)
                         -\operatorname{Li}_2(-\kappa)}{2a h_a}.}
 \label{eq:comp-supp-box-dilog}
\end{equation}
The continuous limit at $a=0$ is $1$. All arguments in this formula
are real and lie between $-1$ and $1$, so no branch continuation
is needed for the physical spacelike momentum transfer.
The imaginary part of $D_0$ for $D>0$ is
$\pi g_0(t)/D$. This derivation fixes both the infrared logarithm
and the finite part without importing a Compton correction factor.

The complete spanning matrix, with its common coupling removed, is
therefore
\begin{equation}
 \begin{aligned}
 \widehat W={}&N_0D_0+\int_0^1\dd b\int_0^1\dd z\,(1-b)
 \left[\sum_{n=1}^3N_nR_n-\frac14(L_0F_1+L_1F_2)\right].
 \end{aligned}
 \label{eq:comp-supp-spanning-integrated}
\end{equation}
The logarithmic $b=0$ endpoint of the finite terms is integrable.
The matrices $N_n,L_n$ follow simply by multiplying the three
linear electron numerators given above.

\subsection{Integration of the self-energy and vertex}
For the self-energy let $S=P^2$, and in the moments
\eqref{eq:comp-supp-radial-moments} set $A=S$, $B=1-S$.
Integrating the two logarithms in
$\Delta_\Sigma=v[1-(1-v)S]$ gives, after removing the common
ultraviolet constant,
\begin{align}
 \widehat a(S)&=-\frac12-2SF_1+SF_2,\qquad
 \widehat b(S)=2+4SF_1,\\
 \widehat c_m&=4,\qquad \widehat c_Z=-4-4\ln\lambda,\\
 \widehat\Sigma_{\mathrm R}(P)
 &=\widehat a(S)\slashed P+[\widehat b(S)-4]\mathbf1
                         -(\slashed P-1)\widehat c_Z.
\end{align}
A hat means that $\alpha_{\mathrm T}/(4\pi)$ has been factored out.
The expression is the photon-mass asymptotic value, including all
terms that survive the zero-mass limit.

At the vertex write $x=\rho v$, $y=\rho(1-v)$ and $z=1-\rho$.
Because the external photon has $q^2=0$,
\begin{equation}
 \overline K=vP'+(1-v)P,\qquad
 A=\overline K^2=vP'^2+(1-v)P^2,\qquad B=1-A.
\end{equation}
Expand $N(\rho\overline K)=N_0+\rho N_1+\rho^2N_2$.
Integration by parts gives
$\int_0^1\rho\ln(A\rho+B-\ii0)\dd\rho=-AF_2/2$,
since $A+B=1$. The renormalized vertex becomes the single integral
\begin{equation}
 \widehat V_{\mathrm R}^{\mu}
 =\int_0^1\dd v\left[-\sum_{n=0}^2N_n^{\mu}F_n
                  +\gamma^\mu\left(-\frac32+AF_2\right)\right]
                  +\widehat c_Z\gamma^\mu.
 \label{eq:comp-supp-vertex-integrated}
\end{equation}
The constant $-3/2$ includes the finite terms from the
$d$-dimensional numerator. These integrated matrices are inserted
into the same eight-diagram sum derived earlier.

\subsection{Closed evaluation of the vertex parameter}
Each real-photon vertex has one on-shell electron momentum $P_a$
and one off-shell momentum $P_b$. Put $P_a^2=1$, $P_b^2=S$,
$D=S-1$, and choose $v$ to multiply the off-shell momentum.
The change $v\mapsto1-v$ exchanges the original parameter ordering
when necessary. Then
\begin{equation}
 \overline K=P_a+v(P_b-P_a),\qquad A=1+Dv,\qquad B=-Dv.
\end{equation}
Because the numerator is at most quadratic in $v$, only six
scalar moments are needed,
\begin{equation}
 J_{nk}(S)=\int_0^1 v^k F_n(1+Dv,-Dv)\dd v,
 \qquad (n,k)=(0,0),(1,0),(1,1),(2,0),(2,1),(2,2).
\end{equation}
Changing variable from $v$ to $A$ yields an analytic primitive,
\begin{equation}
 J_{nk}(S)=\frac{P_{nk}(S)-P_{nk}(1)}{(S-1)^{k+1}},
 \qquad P'_{nk}(A)=(A-1)^k F_n(A,1-A).
 \label{eq:comp-supp-vertex-closed-moments}
\end{equation}
The radial recurrences reduce this derivative to rational functions
and rational functions times $\ln(1-A-\ii0)$. For example,
$F_0=-\ln(1-A-\ii0)/A$ integrates to
$\operatorname{Li}_2(A+\ii0)$. Integration by parts handles the
higher inverse powers of $A$.

Write $\ell=\ln(1-A-\ii0)$ and
$d_2=\operatorname{Li}_2(A+\ii0)$. The six primitives are
\begin{align}
 P_{00}(A)&=d_2,\\
 P_{10}(A)&=d_2+(1-A^{-1})\ell,\\
 P_{11}(A)&=2A-2d_2+(-A+A^{-1})\ell-1,\\
 P_{20}(A)&=d_2+\left(\frac32-\frac2A+\frac1{2A^2}\right)\ell
                         +\frac1{2A},\\
 P_{21}(A)&=\frac52A-3d_2
  +\left(-A-\frac32+\frac3A-\frac1{2A^2}\right)\ell-1-\frac1{2A},\\
 P_{22}(A)&=A^2-\frac{15}2A+6d_2
  +\left(-\frac{A^2}2+4A-\frac4A+\frac1{2A^2}\right)\ell
  +4+\frac1{2A}.
\end{align}
Their continuous boundary values are
\begin{align}
 (P_{00},P_{10},P_{11})(1)
   &=\left(\frac{\pi^2}6,\frac{\pi^2}6,1-\frac{\pi^2}3\right),\\
 (P_{20},P_{21},P_{22})(1)
   &=\left(\frac12+\frac{\pi^2}6,1-\frac{\pi^2}2,-2+\pi^2\right).
\end{align}
Although individual terms have inverse powers of $A$, each
primitive is regular at $A=0$. Its values there are respectively
$0,1,-2,7/4,-15/4,31/4$. These limits allow the crossed channel to
pass through $S=0$ without a spurious singularity. For $S>1$, the
logarithm is $\ln(S-1)-\ii\pi$ and the dilogarithm is approached
from the upper half-plane. For $S<1$, both functions are real.
The point $S=1$ is an infrared endpoint and requires the regulator
before taking the limit.

The branch choices follow directly from the loop denominator.
After the momentum integration its quadratic form is
$\Delta-\ii0$. At fixed $0<\rho,v<1$, taking
$S\to S+\ii\eta$, with $\eta>0$, gives
\begin{equation}
 \Delta=\rho[\rho-(S-1)v(1-\rho)],\qquad
 \operatorname{Im}\Delta=-\eta\rho v(1-\rho)<0.
\end{equation}
Thus the continuation approaches the physical cut from the same
side as the original Feynman denominators. In the radial result
this is $A\to A+\ii0$ and $1-A\to1-A-\ii0$. For $A>1$,
the latter quantity approaches the negative real axis from below,
fixing its argument to $-\pi$. Finally,
\begin{equation}
 \frac{\dd}{\dd A}\operatorname{Li}_2(A+\ii0)
      =-\frac{\ln(1-A-\ii0)}{A},\qquad
 \operatorname{Im}\operatorname{Li}_2(A+\ii0)=\pi\ln A
 \quad(A>1).
\end{equation}
The imaginary part follows by integration from $A=1$, where it
vanishes. It is therefore fixed by the causal boundary condition,
not by selecting a sign to match a reference value.

To assemble the vertex, expand the numerator coefficients as
\begin{equation}
 N_1^\mu(v)=U_0^\mu+vU_1^\mu,\qquad
 N_2^\mu(v)=V_0^\mu+vV_1^\mu+v^2V_2^\mu.
\end{equation}
These matrices are obtained by inserting
$\slashed{\overline K}=\slashed P_a+v(\slashed P_b-\slashed P_a)$
in the numerator of \eqref{eq:comp-supp-vertex-integrated}.
The fully integrated vertex is
\begin{equation}
 \begin{aligned}
 \widehat V_{\mathrm R}^\mu={}&-N_0^\mu J_{00}
 -U_0^\mu J_{10}-U_1^\mu J_{11}
 -V_0^\mu J_{20}-V_1^\mu J_{21}-V_2^\mu J_{22}\\
 &+\gamma^\mu\left[-\frac32+J_{20}+D J_{21}+\widehat c_Z\right].
 \end{aligned}
 \label{eq:comp-supp-vertex-closed}
\end{equation}
Thus no vertex parameter integral remains. The counterterm
$\widehat c_Z=-4-4\ln\lambda$ supplies the on-shell subtraction
in the same convention as the self-energy above.

\subsection{Closed evaluation of the spanning matrix}
\label{sec:comp-supp-box-tensor}
The finite two-parameter terms in
\eqref{eq:comp-supp-spanning-integrated} can be integrated without
constructing a separate primitive for each numerator coefficient.
The numerator is a polynomial of degree three in the loop momentum.
Differences of propagator denominators therefore reduce its tensor
integrals to a small set of scalar integrals. We derive this reduction
before listing the scalar answers. All momenta remain in units of the
electron mass.

Write the four denominators as
\begin{equation}
 d_i=(\ell-q_i)^2-m_i^2+\ii0,\qquad
 (q_0,q_1,q_2,q_3)=(0,p,Q,p'),\qquad
 (m_0^2,m_1^2,m_2^2,m_3^2)=(0,1,1,1).
\end{equation}
The photon mass is kept in the infrared-divergent scalar integrals.
Setting it to zero in the algebraic coefficients omits only terms
that vanish as $\lambda\to0$. For an ordered subset $I$ of denominator
indices define
\begin{equation}
 T_I^{\mu_1\cdots\mu_r}
 =\frac{e^{\gamma_E\epsilon}\mu^{2\epsilon}}{\ii\pi^{d/2}}
   \int\dd^d\ell\,
   \frac{\ell^{\mu_1}\cdots\ell^{\mu_r}}
        {\prod_{i\in I}d_i},\qquad d=4-2\epsilon.
 \label{eq:comp-supp-box-tensor-definition}
\end{equation}
The factor $e^{\gamma_E\epsilon}$ fixes the normalization so that
the common pole is $U=1/\epsilon$, without an accompanying
$-\gamma_E$. Since the full box is ultraviolet finite, this common
factor does not change its four-dimensional value. We choose
$\mu=m$ in the dimensionless finite parts below.

\paragraph{Contractions from denominator differences.}
Choose an index $a\in I$ as the origin and put
$L=\ell-q_a$. Let $v_i=q_i-q_a$ for $i\in I\setminus\{a\}$.
Define the Gram matrix, its dual vectors and the transverse projector
by
\begin{equation}
 G_{ij}=v_i\cdot v_j,\qquad
 w_i^\mu=\sum_j(G^{-1})_{ij}v_j^\mu,\qquad
 P^{\mu\nu}=\eta^{\mu\nu}-\sum_iw_i^\mu v_i^\nu.
 \label{eq:comp-supp-box-projector}
\end{equation}
These formulas first apply when the vectors $v_i$ are linearly
independent and $\det G\ne0$. Singular limits are taken in the
combined tensor, not in isolated inverse-Gram coefficients.
Thus $w_i\cdot v_j=\delta_{ij}$ and
$P^\mu{}_{\mu}=d-(|I|-1)$. The two required identities are
\begin{equation}
 2L\cdot v_i=d_a-d_i+f_i,\qquad
 L^2=d_a+m_a^2,\qquad
 f_i=v_i^2+m_a^2-m_i^2.
 \label{eq:comp-supp-box-denominator-differences}
\end{equation}
For the following formulas every tensor uses the numerator $L$.
A pinched tensor uses this same $L$, even when its denominator set
no longer contains $a$. This convention avoids an implicit shift
of the numerator in the reduction.

Suppressing the set $I$, the first identity gives
\begin{equation}
 V_i^{\mu_2\cdots\mu_r}
 \equiv v_{i\mu_1}T^{\mu_1\cdots\mu_r}
 =\frac12\left[
 T_{I\setminus\{a\}}^{\mu_2\cdots\mu_r}
 -T_{I\setminus\{i\}}^{\mu_2\cdots\mu_r}
 +f_iT_I^{\mu_2\cdots\mu_r}\right].
 \label{eq:comp-supp-box-pinch-recurrence}
\end{equation}
The loop rank falls by one on the right-hand side. Define
\begin{equation}
 H_i=V_i,\qquad
 H_{ij}=v_{j\nu}V_i^\nu,\qquad
 H_{ijk}=v_{j\nu}v_{k\rho}V_i^{\nu\rho},\qquad
 n=|I|-1.
\end{equation}
Rotations in the space transverse to all $v_i$ then determine
ranks one through three completely,
\begin{align}
 T^\mu&=\sum_iw_i^\mu H_i,\\
 T^{\mu\nu}&=\sum_{ij}w_i^\mu w_j^\nu H_{ij}
                      +P^{\mu\nu}Z,\\
 Z&=\frac{T_{I\setminus\{a\}}+m_a^2T_I
                   -\sum_{ij}(G^{-1})_{ij}H_{ij}}{d-n},\\
 T^{\mu\nu\rho}&=\sum_{ijk}w_i^\mu w_j^\nu w_k^\rho H_{ijk}
       +\sum_i\left(P^{\mu\nu}w_i^\rho
                   +P^{\mu\rho}w_i^\nu
                   +P^{\nu\rho}w_i^\mu\right)Z_i,\\
 Z_i&=\frac{v_{i\mu}
       (T_{I\setminus\{a\}}^\mu+m_a^2T_I^\mu)
                 -\sum_{jk}(G^{-1})_{jk}H_{ijk}}{d-n}.
 \label{eq:comp-supp-box-rank-three}
\end{align}
The equations for $Z$ and $Z_i$ follow by taking a Lorentz trace
and using the second identity in
\eqref{eq:comp-supp-box-denominator-differences}.
To change the origin in a pinched integral, expand each factor
$L=\widetilde L+(q_b-q_a)$ in its numerator. This is a finite
binomial expansion of tensors of lower or equal rank.
Consequently a rank-three box requires only rank-two triangles,
rank-one bubbles and scalar tadpoles.

Although the original box is ultraviolet finite, some pinched
integrals are not. Their poles must be carried through the
recurrence before taking $d=4$. In particular,
\begin{equation}
 \frac{x_{-1}/\epsilon+x_0}{c-2\epsilon}
 =\frac{x_{-1}}{c\epsilon}+\frac{x_0}{c}
                          +\frac{2x_{-1}}{c^2}
                          +\mathcal O(\epsilon).
 \label{eq:comp-supp-box-rational-finite}
\end{equation}
The last term is needed for the finite box. This also specifies
how the calculation handles the rational terms without assuming
that a four-dimensional tensor reduction is sufficient.

\paragraph{Scalar integrals.}
Let $S=Q^2$, $D=S-1$, $T=-t\geq0$, and define
\begin{equation}
 a=\frac{\sqrt T}{2},\qquad h=\sqrt{1+a^2},\qquad
 r=\operatorname{arsinh}a,\qquad \kappa=\frac ah.
\end{equation}
For clarity, denote scalar integrals by
$A_i=T_{\{i\}}$, $B_{ij}=T_{\{i,j\}}$,
$C_{ijk}=T_{\{i,j,k\}}$ and $D_{0123}=T_{\{0,1,2,3\}}$.
The tadpoles and bubbles are
\begin{align}
 A_0&=0,\qquad A_1=A_2=A_3=U+1,\\
 B_{01}&=B_{03}=U+2,\qquad B_{12}=B_{23}=U,\\
 B_{02}&=U+2+\frac{1-S}{S}\ln(1-S-\ii0),\\
 B_{13}&=U+2-\frac{2h}{a}r.
 \label{eq:comp-supp-box-bubbles}
\end{align}
The continuous values at $S=0$ and $T=0$ are used.
These expressions follow by integrating
\begin{equation}
 B_{ij}=U-\int_0^1\dd x\,
 \ln\!\left[xm_i^2+(1-x)m_j^2
                    -x(1-x)(q_i-q_j)^2-\ii0\right].
\end{equation}
For equal masses the rank-one bubble is exactly
$(q_i+q_j)^\mu B_{ij}/2$. This formula is also valid when
$(q_i-q_j)^2=0$, where a one-vector Gram inverse would be
singular. For unequal masses use
\begin{equation}
 T_{\{i,j\}}^\mu=q_i^\mu B_{ij}
 +(q_j-q_i)^\mu
 \frac{A_j-A_i+[(q_j-q_i)^2+m_i^2-m_j^2]B_{ij}}
      {2(q_j-q_i)^2}.
\end{equation}
For $i=0$, $m_j=1$, the coefficient in the second term approaches
$U/2+1/4$ as $(q_j-q_0)^2\to0$.

The four triangles reduce to three distinct functions,
\begin{align}
 C_{012}=C_{023}
 &=-\frac{\operatorname{Li}_2(S+\ii0)-\pi^2/6}{S-1},\\
 C_{123}&=-\frac{2r^2}{T},\\
 C_{013}&=\frac r{ah}\ln\lambda-\frac12\mathcal H(T),\\
 \mathcal H(T)&=
 \frac{2r\ln(2h)-r^2-\operatorname{Li}_2(-e^{-2r})
                        +\operatorname{Li}_2(-1)}{ah}.
 \label{eq:comp-supp-box-triangles}
\end{align}
The overall minus signs in the finite triangles follow from
$\int\dd^d\ell/(\ii\pi^{d/2})$ with three Feynman denominators.
The first line is the previously integrated vertex moment $-J_{00}$.
For the pure-electron triangle, integrate its simplex denominator
$1+T xz$ to obtain $2r^2/T$ before applying this sign.
The infrared triangle follows by writing the sum of its two
electron parameters as $\rho$,
\begin{align}
 C_{013}&=-\int_0^1\dd z\int_0^1
 \frac{\rho\dd\rho}{\rho^2[1+Tz(1-z)]+(1-\rho)\lambda^2-\ii0}\\
 &=g_0(t)\ln\lambda-\frac12\int_0^1
          \frac{\ln[1+Tz(1-z)]}{1+Tz(1-z)}\dd z+o(1).
\end{align}
To evaluate the last integral, set $2z-1=y$ and then
$\kappa y=\tanh u$. Its upper boundary is $u=r$, giving
\begin{align}
 \mathcal H(T)&=\frac1{ah}
 \left[2r\ln h-2\int_0^r\ln\cosh u\,\dd u\right],\\
 \int_0^r\ln\cosh u\,\dd u
 &=\frac{r^2}{2}-r\ln2
  +\frac{\operatorname{Li}_2(-e^{-2r})-\operatorname{Li}_2(-1)}2.
\end{align}
This proves the last line of
\eqref{eq:comp-supp-box-triangles}. Its boundary values are
$C_{123}\to-1/2$ and $\mathcal H(T)\to0$ at $T=0$.
Finally the scalar box is the result already derived in
\eqref{eq:comp-supp-box-dilog},
\begin{equation}
 D_{0123}=\frac1D\left[
 \frac r{ah}\ln\frac{\lambda}{-D-\ii0}
 +\frac{\operatorname{Li}_2(\kappa)-\operatorname{Li}_2(-\kappa)}{2ah}
 \right]+o(1).
\end{equation}
No parameter integral remains in this list.

\paragraph{Physical boundary values.}
The parameter denominators have the boundary value
$\Delta-\ii0$ inherited from the original propagators.
For $S>1$ this gives
$\ln(1-S-\ii0)=\ln(S-1)-\ii\pi$.
Since $\dd\operatorname{Li}_2(S)/\dd S=-\ln(1-S)/S$,
the corresponding dilogarithm is approached from the upper
half-plane. This fixes its imaginary part by continuation from
$S<1$, without choosing a sign from a numerical comparison.
The functions of $T\geq0$ in the other lines are real.
In particular $-e^{-2r}\in[-1,0)$ and $\kappa\in[0,1)$,
so they do not cross a dilogarithm cut.

\paragraph{Assembly of the matrix.}
Expand the original numerator as
\begin{equation}
 N(\ell)=\sum_{r=0}^3
   \mathcal N_{\mu_1\cdots\mu_r}
           \ell^{\mu_1}\cdots\ell^{\mu_r}
 =\gamma^\alpha(\slashed p'-\slashed\ell+1)
 B(\slashed Q-\slashed\ell+1)A
 (\slashed p-\slashed\ell+1)\gamma_\alpha,
\end{equation}
where $A$ and $B$ are the incident and outgoing photon vertex
matrices in the chosen electron ordering. The complete closed
spanning matrix is then
\begin{equation}
 \boxed{\widehat W=\sum_{r=0}^3
       \mathcal N_{\mu_1\cdots\mu_r}
       \left.T_{\{0,1,2,3\}}^{\mu_1\cdots\mu_r}
       \right|_{\epsilon^0}.}
 \label{eq:comp-supp-box-closed-matrix}
\end{equation}
Equations~\eqref{eq:comp-supp-box-pinch-recurrence}--
\eqref{eq:comp-supp-box-triangles} give every coefficient through
finite algebra and logarithms or dilogarithms.
The ultraviolet poles cancel in each box tensor.
At exactly forward or backward kinematics the three-vector Gram
matrix is singular. The physical answer is the continuous limit
of the closed result at interior scattering angles.
This completes the loop-parameter integrations of the spanning
matrix. The subsequent spin interference and elastic-angle
integration are separate steps toward the total cross section.

\section{Virtual interference and angular integration}
\subsection{Closed interference coefficients for the box}
\label{sec:comp-supp-box-interference}
For the angular integration, introduce positive variables
$x=s-1$ and $y=1-u$, so $t=y-x$. Contract each external photon
pair with $-\eta_{\mu\nu}$ and average over the two initial electron
spins and two initial photon polarizations. Individual diagram
groups do not satisfy the photon Ward identity. The same covariant
polarization contraction must therefore be used for the box,
vertex and self-energy groups before they are added.

Let $\mathcal T_{\mathrm{box},s}$ denote the bilinear between
the $s$-ordered spanning matrix and the full Born matrix, with
$\alpha_{\mathrm T}/(4\pi)$ removed from the former. Define finite
scalar parts $\overline A_1=A_1-U$ and
$\overline B_{ij}=B_{ij}-U$. Direct Clifford traces and the preceding
reduction give
\begin{equation}
 \mathcal T_{\mathrm{box},s}
 =f_A\overline A_1+f_a\overline B_{01}
 +f_b\overline B_{02}+f_t\overline B_{13}
 +f_c C_{012}+f_d C_{013}+f_e C_{123}
 +f_DD_{0123}+f_0.
 \label{eq:comp-supp-box-interference-closed}
\end{equation}
The coefficients are real rational functions,
\begin{align}
 f_A&=\frac{4(xy-x+y)}{xy(x+1)},\\
 f_a&=-\frac{8P_a}{x^2y(x-y+4)},\\
 f_b&=\frac{4P_b}{x^2y(x+1)},\\
 f_t&=-\frac{4(x-y)(xy-x+y)}{xy(x-y+4)},\\
 f_c&=\frac{4(2x^2-2xy-2x+y^2-8)}{y},\\
 f_d&=-\frac{2(2x^3-4x^2y+2x^2+xy^2-6xy-12x+y^3-4y)}{xy},\\
 f_e&=-\frac{2(x-y)(2x^2-2xy-2x+y^2-2y-4)}{xy},\\
 f_D&=\frac{2P_D}{xy},\qquad
 f_0=\frac{4(xy-x+y)}{xy}.
\end{align}
Here the polynomials in the longer coefficients are
\begin{align}
 P_a={}&x^3y-2x^2y^2+6x^2y+xy^3-5xy^2-2xy
                +8x+y^3-2y^2-8y,\\
 P_b={}&3x^3y-x^3-2x^2y^2+3x^2y+4x^2
                -4xy^2-4xy+4x-2y^2-4y,\\
 P_D={}&2x^4-4x^3y+3x^2y^2+2x^2y-24x^2
                -xy^3-2xy^2+24xy-16x-8y^2+16y.
\end{align}
The $u$-ordered coefficient of each corresponding scalar integral
is obtained by the simultaneous replacement
$(x,y)\mapsto(-y,-x)$. In that ordering the off-shell invariant
in $\overline B_{02}$, $C_{012}$ and $D_{0123}$ is $S=1-y$.
This relation follows from exchanging the two photon insertions
in the original numerator. It also follows by independent trace
evaluation of that ordering.

The calculation of these coefficients uses only the Clifford
identity $\{\gamma^\mu,\gamma^\nu\}=2\eta^{\mu\nu}$ and the tensor
recurrences above. For example, the scalar trace numerator is
at most cubic in $\ell$. Each factor of $\ell^2$ cancels $d_0$,
and each contraction of $\ell$ with an external momentum is
reduced using \eqref{eq:comp-supp-box-denominator-differences}.
There are seventeen scalar monomials before this reduction in
each photon ordering. The coefficients multiplying $U$ cancel
exactly. Thus \eqref{eq:comp-supp-box-interference-closed} is an
explicit finite expression, including its infrared logarithm.

\subsection{A rational variable for the elastic angle}
For fixed $x>0$ the physical interval is
$x/(1+x)\leq y\leq x$. The substitution
\begin{equation}
 v=\sqrt{\frac{x-y}{x-y+4}},\qquad
 y=\frac{x-(x+4)v^2}{1-v^2},\qquad
 0\leq v\leq\frac{x}{x+2},\qquad
 -\dd y=\frac{8v}{(1-v^2)^2}\dd v
 \label{eq:comp-supp-box-angle-rationalization}
\end{equation}
removes the square root from every spacelike scalar integral.
In particular $r=\operatorname{artanh}v$ and
$ah=v/(1-v^2)$. This substitution follows directly from the
quadratic $T(T+4)$ in the scalar box, and its upper boundary follows
from the backward-scattering value $T=x^2/(1+x)$.

The generalized polylogarithms (GPLs) used here are defined recursively by
\begin{align}
 G(\varnothing;v)&=1,\\
 G(a_1,\ldots,a_n;v)
 &=\int_0^v\frac{\dd t}{t-a_1}G(a_2,\ldots,a_n;t),\\
 G(\underbrace{0,\ldots,0}_{n};v)&=\frac{\ln^n v}{n!}.
\end{align}
The last line specifies the logarithmically singular lower boundary
rather than an ordinary convergent integral. The ordered list
$(a_1,\ldots,a_n)$ is called a word, its entries are letters, and its
weight is its length $n$. The alphabet is the finite set of letters
needed by a given expression. These terms describe the iterated
integrals, not additional physical assumptions.
For a real letter outside the integration interval,
$G(a;v)=\ln(1-v/a)$. Letters on the path inherit their displacement
from the original Feynman denominators, as specified for the crossed
triangle below.

At a logarithmic endpoint, a regularized value means the constant
term after explicitly separating powers of the endpoint logarithm.
This choice is also called a tangential base point. At the origin
it is fixed by the all-zero formula above and by the vanishing limit
of words containing a nonzero letter. A full convergent combination
is formed before taking an endpoint limit. Define
\begin{equation}
 R(v)=\frac{G(-1;v)-G(1;v)}2,
 \qquad Z=\sqrt{\frac{x}{x+4}}.
\end{equation}
Then the needed logarithms and spacelike functions have the
explicit word representations
\begin{align}
 \ln y&=\ln x+G(Z;v)+G(-Z;v)-G(1;v)-G(-1;v),\\
 g_0(t)&=\frac{1-v^2}{v}R(v),\\
 g_a(t)&=\frac{1-v^2}{2v}\left[G(0,-1;v)-G(0,1;v)\right],\\
 \mathcal H(T)&=\frac{1-v^2}{2v}
 \left[G(1,1;v)-G(1,-1;v)+G(-1,1;v)-G(-1,-1;v)\right],\\
 C_{123}&=-\frac{1-v^2}{2v^2}R(v)^2,\qquad
 \overline B_{13}=2-\frac{2R(v)}v.
\end{align}
Products of words are reduced by interleaving their letters while
preserving the order within each word. This identity follows by
partitioning the product of their ordered integration domains.
For example $G(a)G(b)=G(a,b)+G(b,a)$.
These relations turn the closed box interference into rational
functions of $v$ multiplying words of weight at most two. They
provide the starting point for its remaining elastic-angle integral.

\subsection{Integrating the crossed rational subgroup first}
\label{sec:comp-crossed-first}
Before integrating the crossed box in the rationalized angle, separate
the scalar terms whose coefficients and function arguments are rational
in the original invariant $y$.
Transforming those terms to the angle variable would introduce an auxiliary root $\rho_u=\sqrt{(x-1)/(x+3)}$, whose
apparent threshold at $x=1$ is absent from the original scalar
integral. This root is introduced by applying the spacelike
variable $v$ to a subgroup that is already rational in the
original invariant $y$.
This subgroup can be integrated before that change of variable.
Here $x=s-1>0$ and $x/(1+x)\leq y\leq x$, in the mass units
used for the box calculation.

For this subgroup write $A=\overline A_1$, $B_s=\overline B_{02}$,
and $C_s=C_{012}$, evaluated at the crossed invariant $S=1-y$.
The subscript $s$ in these local scalar names does not label the
photon ordering. Include also the explicit constant in the crossed-box
trace. The scalar functions are
\begin{equation}
 A=1,\qquad B_s=2+\frac{y\ln y}{1-y},\qquad
 C_s=\frac{\operatorname{Li}_2(1-y)-\pi^2/6}{y}.
\end{equation}
Define the two GPLs
\begin{equation}
 F(y)=G(1,0;y)=\frac{\pi^2}{6}-\operatorname{Li}_2(1-y),
 \qquad W(y)=G(0,1,0;y).
\end{equation}
Their differential equations are
$F'(y)=\ln y/(y-1)$ and $W'(y)=F(y)/y$.
In particular, $F$ is real and analytic for every $y>0$.
The apparent pole of its derivative at $y=1$ is removable.
These statements fix the continuation of the GPL combination
without assigning separate imaginary parts at an artificial pole.

The selected crossed-box integrand is
\begin{equation}
 I_{\mathrm{rat}}(x,y)=R(x,y)+Q(x,y)\ln y-D(x,y)F(y),
\end{equation}
where exact partial fractions of the trace coefficients give
\begin{align}
 R={}&\frac{16(x-2)}{y^2}+\frac{28x+12}{x}
       +\frac4{x(y-1)}-\frac{8(2x^2+x-4)}{xy},\\
 Q={}&\frac{8(x-2)}y+\frac{4(x-6)}{x(y-1)}
       -\frac{12x+4}{x}-\frac4{x(y-1)^2},\\
 D={}&-\frac{8y}x+\frac{8x-8}x-\frac{4(x^2-8)}{xy}.
\end{align}
The only apparent double pole is integrated with its accompanying
simple pole. The exact identity
\begin{equation}
 \frac1{y-1}-\frac{\ln y}{(y-1)^2}
 =\frac{\dd}{\dd y}\frac{\ln y}{y-1}+\frac1y
\end{equation}
shows explicitly why no isolated $\ln(1-y)$ is needed.
Integration by parts of the two polynomial coefficients of $F$
then gives the finite primitive
\begin{align}
 P(x,y)={}&-\frac{16(x-2)}y+\frac{(32x+28)y+y^2}{x}\\
 &+\frac{-16x^2-8x+36-(4x+16)y-2y^2}{x}\ln y
       +4(x-2)\ln^2y\\
 &+\frac4x\frac{\ln y}{y-1}
       +\frac4x[y^2+(2-2x)y+3x-9]F(y)\\
 &+\frac{4(x^2-8)}xW(y).
\end{align}
Direct differentiation, using the two displayed GPL differential
equations, gives $\partial_yP=I_{\mathrm{rat}}$ as a rational
identity in the independent symbols $\ln y,F,W$.
The contribution to the unnormalized crossed angular integral is
therefore exactly
\begin{equation}
 J_{u,\mathrm{rat}}(x)=P(x,x)-P\left(x,\frac{x}{1+x}\right).
\end{equation}
Since $\dd c=2s\,\dd y/x^2$ and the interference contributes
$(4s)^{-1}\int\operatorname{Re}\mathcal T\,\dd c$, this subgroup
enters the cross-section coefficient as $J_{u,\mathrm{rat}}/(2x^2)$.
At $x=1$ the upper endpoint quotient is assigned
its analytic value $\ln y/(y-1)=1$, so the same expression covers
the apparent root collision there.

The lower endpoint introduces no new alphabet. Write
$f(x)=x/(1+x)$. The first identity below follows directly from
the logarithm of this quotient. For the other two, differentiation
and the boundary values $F(0)=W(0)=0$ prove
\begin{align}
 \ln f(x)&=G(0;x)-G(-1;x),\\
 F(f(x))&=-G(-1,0;x)+G(-1,-1;x),\\
 W(f(x))&=-G(0,-1,0;x)+G(0,-1,-1;x)\\
        &\quad+G(-1,-1,0;x)-G(-1,-1,-1;x).
\end{align}
Thus this entire subgroup has constant letters $\{-1,0,1\}$.
The remaining crossed masters are
$B_a=\overline B_{01}$, $B_t=\overline B_{13}$,
$C_t=C_{013}$, $C_m=C_{123}$, and $D_b=D_{0123}$.
Their rational trace coefficients have denominators built only
from $x$, $y$, and $x-y+4$. In particular none has a pole at
$y=1$. After the spacelike substitution their words involve
$0,\pm1,\pm\sqrt{x/(x+4)}$ and do not involve $\rho_u$.
Consequently $\rho_u$ disappears within the crossed box itself,
before adding soft or real-radiation contributions. This removes an auxiliary root within the virtual calculation.

There is also no remaining positive interior letter in the angular
remainder. If $V=x/(x+2)$ and $z=\sqrt{x/(x+4)}$, then
\begin{equation}
 z^2-V^2=\frac{4x}{(x+4)(x+2)^2}>0.
\end{equation}
Thus $0<V<z<1$. The remainder can be integrated along the real
$v$ interval without detouring around an artificial crossed root.

To fix the lower boundary in the subsequent function-basis
conversion, the series at positive $x\to0$ follows from
\begin{align}
 F(y)&=\sum_{n=1}^{\infty}y^n
             \left(\frac1{n^2}-\frac{\ln y}{n}\right),\\
 W(y)&=\sum_{n=1}^{\infty}y^n
             \left(\frac2{n^3}-\frac{\ln y}{n^2}\right).
\end{align}
They give
\begin{align}
 J_{u,\mathrm{rat}}(x)
 ={}&x(32\ln x-32)+x^2\left(\frac{92}{3}-24\ln x\right)\\
 &+x^3\left(\frac{88}{3}\ln x-\frac{236}{9}\right)
             +O(x^4\ln x).
\end{align}
In particular this unnormalized subgroup vanishes at the lower
boundary. The series fixes the boundary constant from the derived
integrand itself, without numerical fitting or a literature value.


\subsection{Angular primitives and endpoint assembly}
The preferred construction applies the finite primitive rules below to
the direct box and to the seventeen-word crossed remainder, adding the
rational subgroup boundary obtained above. For an independent algebraic
representation, one may instead transform the entire crossed integrand.
This introduces twenty-nine words. We give its continuation explicitly
so that both representations have a fully specified physical branch.

The crossed real-photon triangle is also a weight-two expression.
Put $\eta^2=(x-1)/(x+3)$ and define coefficients
\begin{equation}
 (c_Z,c_{-Z},c_1,c_{-1})=(1,1,-1,-1),\qquad
 (d_\eta,d_{-\eta},d_1,d_{-1})=(-1,-1,1,1).
\end{equation}
Differentiating its dilogarithm gives
\begin{equation}
 \frac{\dd}{\dd v}\operatorname{Li}_2(1-y(v)+\ii0)
 =\frac{y'(v)}{1-y(v)+\ii0}\ln(y(v)-\ii0),\qquad
 \frac{y'}{1-y}=\sum_{a\in\{\eta,-\eta,1,-1\}}\frac{d_a}{v-a}.
\end{equation}
Integrating each simple kernel therefore gives the explicit result
\begin{equation}
 \begin{aligned}
 \operatorname{Li}_2(1-y(v)+\ii0)-\frac{\pi^2}{6}
 ={}&\operatorname{Li}_2(1-x)-\frac{\pi^2}{6}\\
 &+\sum_a d_a\left[\ln x\,G(a;v)
                +\sum_b c_bG(a,b;v)\right].
 \end{aligned}
 \label{eq:comp-supp-box-crossed-triangle-words}
\end{equation}
The sums have exactly the four entries listed above.
For $x>1$ the root $\eta>0$ lies inside the angular interval.
Indeed, with $v_{\max}=x/(x+2)$,
\begin{equation}
 v_{\max}^2-\eta^2=\frac{4}{(x+2)^2(x+3)}>0,\qquad
 Z^2-v_{\max}^2=\frac{4x}{(x+4)(x+2)^2}>0.
\end{equation}
Thus the only positive auxiliary root crossed before the endpoint
is $\eta$, while the zero $Z$ of $y$ is outside the path.
It is an intermediate zero of $1-y$, not a singularity of the
triangle: $\ln y$ vanishes there. Its displacement follows from
$1-y+\ii0=0$. Since $y'(\eta)<0$, the positive root is below
the real axis and the negative root is above it. More explicitly,
\begin{equation}
 y'(v)=-\frac{8v}{(1-v^2)^2},\qquad
 v_{\rm pole}=\eta+\frac{\ii\varepsilon}{y'(\eta)}
                   +O(\varepsilon^2),\qquad \varepsilon>0.
\end{equation}
Here $\varepsilon$ is the infinitesimal Feynman displacement, not
a soft-energy separator. The GPL terms
in \eqref{eq:comp-supp-box-crossed-triangle-words} use those
boundary values together. Their apparent singularity cancels.
For $x<1$ the two roots are off the real integration interval.
The value at $x=1$ is obtained by continuity of the combined
expression. The reality of the combined result can also be proved
without evaluating any GPL. Throughout the path $y>0$, hence
$1-y<1$ is on the real branch of the dilogarithm. At $y=1$,
$\ln y/(1-y)\to-1$, so its displayed derivative is regular.
That derivative and the real initial value at $v=0$ uniquely fix
a real function on the full interval. Consequently the imaginary
parts introduced by the separated GPL terms must cancel.

\paragraph{Finite integration of the words.}
The remaining rational coefficients need no new special functions.
For a rational function $F(v)$, separate its polynomial and pole
parts,
\begin{equation}
 F(v)=\sum_{k\geq0}p_kv^k
       +\sum_a\sum_{j=1}^{m_a}\frac{r_{aj}}{(v-a)^j}
       =P'(v)+\sum_a\frac{r_{a1}}{v-a},
\end{equation}
where all sums are finite and
\begin{equation}
 P(v)=\sum_{k\geq0}\frac{p_kv^{k+1}}{k+1}
       -\sum_a\sum_{j=2}^{m_a}
                    \frac{r_{aj}}{(j-1)(v-a)^{j-1}}.
\end{equation}
The coefficients $r_{aj}$ are fixed algebraically by derivatives
of $(v-a)^{m_a}F(v)$ at $v=a$. This form is the rational
Hermite reduction, written here as partial fractions.
It gives the following finite primitive for the empty word and
for words of lengths one and two,
\begin{align}
 \Phi_{\varnothing}[F](v)
   &=P(v)+\sum_\alpha r_{\alpha1}G(\alpha;v),\\
 \Phi_a[F](v)
   &=P(v)G(a;v)+\sum_\alpha r_{\alpha1}G(\alpha,a;v)
                            -\Phi_{\varnothing}[P/(v-a)](v),\\
 \Phi_{a,b}[F](v)
   &=P(v)G(a,b;v)+\sum_\alpha r_{\alpha1}G(\alpha,a,b;v)
                            -\Phi_b[P/(v-a)](v).
 \label{eq:comp-supp-box-gpl-primitive}
\end{align}
Each recursive use shortens the word. The displayed formulas
therefore terminate after at most three algebraic steps.
Differentiating them gives
$\Phi_w'[F](v)=F(v)G(w;v)$, since the derivative of a nonempty
word removes its first letter and divides by its kernel denominator.
Thus these formulas describe closed GPL expressions of weight
at most three, not a residual parameter quadrature.

\paragraph{Endpoint expression for the box contribution.}
For each photon ordering $j\in\{s,u\}$, the preceding explicit
scalar formulas uniquely specify the finite set of rational
coefficients $F_{jw}(v;x)$ through
\begin{equation}
 \frac{8v}{(1-v^2)^2}
       \operatorname{Re}\mathcal T_{\mathrm{box},j}(x,y(v))
       =\sum_{|w|\leq2}F_{jw}(v;x)G(w;v).
\end{equation}
There are nine input words in the $s$ ordering and twenty-nine
in the $u$ ordering when products have been interleaved as above.
Apply \eqref{eq:comp-supp-box-gpl-primitive} to these rational
coefficients and define the resulting finite GPL sums by
\begin{equation}
 \mathcal F_j(v;x)=\sum_w\Phi_w[F_{jw}](v).
\end{equation}
All partial-fraction coefficients used here have the explicit
algebraic definition
\begin{equation}
 r_{aj}=\frac1{(m_a-j)!}
 \left.\left(\frac{\dd}{\dd v}\right)^{m_a-j}
                  [(v-a)^{m_a}F(v)]\right|_{v=a}.
\end{equation}
Thus a coefficient need not be expanded into a long quotient of
polynomials and square roots in order to define or evaluate the
closed answer.

The lower boundary is also determined by finite algebra.
Write the rational coefficient of a word in $\mathcal F_j$ as
$f_w(v)$, with pole order $p_w$ at zero. For a word that does not
end in zero, let $g_{w,n}$ be its Taylor coefficient of $v^n$.
The coefficients required at this boundary follow from
\begin{align}
 g_{\varnothing,n}&=\delta_{n0},\qquad g_{aw,0}=0,\\
 g_{aw,n}&=-\frac1n\sum_{k=0}^{n-1}
                   \frac{g_{w,k}}{a^{n-k}}\quad(a\ne0),\\
 g_{0w,n}&=\frac{g_{w,n}}n\quad(n\geq1).
\end{align}
Only coefficients up to $n=p_w$ are needed. The single logarithm
$G(0;v)=\ln v$ has zero constant term. Its divergent coefficient
cancels in the combined primitive, as required by the regular
original integrand at forward scattering. Define
\begin{equation}
 \mathcal L_j(x)=\sum_w\sum_{n=0}^{p_w}
 \frac{g_{w,n}}{(p_w-n)!}
 \left.\left(\frac{\dd}{\dd v}\right)^{p_w-n}
                    [v^{p_w}f_w(v)]\right|_{v=0}.
\end{equation}
This is an explicit finite sum of derivatives of rational functions.
It does not require an integration or a numerical endpoint limit.
The integrated box contribution to
$\sigma_1=\alpha_{\mathrm T}^3\mathcal C/m^2$ is
\begin{equation}
 \boxed{\mathcal C_{\mathrm{box}}(x;\lambda)
 =\frac1{2x^2}\sum_{j=s,u}
 \left[\mathcal F_j\!\left(\frac{x}{x+2};x\right)
                         -\mathcal L_j(x)\right].}
 \label{eq:comp-supp-box-angular-closed}
\end{equation}
The factor $1/(2x^2)$ follows from
$\dd c=2(x+1)\dd y/x^2$ and the interference normalization
$\mathcal C_{\mathrm{virt}}=(4s)^{-1}
\int\dd c\,\operatorname{Re}\mathcal T$.
For the preferred shorter construction let $\mathcal F_{u,\mathrm{rem}}$
be the same finite primitive operation applied only to the remainder
specified in Section~\ref{sec:comp-crossed-first}, and let
$\mathcal L_{u,\mathrm{rem}}$ be its lower boundary computed by the
same finite series rule. The full crossed boundary difference obeys
\begin{equation}
 \begin{aligned}
 \mathcal F_u(v_{\max};x)-\mathcal L_u(x)
 ={}&P(x,x)-P\!\left(x,\frac{x}{1+x}\right)\\
 &+\mathcal F_{u,\mathrm{rem}}(v_{\max};x)
       -\mathcal L_{u,\mathrm{rem}}(x).
 \end{aligned}
\end{equation}
Here $P(x,y)$ is the rational-subgroup primitive derived above.
The equality follows by differentiating the two constructions and using
the same original angular endpoints. It reduces the crossed angular
input from twenty-nine to seventeen words and avoids an internal
$\rho_u$ pole in this construction.

The expression still contains the infrared logarithm. It is combined
with the other virtual diagrams and the real contribution before
removing the photon mass.

\paragraph{Continuation through the apparent crossed-channel pole.}
For $x>1$, the positive roots obey $\eta<v_{\max}<Z$ because
\begin{equation}
 v_{\max}^2-\eta^2=\frac4{(x+2)^2(x+3)}>0,\qquad
 Z^2-v_{\max}^2=\frac{4x}{(x+4)(x+2)^2}>0.
\end{equation}
The original triangle requires the root $\eta-\ii0$.
An equivalent GPL evaluation starts at zero and deforms the path
into the upper half-plane before returning to $v_{\max}$.
In particular,
\begin{equation}
 G(\eta-\ii0;v_{\max})
 =\ln\!\left(\frac{v_{\max}}\eta-1\right)-\ii\pi.
\end{equation}
This follows by writing the weight-one integral as
$\log(v_{\max}-\eta+\ii0)-\log(-\eta+\ii0)$: the
first argument has phase zero and the second has phase $+\pi$.
Discarding this imaginary part term by term would change the
analytic continuation. Instead all words use the same deformed
path. Their imaginary parts cancel in the complete real
crossed-channel interference.

\subsection{Angular integration of the vertex and self-energy terms}
\label{sec:comp-supp-vertex-angle}
The closed vertex matrices can now be interfered with the Born
amplitude and integrated over the elastic scattering angle.
Define $x=s-1$ and $y=1-u=x+t$, with $c$ the center-of-mass
photon scattering cosine. Since $t=-x^2(1-c)/(2s)$, the physical
limits and Jacobian are
\begin{equation}
 y_- =\frac{x}{1+x},\qquad y_+=x,\qquad
 \dd c=\frac{2s}{x^2}\dd y.
\end{equation}
Let $T_s$ and $T_u$ be the spin-averaged interferences from the two
vertex insertions and the self-energy insertion in each ordering,
with the common loop factor removed as above. We contract each
external photon pair with $-\eta_{\mu\nu}$. These diagram subsets
are not separately gauge invariant. Physical transverse polarization
sums can replace these contractions only after including the
spanning graphs, when the full Ward identities apply.

The traces use only the Clifford anticommutation relation and
on-shell momentum identities. They are recorded by
\texttt{compton\_virtual\_trace\_analytic.py}. Both vertex insertions
in a given ordering give the same interference. For the $s$ ordering,
\begin{equation}
 T_s(x,y)=\sum_f f(1+x)
       \left[\frac{a_f(x)}y+b_f(x)+c_f(x)y\right],
\end{equation}
where the symbols $f$ in the first column of
Table~\ref{tab:comp-vt-angle-coefficients} mean the moments and
self-energy coefficients already evaluated above. The entries
$\widehat c_Z$ and $1$ are constants, not functions of $1+x$.
\begin{table}[!htbp]
\centering
\begin{tabular}{c|ccc}
$f$ & $a_f$ & $b_f$ & $c_f$\\
\hline
$J_{00}$ & $16(x^2+2x+4)/x$ & $-8(x^3+2x^2+8x+8)/x^2$ & $0$\\
$J_{10}$ & $-8(x^2+2x+8)/x$ & $8(x^3+x^2+6x+8)/x^2$ & $16/x$\\
$J_{11}$ & $-8x$ & $8(x+2)$ & $-8/x$\\
$J_{20}$ & $-4(2x^2+7x+12)/x$ & $4(2x^2+13x+12)/x^2$ & $-4/x$\\
$J_{21}$ & $4(x-4)$ & $-4(2x^2+x-4)/x$ & $4(x+2)/x$\\
$\widehat a$ & $-4(x+2)^2/x^2$ & $2(3x^2+12x+8)/x^3$ & $2(x+1)/x^2$\\
$\widehat b$ & $-2(x^2+4x+8)/x^2$ & $16(x+1)/x^3$ & $2/x^2$\\
$\widehat c_Z$ & $-2(x+4)/x$ & $2(3x+4)/x^2$ & $2/x$\\
$1$ & $2(7x^2+28x+32)/x^2$ & $-2(x+4)(9x+8)/x^3$ & $-2(3x+4)/x^2$\\
\end{tabular}
\caption{Exact trace coefficients for the two $s$-channel vertex
insertions and the self-energy insertion together.}
\label{tab:comp-vt-angle-coefficients}
\end{table}
There is no $J_{22}$ contribution after the trace. The integral is
therefore elementary once the moments have been evaluated,
\begin{equation}
 I_s=\sum_f f(1+x)\left[a_f\ln(1+x)
       +b_f\frac{x^2}{1+x}
       +\frac{c_f x^2}{2}\left(1-\frac1{(1+x)^2}\right)\right].
 \label{eq:comp-vt-s-angle}
\end{equation}

For the crossed ordering the moment arguments depend on the
integration variable, $S=1-y$. Substituting the explicit moments
and self-energy coefficients reduces $T_u$ to rational functions
multiplying $1$, $\ln y$ and $\operatorname{Li}_2(1-y)$.
Partial fractions and integration by parts give the following
closed primitive. Define the real functions
\begin{align}
 H_{10}(y)&=\frac{\pi^2}6-\operatorname{Li}_2(1-y),\\
 H_{010}(y)&=2\operatorname{Li}_3(y+\ii0)
                        -\ln y\operatorname{Li}_2(y+\ii0).
\end{align}
Their derivatives are $H'_{10}=\ln y/(y-1)$ and
$H'_{010}=H_{10}/y$. For $y>1$, continuation from the upper
half-plane gives imaginary parts $\pi\ln y$ and
$\pi\ln^2y/2$ for the dilogarithm and trilogarithm respectively.
They cancel in $H_{010}$. Thus these expressions are real for all
$y>0$, including when the physical interval crosses $y=1$.

The primitive is
\begin{align}
 F_u(x,y)={}&-\frac{2(x^2+8)}x\ln^2y
       +\frac{8(x+1)}x H_{010}(y)
       +P(x,y)\ln y+Q(x,y)H_{10}(y)+R(x,y),\\
 P(x,y)={}&-\frac2{y-1}-\frac{4(x+8)}y
       +\frac{8x+12}{x}y
       +\frac{2\widehat c_Z(x^2-3x-4)+2x^2+10x}{x},\\
 Q(x,y)={}&-\frac{4x}{y^2}+\frac{16}y
                         +\frac{2x^2-18x+20}{x},\\
 R(x,y)={}&\frac{4(-2\widehat c_Z+x-8)}y
                 +\frac{2\widehat c_Z-12x-16}{x}y.
 \label{eq:comp-vt-u-primitive}
\end{align}
The apparent pole multiplying $\ln y$ at $y=1$ is removable,
since $\ln y/(y-1)\to1$. Direct differentiation recovers the
complete crossed-ordering interference, including its finite
on-shell subtraction. Consequently the angular integral of this
virtual subset is
\begin{equation}
 \boxed{\mathcal C_{\mathrm{vertex}+\Sigma}
 =\frac1{2x^2}\operatorname{Re}
       [I_s+F_u(x,x)-F_u(x,x/(1+x))].}
 \label{eq:comp-vt-integrated}
\end{equation}
This is an analytic contribution to the total cross section, with
no loop-parameter or angular integral left. It still contains
$\widehat c_Z=-4-4\ln\lambda$. Its regulator dependence is canceled
only in the complete sum with the spanning and real contributions.

\subsection{Algebraic cancellation of the infrared logarithm}
The explicit trace coefficients allow the infrared cancellation to
be checked as an identity before any angular integration. Write
$x=s-1$, $y=1-u$, $t=y-x$ and
\begin{equation}
 \mathcal B=2\left[\frac{x}{y}+\frac{y}{x}
 -4\left(\frac1y-\frac1x\right)
 +4\left(\frac1y-\frac1x\right)^2\right],\qquad
 z_e=1+\frac{x-y}{2}.
\end{equation}
This is the Born spin average divided by $e_{\mathrm T}^4$.
For comparison with the soft current, define
\begin{equation}
 b_e=\frac{z_e\operatorname{arcosh}z_e}{\sqrt{z_e^2-1}},
 \qquad b_e\big|_{z_e=1}=1.
\end{equation}
The six vertex and self-energy insertions have total coefficient
$\mathcal B$ multiplying $\widehat c_Z$ in their interference.
Since $\widehat c_Z=-4-4\ln\lambda$, their infrared part is
\begin{equation}
 \left.\mathcal T_{\mathrm{vertex}+\Sigma}\right|_{\ln\lambda}
 =-4\mathcal B\ln\lambda.
\end{equation}
In each spanning graph only the scalar triangle $C_t=C_{013}$ and
scalar box $D_b=D_{0123}$ contain $\ln\lambda$. Their coefficients are respectively
$g_0(t)$ and $g_0(t)/(Q^2-1)$. Let $c_{t,s}$ and $c_{b,s}$ denote
the rational trace coefficients multiplying these two scalar
integrals in the $s$ ordering, and use subscript $u$ for the crossed
ordering. The explicit coefficients derived in the spanning
calculation satisfy the rational identity
\begin{equation}
 c_{t,s}+\frac{c_{b,s}}x+c_{t,u}-\frac{c_{b,u}}y
             =4\mathcal B z_e.
 \label{eq:comp-supp-box-IR-identity}
\end{equation}
The script \nolinkurl{check_compton_infrared_analytic.py} checks
this equation symbolically, giving an identically zero numerator.
Thus the full interference has infrared part
\begin{equation}
 \left.\operatorname{Re}\mathcal T\right|_{\ln\lambda}
 =4\mathcal B[z_e g_0(t)-1]\ln\lambda
 =4\mathcal B(b_e-1)\ln\lambda.
\end{equation}
The last equality follows from
$z_e=\cosh[2\operatorname{arsinh}(\sqrt{-t}/2)]$ and the explicit
formula for $g_0(t)$. It is not an inference from numerical agreement.

After angular integration, the coefficient of $\ln\lambda$ in
$\mathcal C_{\mathrm{virt}}$ is therefore exactly
\begin{equation}
 A(s)=\frac1s\int_{-1}^1\mathcal B(b_e-1)\dd c.
\end{equation}
The soft calculation below evaluates this angular integral in
closed form. Its real-emission matching contribution has coefficient
$-A(s)\ln\lambda$, canceling the virtual term. This fixes the
regulator-dependent part of the total before the remaining finite
real-energy integrations are performed.


\section{Soft subtraction and its normalization}
\subsection{Local subtraction in the real phase space}
Work in the center-of-mass frame, and define
\begin{equation}
 E=\frac{s+1}{2\sqrt s},\qquad q=\frac{s-1}{2\sqrt s},\qquad
 p=(E,0,0,-q),\qquad k=(q,0,0,q).
\end{equation}
Order the outgoing photons by energy. Let $k_1=\omega n_1$ be
the lower-energy photon, $k_2=\omega_2 n_2$, and
$n_i=(1,\widehat{\bm n}_i)$. Restricting to $\omega<\omega_2$
accounts for identical photons without a factor $1/2!$.
Put $c_{12}=\widehat{\bm n}_1\cdot\widehat{\bm n}_2$ and
\begin{align}
 D_\gamma&=\sqrt s-\omega(1-c_{12}),\qquad
 \omega_2=\frac{s-1-2\sqrt s\,\omega}{2D_\gamma},\\
 \omega_*&=\frac{s-1}{2\sqrt s+
                 \sqrt{4s-2(1-c_{12})(s-1)}},\qquad 0<\omega<\omega_*.
\end{align}
The on-shell condition for $p'=p+k-k_1-k_2$ determines $\omega_2$.
The upper limit follows by setting $\omega=\omega_2$.
Integrating the electron momentum delta function gives
\begin{equation}
 \dd\Phi_3=\frac{\omega\omega_2}{8(2\pi)^5D_\gamma}
                \dd\omega\dd\Omega_1\dd\Omega_2.
\end{equation}
Rotational symmetry about the incoming photon integrates the
azimuth of $n_2$ to $2\pi$. Denote its remaining polar cosine by
$c_2$ and the relative azimuth of $n_1$ by $\phi$.

Let $\mathcal R=\overline{|\mathcal M_{\gamma\gamma}|^2}/e_{\mathrm T}^6$
be computed from the six diagrams, and
$\mathcal B=\overline{|\mathcal M_0|^2}/e_{\mathrm T}^4$.
The elastic electron paired with direction $n_2$ is
$p'_B=(E,-q\widehat{\bm n}_2)$. Define
\begin{equation}
 j_n=\frac{p'_B}{p'_B\cdot n_1}-\frac p{p\cdot n_1},\qquad
 \mathcal A_{\mathrm{loc}}=\mathcal B(-j_n^2)\frac q{\sqrt s}.
\end{equation}
Then $(\omega\omega_2/D_\gamma)\mathcal R\to\mathcal A_{\mathrm{loc}}/\omega$.
Subtract this term at every phase-space point, rather than
integrating two separately large soft contributions. With
$u=\omega/\omega_*$, the finite real part is
\begin{equation}
 \begin{aligned}
 \mathcal C_{\mathrm R}=
 \frac1{4\pi(s-1)}\int_{-1}^1\dd c_2\int\dd\Omega_1
 \bigg\{\int_0^1\dd u\left[
 \frac{\omega_*^2u\,\omega_2}{D_\gamma}\mathcal R
          -\frac{\mathcal A_{\mathrm{loc}}}{u}\right]
       +\mathcal A_{\mathrm{loc}}\ln\omega_*\bigg\}.
 \end{aligned}
 \label{eq:comp-supp-real-subtracted}
\end{equation}
The logarithms have dimensionless arguments because energies are
measured in units of $m$. The subtracted integral is finite at
$u=0$. Both hard recoil and all spin-dependent emission terms remain.

\subsection{A finite integral for the total correction}
For the elastic momenta let $z_e=p\cdot p'_B=1-t/2$ and
$b_e=z_e\operatorname{arcosh}(z_e)/\sqrt{z_e^2-1}$.
The soft angular integrals in the center-of-mass frame reduce to
one real parameter integral. Define
\begin{align}
 M_x&=1-tx(1-x),\quad v_x=\sqrt{1-M_x/E^2},\quad \beta=q/E,\\
 K&=z_e\int_0^1\frac{\dd x}{M_x}
              \frac{\operatorname{artanh}v_x}{v_x}
                   -\frac{\operatorname{artanh}\beta}{\beta}.
\end{align}
Let $\varepsilon$ be an auxiliary soft-photon energy boundary in
this frame, with $\lambda\ll\varepsilon$ and $\varepsilon$ small
relative to the hard energies. The massive-photon radial integral is
\begin{equation}
 \int_0^{\sqrt{\varepsilon^2-\lambda^2}}
 \frac{r^2\dd r}{\sqrt{r^2+\lambda^2}(M_xr^2+E^2\lambda^2)}
 =\frac1{M_x}\left[
 \operatorname{arcosh}\frac\varepsilon\lambda
 -\frac{\operatorname{artanh}(v_x\sqrt{1-\lambda^2/\varepsilon^2})}{v_x}
 \right].
\end{equation}
Thus $\delta_{\mathrm s}/\alpha_{\mathrm T}
=2[(b_e-1)\ln(2\varepsilon/\lambda)-K]/\pi$ in the
zero-mass limit. Let
$\mathcal T=\mathcal T(\widehat R_1,R_0)$ be the spin-averaged
bilinear in \eqref{eq:comp-supp-trace-bilinear}, with $m=1$. Combining the soft integral
with the analytic subtraction in
\eqref{eq:comp-supp-real-subtracted} cancels $\varepsilon$ before
the remaining integrations. The result is
\begin{align}
 \mathcal C_{\mathrm{VS}}
 &=\frac1s\int_{-1}^1\dd c_2\left[
 \frac14\operatorname{Re}\mathcal T
 +\mathcal B\left\{(b_e-1)\ln\frac2\lambda-K\right\}\right],\\
 \sigma_1&=\boxed{\frac{\alpha_{\mathrm T}^3}{m^2}
             (\mathcal C_{\mathrm{VS}}+\mathcal C_{\mathrm R}).}
 \label{eq:comp-supp-independent-total}
\end{align}
In the normalization used here, the soft part of the virtual
interference obeys
\begin{equation}
 \left.\operatorname{Re}\mathcal T\right|_{\ln\lambda}
      =4\mathcal B(b_e-1)\ln\lambda.
\end{equation}
It cancels the coefficient of $\ln\lambda$ in
$\mathcal C_{\mathrm{VS}}$ point by point in the elastic angle.
There is no photon-energy acceptance cut or numerical soft separator
in this final expression. Its ingredients are the independently
derived amplitudes and the displayed integrations. The analytic
total in Equation (6) of Lee, Schwartz and Zhang
\cite{qed-processes-lsz} supplies a separate final comparison.

\subsection{The finite soft function in dilogarithms}
\label{sec:comp-supp-soft-analytic}
The soft-photon calculation leaves one finite parameter integral at
fixed elastic scattering angle. It can be integrated before performing
the elastic angular integral. In this section momenta and energies are
in units of the electron mass, so $s>1$. Define
\begin{align}
 E&=\frac{s+1}{2\sqrt{s}},\qquad q=\frac{s-1}{2\sqrt{s}},
 \qquad \beta=\frac qE,\qquad t=-2q^2(1-c),\\
 z_e&=1-\frac t2,\qquad M_x=1-tx(1-x),
 \qquad v_x=\sqrt{1-M_x/E^2}.
\end{align}
Here $c$ is the cosine of the elastic photon scattering angle in the
center-of-mass frame. The function to be integrated is
\begin{equation}
 K=z_e I-\frac{\operatorname{artanh}\beta}{\beta},
 \qquad
 I=\int_0^1\frac{\operatorname{artanh}v_x}{v_xM_x}\,\dd x.
 \label{eq:comp-soft-I-definition}
\end{equation}
The ratio $\operatorname{artanh}v_x/v_x$ equals one at $v_x=0$ by continuity.
Every quantity in this integral is real in the physical angular interval.

\paragraph{Rationalizing the square root.}
For an interior angle introduce
\begin{equation}
 a=\frac{\sqrt{-t}}2,\qquad
 \kappa=\frac{\sqrt{q^2-a^2}}E,\qquad
 0<a<q,\qquad 0<\kappa<\beta<1.
\end{equation}
Symmetry about $x=1/2$ allows the substitution $x=(1+y)/2$ with
$0\leq y\leq1$. Then
\begin{equation}
 M_x=1+a^2-a^2y^2,\qquad
 v_x^2=\kappa^2+\frac{a^2y^2}{E^2},\qquad
 I=\frac1{aE}\int_\kappa^\beta
 \frac{\operatorname{artanh}v}{(1-v^2)\sqrt{v^2-\kappa^2}}\,\dd v.
\end{equation}
Set $v=\kappa\cosh u=\kappa(r+r^{-1})/2$ with $r=e^u$.
The result contains only rational functions and logarithms,
\begin{align}
 I&=\frac1{aE}\int_1^{R}
 \frac{2r}{P_+(r)P_-(r)}
 \ln\frac{P_+(r)}{P_-(r)}\,\dd r,
 \label{eq:comp-soft-rational-integral}\\
 P_+(r)&=\kappa r^2+2r+\kappa,\qquad
 P_-(r)=-\kappa r^2+2r-\kappa,\\
 R&=\frac{q+a}{\sqrt{q^2-a^2}},\qquad Y=R-1.
\end{align}
This transformation selects the positive square root because $y\geq0$
and $u\geq0$. Thus it does not require analytic continuation from another
kinematic region.

\paragraph{A real dilogarithmic primitive.}
Let $d=\sqrt{1-\kappa^2}$ and list the four roots with their signs,
\begin{align}
 (z_1,z_2,z_3,z_4)&=
 \left(-\frac{1+d}{\kappa},-\frac{\kappa}{1+d},
 \frac{\kappa}{1+d},\frac{1+d}{\kappa}\right),\\
 (\epsilon_1,\epsilon_2,\epsilon_3,\epsilon_4)&=(1,1,-1,-1),
 \qquad h_j=z_j-1,\\
 A_j&=\frac1{4-2\kappa^2(z_j^2+1)},\qquad
 L_0=\ln\frac{1+\kappa}{1-\kappa}.
\end{align}
Partial fractions and factorization give
\begin{align}
 \frac{2r}{P_+(r)P_-(r)}&=\sum_{j=1}^4\frac{A_j}{r-z_j},\\
 \ln\frac{P_+(1+y)}{P_-(1+y)}
 &=L_0+\sum_{j=1}^4\epsilon_j\ln(1-y/h_j).
\end{align}
The two negative roots and the smaller positive root have $h_j<0$.
The remaining root satisfies $z_4>R$, because $v(R)=\beta<1$
whereas $v(z_4)=1$. Consequently every factor $1-y/h_j$ is positive
for $0\leq y\leq Y$. This proves that all logarithms in the transformed
integral are real and that no pole lies on the integration path.

Define the real dilogarithm on the whole real axis by
\begin{equation}
 \mathcal L_2(x)=\operatorname{Re}\operatorname{Li}_2(x+\ii0),
 \qquad
 \frac{\dd\mathcal L_2(x)}{\dd x}=-\frac{\ln|1-x|}{x}.
\end{equation}
For $x>1$ it can be evaluated without a complex branch choice as
\begin{equation}
 \mathcal L_2(x)=\frac{\pi^2}{6}-\ln x\ln(x-1)
 -\operatorname{Li}_2(1-x),\qquad x>1.
\end{equation}
This identity follows by differentiating both sides and matching their
continuous limits at $x=1$. For $x\leq1$ the ordinary real dilogarithm
applies. Thus an argument larger than one does not introduce an imaginary
contribution to $K$. For $g$ and $h$ drawn from the $h_j$, put
\begin{align}
 \mathcal J(Y;g,h)
 &=\ln|1-g/h|\ln|1-Y/g|
 -\mathcal L_2\!\left(\frac{Y-g}{h-g}\right)
 +\mathcal L_2\!\left(-\frac{g}{h-g}\right),
 \qquad g\ne h,\\
 \mathcal J(Y;g,g)&=\frac12\ln^2|1-Y/g|.
\end{align}
Differentiation establishes the required primitive directly,
\begin{equation}
 \frac{\partial\mathcal J(Y;g,h)}{\partial Y}
 =\frac{\ln(1-Y/h)}{Y-g},\qquad \mathcal J(0;g,h)=0.
\end{equation}
Indeed, for $g\ne h$ the two logarithms produced by differentiation
combine to $\ln|(h-Y)/h|$. Positivity proved above removes the absolute
value from this last expression. The closed result is therefore
\begin{align}
 W_j&=L_0\ln|1-Y/h_j|
 +\sum_{k=1}^4\epsilon_k\mathcal J(Y;h_j,h_k),\\
 K&=\boxed{\frac{z_e}{aE}\sum_{j=1}^4 A_jW_j
 -\frac{\operatorname{artanh}\beta}{\beta}.}
 \label{eq:comp-soft-K-dilog}
\end{align}
This is a finite sum of logarithms and dilogarithms with explicit
arguments. No parameter quadrature remains in the finite soft function.

\paragraph{Angular endpoints.}
At $t=0$, $M_x=1$ and $v_x=\beta$ throughout the interval, which gives
\begin{equation}
 K(s,0)=0.
\end{equation}
At $t=-4q^2$ the rationalizing variable used above degenerates because
$\kappa=0$. The endpoint can instead be integrated directly. Set
$w=\operatorname{artanh}\beta$. Then
\begin{align}
 I\big|_{t=-4q^2}
 &=\frac1{\beta E^2}\int_0^\beta
 \frac{\operatorname{artanh}v}{v(1-v^2)}\,\dd v
 =\frac{J_\beta}{\beta E^2},\\
 J_\beta&=\frac{\operatorname{Li}_2(\beta)-\operatorname{Li}_2(-\beta)}2
 +w\ln(2\cosh w)-\frac{w^2}2
 -\frac12\operatorname{Li}_2(-e^{-2w})-\frac{\pi^2}{24},\\
 K(s,-4q^2)&=\frac{1+2q^2}{\beta E^2}J_\beta-\frac w\beta.
 \label{eq:comp-soft-backward-endpoint}
\end{align}
To obtain the second line, split $1/[v(1-v^2)]$ into
$1/v+v/(1-v^2)$. The first term integrates to the difference of two
dilogarithms. In the second term $v=\tanh u$ gives
$\int_0^w u\tanh u\,\dd u$, which is integrated by parts using
$\ln\cosh u=u+\ln(1+e^{-2u})-\ln2$.
All dilogarithm arguments in this endpoint expression are less than one.

\subsection{The integrated infrared coefficient}
The coefficient multiplying the soft logarithm can also be integrated
analytically over the elastic angle. Define the spin-averaged Born
squared amplitude with its factor $e_{\mathrm T}^4$ removed by $\mathcal B$.
In terms of $x=s-1$ and $y=1-u=x+t$, its invariant expression is
\begin{equation}
 \mathcal B=2\left[\frac xy+\frac yx
 -4\left(\frac1y-\frac1x\right)
 +4\left(\frac1y-\frac1x\right)^2\right].
\end{equation}
The forward limit is $\mathcal B=4$. At threshold this expression tends
to $2(1+c^2)$, which also checks the normalization. Write the required
angular coefficient as
\begin{equation}
 \mathcal C_{\mathrm{IR}}(s)=\int_{-1}^1
 \mathcal B(s,t)\,[b_e(t)-1]\,\dd c,
 \qquad b_e(t)=\frac{z_e\operatorname{arcosh}z_e}{\sqrt{z_e^2-1}}.
\end{equation}
The rapidity variable $r=\exp(\operatorname{arcosh}z_e)$ removes the
square root in $b_e$. Reversing the integration limits gives
\begin{align}
 t&=2-r-r^{-1},\qquad y=s+1-r-r^{-1},\qquad 1\leq r\leq s,\\
 \mathcal C_{\mathrm{IR}}(s)
 &=\frac{2s}{x^2}\int_1^s
 \left[R_+(r)\ln r-R_-(r)\right]\dd r,\\
 R_\eta(r)&=\mathcal B(s,2-r-r^{-1})(1+\eta r^{-2}),
 \qquad \eta=\pm1.
\end{align}
The last integral is a rational function times a single logarithm.
For completeness all its partial-fraction coefficients can be stated
explicitly. Put
\begin{align}
 A&=x-4-\frac8x,\qquad B_0=\frac2x\left(s+5+\frac4x\right),\\
 z_\pm&=\frac{s+1\pm\sqrt{(s-1)(s+3)}}2,\\
 d_z&=-1+z^{-2},\qquad e_z=-2z^{-3},
 \qquad f_{\eta z}=1+\eta z^{-2}.
\end{align}
Then
\begin{align}
 R_\eta(r)&=\sum_{n=-3}^{1}p_{\eta n}r^n
 +\sum_{z=z_\pm}\left[\frac{c_{\eta z}}{r-z}
 +\frac{d_{\eta z}}{(r-z)^2}\right],\\
 p_{\eta,1}&=-\frac2x,\quad p_{\eta,0}=B_0,
 \quad p_{\eta,-3}=-\frac{2\eta}x,\quad
 p_{\eta,-2}=\eta B_0,\\
 p_{\eta,-1}&=-\frac2x+\eta\left(-\frac2x-2A\right),\\
 d_{\eta z}&=\frac{8f_{\eta z}}{d_z^2},\qquad
 c_{\eta z}=f_{\eta z}\left(\frac{2A}{d_z}
 -\frac{8e_z}{d_z^3}\right)-\frac{16\eta}{z^3d_z^2}.
\end{align}
The coefficients follow by expanding $y$ to second order at each of its
zeros. Both zeros are outside the interval: $0<z_-<1<s<z_+$.
No pole displacement or continuation is required for this angular integral.

Define elementary power primitives and the two pole primitives by
\begin{align}
 U_n(r)&=\frac{r^{n+1}}{n+1},\qquad
 V_n(r)=r^{n+1}\left[\frac{\ln r}{n+1}-\frac1{(n+1)^2}\right],
 \qquad n\ne-1,\\
 U_{-1}(r)&=\ln r,\qquad V_{-1}(r)=\frac12\ln^2r,\\
 H_1(r,z)&=\ln|r-z|,\qquad H_2(r,z)=-\frac1{r-z},\\
 J_1(r,z)&=\ln r\ln|1-r/z|+\mathcal L_2(r/z),\\
 J_2(r,z)&=-\frac{\ln r}{r-z}
 +\frac{\ln|r-z|-\ln r}{z}.
\end{align}
Their derivatives are, respectively, $r^n$, $r^n\ln r$,
$(r-z)^{-j}$, and $(r-z)^{-j}\ln r$. The integrated result is
\begin{align}
 \mathcal F_+(r)&=\sum_{n=-3}^1p_{+,n}V_n(r)
 +\sum_{z=z_\pm}[c_{+z}J_1(r,z)+d_{+z}J_2(r,z)],\\
 \mathcal F_-(r)&=\sum_{n=-3}^1p_{-,n}U_n(r)
 +\sum_{z=z_\pm}[c_{-z}H_1(r,z)+d_{-z}H_2(r,z)],\\
 \Delta\mathcal F_\eta&=\mathcal F_\eta(s)-\mathcal F_\eta(1),\\
 \mathcal C_{\mathrm{IR}}(s)&=\boxed{\frac{2s}{(s-1)^2}
 [\Delta\mathcal F_+-\Delta\mathcal F_-].}
 \label{eq:comp-soft-integrated-ir}
\end{align}
This expression contains no angular integral. The cancellation of the total soft logarithm can therefore be checked
with its coefficient already integrated analytically.

\subsection{Exchanging the finite soft angular integrals}
\label{sec:comp-soft-angular-exchange}
The remaining finite soft angular coefficient is
\begin{equation}
 \mathcal C_K(s)=\int_{-1}^1\mathcal B(s,t)K(s,t)\,\dd c.
\end{equation}
Rather than integrate the individual dilogarithms in
\eqref{eq:comp-soft-K-dilog}, it is useful to exchange the original
parameter and angular integrations. Put $T=-t$, $h=1+T/4$, and
\begin{equation}
 f(M)=\frac{\operatorname{artanh}\sqrt{1-M/E^2}}
 {M\sqrt{1-M/E^2}},\qquad
 I(T)=\frac1{2\sqrt{h-1}}\int_1^h\frac{f(M)}{\sqrt{h-M}}\,\dd M.
\end{equation}
This is obtained from $M=h-(h-1)y^2$ in the original parameter integral.
Since $1\leq h\leq E^2$, interchanging the two integrations yields
\begin{equation}
 \int_0^{4q^2}\mathcal B(T)(1+T/2)I(T)\,\dd T
 =2\int_1^{E^2}f(M)\,\dd M
 \int_M^{E^2}\frac{\mathcal B(h)(2h-1)}
 {\sqrt{(h-1)(h-M)}}\,\dd h.
\end{equation}
All endpoint singularities here are integrable. The substitution
$u=\sqrt{(h-M)/(h-1)}$ makes the inner integral rational. Its upper
limit is $U=\sqrt{E^2-M}/q$. Define
\begin{equation}
 k^2=\frac{1+xU^2}{s},\qquad x=s-1,
 \qquad h=\frac{M-u^2}{1-u^2}.
\end{equation}
The rational function that occurs is
\begin{equation}
 \frac{\mathcal B(h)(2h-1)}{1-u^2}
 =\frac{a_1}{1-u^2}+\frac{a_2}{(1-u^2)^2}
 +\frac{a_3}{(1-u^2)^3}
 +\frac{b_1}{k^2-u^2}+\frac{b_2}{(k^2-u^2)^2},
\end{equation}
with explicit coefficients
\begin{align}
 a_1&=-\frac{x^3-6x^2-16x-8}{x^2},
 &a_2&=\frac{(1-k^2)(x^2+2x+4)}x,\\
 a_3&=-x(1-k^2)^2,
 &b_1&=\frac{x^3-2x^2-16x-8}{x^2},\\
 b_2&=\frac{4(1-k^2)(x+2)}{x^2}.
\end{align}
The elementary primitives needed at the upper limit are
\begin{align}
 F_1(U;k)&=\frac{\operatorname{artanh}(U/k)}k,\\
 F_2(U;k)&=\frac{U}{2k^2(k^2-U^2)}+\frac{F_1(U;k)}{2k^2},\\
 F_3(U;k)&=\frac{U}{4k^2(k^2-U^2)^2}+\frac{3F_2(U;k)}{4k^2}.
\end{align}
They vanish at $U=0$ and differentiate to $(k^2-U^2)^{-n}$ when $k$
is held fixed. In the inner integration it is indeed fixed. Consequently
\begin{align}
 H(U)&=a_1F_1(U;1)+a_2F_2(U;1)+a_3F_3(U;1)
 +b_1F_1(U;k)+b_2F_2(U;k),\\
 \mathcal C_K(s)&=\frac4\beta\int_0^1
 \frac{\operatorname{artanh}(\beta U)}{E^2(1-\beta^2U^2)}
 H(U)\,\dd U-\frac{\operatorname{artanh}\beta}{\beta}
 \mathcal C_0(s),
 \label{eq:comp-soft-K-angular-reduced}\\
 \mathcal C_0(s)&=\int_{-1}^1\mathcal B\,\dd c
 =\frac{4s}{x^2}[P(x)-P(x/s)],\\
 P(y)&=\frac{y^2}{2x}+\left(\frac4x+\frac4{x^2}\right)y
 +\left(x-4-\frac8x\right)\ln y-\frac4y.
\end{align}
The square root $k$ is positive, and
$k^2-U^2=(1-U^2)/s>0$ for $0\leq U<1$. Both inverse hyperbolic tangents
in $H$ are therefore real. This reduction integrates the original elastic
angle analytically before performing the remaining soft parameter integral.

The one remaining square root can itself be rationalized by
\begin{align}
 w&=\sqrt{x}U+\sqrt{1+xU^2},
 &1\leq w&\leq\sqrt{x}+\sqrt{s},\\
 U&=\frac{w-w^{-1}}{2\sqrt{x}},
 &k&=\frac{w+w^{-1}}{2\sqrt{s}}.
\end{align}
Thus \eqref{eq:comp-soft-K-angular-reduced} becomes an integral of
rational functions of $w$ multiplying one or two logarithms of rational
functions. It requires at most weight-three multiple polylogarithms.
The endpoints and all logarithm arguments have been obtained directly
from the physical interval, without selecting branches by comparison
with a numerical answer.

\paragraph{Reduction to a finite polylogarithmic expression.}
It is useful to state an algebraic reduction which avoids printing a
large list of weight-three terms. It terminates after at most two
reductions for the expression below and contains no quadrature.
Set $z=w-1$, $Z=\sqrt{s}+\sqrt{x}-1$, and regard $U$ and $k$ as the
rational functions of $w=1+z$ just obtained. Write
\begin{align}
 H&=h_0+h_1\operatorname{artanh}U
       +h_k\operatorname{artanh}(U/k),\\
 h_0&=\frac{a_2U}{2(1-U^2)}
 +a_3\left[\frac{U}{4(1-U^2)^2}
 +\frac{3U}{8(1-U^2)}\right]
 +\frac{b_2U}{2k^2(k^2-U^2)},\\
 h_1&=a_1+\frac{a_2}2+\frac{3a_3}8,
 \qquad h_k=\frac{b_1}k+\frac{b_2}{2k^3},\\
 R_j(z)&=\frac{4U'(z)h_j}{\beta E^2(1-\beta^2U^2)},
 \qquad j=0,1,k.
\end{align}
Define the positive root
\begin{equation}
 r=\sqrt{1+(\sqrt{x}/\beta)^2}-\sqrt{x}/\beta,\qquad 0<r<1.
\end{equation}
Each $R_j$ is explicitly rational. Its poles in the $w$ plane belong to
\begin{equation}
 w=0,\quad \pm\ii,\quad \pm r,\quad \pm r^{-1},
\end{equation}
The orders at $w=0$ are three for $R_0$ and
four for $R_1$, and $R_k$ has no pole there. At $w=\pm\ii$, $R_0$ has
simple poles and $R_k$ has double poles. The four real poles are simple.
The polynomial parts of $R_0$, $R_1$, and $R_k$ have degrees one, two,
and no polynomial part, respectively. These statements follow by
substituting the displayed rational expressions for $U$ and $k$.
Apparent endpoint poles in $h_0$
cancel because $1-k^2=x(1-U^2)/s$ and $k^2-U^2=(1-U^2)/s$.

\paragraph{Reducing the kernel before integration.}
Put $X=x=s-1$ in this paragraph. An integration by parts removes the
poles at $w=\pm\ii$ before any polylogarithmic primitive is constructed.
It combines the elementary and root-dependent parts of the original
soft kernel, so no integrated cross section is needed.

Keep $X=s-1$, $\beta=X/(X+2)$ and $E^2=(X+2)^2/(4s)$, and define
\begin{equation}
 k^2=\frac{1+XU^2}{s},\qquad d=1-\beta^2U^2,\qquad
 a=\operatorname{artanh}(\beta U),\qquad
 f=\operatorname{artanh}(U/k).
\end{equation}
In this section $d$ denotes this rational function, rather than the
spacetime dimension.  On $0<U<1$, all square roots are positive and
both functions are real.  Separate the elementary part of $H$ as
\begin{align}
 h_0&=\overline h_0+\frac{2(X+2)U}{Xk^2},\\
 \overline h_0&=U\left[\frac{X^2+2X+4}{2s}
 -\frac{X^3}{4s^2}-\frac{3X^3(1-U^2)}{8s^2}\right],\\
 h_k&=\frac{b_1}{k}+\frac{2(X+2)(1-U^2)}{Xsk^3},\qquad
 b_1=\frac{X^3-2X^2-16X-8}{X^2}.
\end{align}
The identities
\begin{equation}
 \frac{\dd f}{\dd U}=\frac1{k(1-U^2)},\qquad
 p(U)=\frac{2(X+2)U(1-U^2)}X
\end{equation}
give the exact integration-by-parts relation
\begin{equation}
 \begin{aligned}
 \frac a d\left[\frac{2(X+2)U}{Xk^2}+h_kf\right]
 ={}&\frac{\dd}{\dd U}\left(\frac{apf}{dk}\right)\\
 &+\frac f k\left[
 a\left(\frac{b_1}d+\frac{16sU^2}{X(X+2)d^2}\right)
 -\frac{2U(1-U^2)}{d^2}\right].
 \end{aligned}
 \label{eq:comp-ck-pole-removal}
\end{equation}
Its boundary term vanishes at both ends.  At $U=0$, $a$ and $f$
vanish linearly.  At $U\to1$, $p=O(1-U)$ while $f$ grows only as
$\ln(1-U)$, and $d\to1-\beta^2>0$.  There is therefore no boundary
constant to supply from a different calculation.

Use the previous rational variable
$U=(w-w^{-1})/(2\sqrt X)$.  The remaining root factor cancels from
its differential,
\begin{equation}
 \frac{\dd U}{k}=\sqrt{\frac sX}\frac{\dd w}{w}.
\end{equation}
Writing $z=w-1$, define four rational kernels
\begin{align}
 \widetilde R_0&=\frac4{\beta E^2}\frac{U'\overline h_0}{d},&
 \widetilde R_1&=\frac4{\beta E^2}\frac{U'h_1}{d},\\
 \widetilde R_k&=\frac4{\beta E^2}\sqrt{\frac sX}\frac1w
 \left[\frac{b_1}d+\frac{16sU^2}{X(X+2)d^2}\right],&
 \widetilde R_f&=-\frac8{\beta E^2}\sqrt{\frac sX}
 \frac{U(1-U^2)}{wd^2}.
\end{align}
Here $U'=\dd U/\dd z$ and
\begin{equation}
 h_1=-\frac{X^3-6X^2-16X-8}{X^2}
 +\frac{(1-k^2)(X^2+2X+4)}{2X}
 -\frac{3X}{8}(1-k^2)^2.
\end{equation}
The logarithms are $L_\beta=\operatorname{artanh}(\beta U)$,
$L_1=\operatorname{artanh}U$, and $L_k=\operatorname{artanh}(U/k)$.
The transformed integrand is
\begin{equation}
 \widetilde R_0 L_\beta+\widetilde R_1L_\beta L_1
 +\widetilde R_kL_\beta L_k+\widetilde R_f L_k.
\end{equation}
Every denominator now factors into powers of $w$ and
\begin{equation}
 Q_w=-Xw^4+(4X^2+18X+16)w^2-X.
\end{equation}
The roots of $Q_w$ are the four real poles
$\pm r,\pm r^{-1}$, where
$r=\sqrt{1+X/\beta^2}-\sqrt X/\beta$.
The factor $1+w^2$ is absent.
Thus the entire GPL expression, including the letters introduced by
partial fractions, has only real letters.


For a finite list of pairs $(a,\ell_a)$ write
$L(z)=\sum_a\ell_aG(a;z)$, where $G$ is the usual multiple
polylogarithm with kernels $1/(z-a)$. The three lists needed here are
\begin{align}
 \mathcal A_\beta&=\tfrac12\{(r-1,+1),(-r^{-1}-1,+1),
 (r^{-1}-1,-1),(-r-1,-1)\},\\
 \mathcal A_1&=\tfrac12\{(W^{-1}-1,+1),(-W-1,+1),
 (W-1,-1),(-W^{-1}-1,-1)\},\\
 \mathcal A_k&=\tfrac12\{(W^{-1}-1,+1),(-W^{-1}-1,+1),
 (W-1,-1),(-W-1,-1)\},\\
 W&=\sqrt{s}+\sqrt{x}.
\end{align}
Here the factor $1/2$ multiplies the weights, not the letters.
Factoring $1\pm\beta U$, $1\pm U$, and $k\pm U$ proves
\begin{equation}
 L_\beta(z)=\operatorname{artanh}(\beta U),\qquad
 L_1(z)=\operatorname{artanh}U,\qquad
 L_k(z)=\operatorname{artanh}(U/k).
\end{equation}
All three vanish at $z=0$, which fixes the constants in the factorized
logarithms. Their real letters lie outside $0<z<Z$.
The two functions with endpoint logarithms are defined by their limits
from $z<Z$.

Here is an explicit algebraic rule for removing the rational factors.
Let $\mathfrak T[R;\boldsymbol a](z)$ be linear in the rational function
$R$. Expand $R$ into its polynomial part and partial fractions. For a pole
$a$ of order $m_a$, the coefficient of $(z-a)^{-n}$ is explicitly
\begin{equation}
 C_{a,n}=\left.\frac1{(m_a-n)!}
 \frac{\dd^{m_a-n}}{\dd z^{m_a-n}}
 [(z-a)^{m_a}R(z)]\right|_{z=a},\qquad 1\leq n\leq m_a.
\end{equation}
The polynomial part is obtained by polynomial division. The known pole
list above makes all these operations finite. For a
nonempty word $\boldsymbol a=(b,\boldsymbol c)$ use
\begin{align}
 \mathfrak T[(z-a)^{-1};\boldsymbol a](z)
 &=G(a,\boldsymbol a;z),\\
 \mathfrak T[(z-a)^{-n};b,\boldsymbol c](z)
 &=-\frac{G(b,\boldsymbol c;z)}{(n-1)(z-a)^{n-1}}
 +\frac1{n-1}\mathfrak T\!\left[
 \frac1{(z-a)^{n-1}(z-b)};\boldsymbol c\right](z),
 \quad n\geq2,\\
 \mathfrak T[z^n;b,\boldsymbol c](z)
 &=\frac{z^{n+1}}{n+1}G(b,\boldsymbol c;z)
 -\frac1{n+1}\mathfrak T\!\left[
 \frac{z^{n+1}}{z-b};\boldsymbol c\right](z),
 \quad n\geq0.
\end{align}
For the empty word, use the elementary rules
\begin{align}
 \mathfrak T[z^n;\varnothing](z)&=\frac{z^{n+1}}{n+1},\\
 \mathfrak T[(z-a)^{-1};\varnothing](z)&=G(a;z),\\
 \mathfrak T[(z-a)^{-n};\varnothing](z)
 &=\frac{(z-a)^{1-n}-(-a)^{1-n}}{1-n},\qquad n\geq2.
\end{align}
Every pole appearing here has $a\ne0$. Differentiation proves that
$\mathfrak T[R;\boldsymbol a]'=R\,G(\boldsymbol a;z)$ and its value
at zero is zero. For repeated poles and polynomial terms the word
length decreases. Simple poles terminate immediately. Hence these
rules specify a finite combination of rational functions and ordinary
multiple polylogarithms, rather than a new unevaluated integral.

The endpoint expression for the finite soft angular coefficient is
\begin{align}
 \mathcal S_0(z)&=\sum_{a\in\mathcal A_\beta}
 \ell_a\mathfrak T[R_0;a](z),\\
 \mathcal S_j(z)&=\sum_{a\in\mathcal A_\beta}
 \sum_{b\in\mathcal A_j}\ell_a\ell_b
 \left\{\mathfrak T[R_j;a,b](z)
 +\mathfrak T[R_j;b,a](z)\right\},\qquad j=1,k,\\
 \mathcal C_K(s)&=\lim_{z\to Z^-}
 [\mathcal S_0(z)+\mathcal S_1(z)+\mathcal S_k(z)]
 -\frac{\operatorname{artanh}\beta}{\beta}\mathcal C_0(s).
 \label{eq:comp-soft-K-angular-GPL}
\end{align}
The two words in $\mathcal S_j$ are the shuffle identity for the product
of two weight-one logarithms. Formula \eqref{eq:comp-soft-K-angular-GPL}
is an algebraically specified weight-three polylogarithmic expression.
An expanded list of its coefficients can be generated using only
partial fractions and the finite rules above. The limit is finite:
the original $U$ integrand has at most an integrable logarithm at $U=1$.

\paragraph{The reduced soft endpoint.}
For an explicit closed endpoint expression, apply the finite
primitive operator $\mathfrak T$ just defined to the four reduced
kernels and the logarithm lists $\mathcal A_\beta,\mathcal A_1,\mathcal A_k$. The notation $\mathfrak T[R;L_aL_b]$ means the finite
sum over their weights after applying
$G(a)G(b)=G(a,b)+G(b,a)$.  Put
\begin{equation}
 \begin{aligned}
 \widetilde\Psi(z)={}&
 \mathfrak T[\widetilde R_0;L_\beta](z)
 +\mathfrak T[\widetilde R_1;L_\beta L_1](z)\\
 &+\mathfrak T[\widetilde R_k;L_\beta L_k](z)
 +\mathfrak T[\widetilde R_f;L_k](z).
 \end{aligned}
\end{equation}
Each primitive is normalized to vanish at $z=0$.  The simplified
answer is
\begin{equation}
 \boxed{\mathcal C_K(s)=
 \lim_{z\to Z^-}\widetilde\Psi(z)
 -\frac{\operatorname{artanh}\beta}{\beta}\mathcal C_0(s),
 \qquad Z=\sqrt s+\sqrt X-1.}
 \label{eq:comp-ck-simplified-endpoint}
\end{equation}
The combined endpoint is finite.  All terms are GPLs of weight at
most three with algebraic endpoint coefficients.  There are two further exact residue cancellations.  If
$P_1$ denotes the primitive of the higher-order pole part of
$\widetilde R_1$ at $w=0$, then
\begin{equation}
 \operatorname*{Res}_{w=0}\widetilde R_1=0,\qquad
 \operatorname*{Res}_{w=0}\left[
 \widetilde R_0-P_1\frac{\dd L_1}{\dd w}\right]=0.
\end{equation}
The first identity removes 32 weight-three words introduced by a
vanishing simple-pole coefficient.  The second removes four
weight-two words.  They follow by substituting the displayed rational
kernels, so no small numerical coefficient is discarded.
After these cancellations, the finite rule produces 190 terms instead
of 306, and removes both imaginary letters.  The same physical interval fixes all logarithms.


\paragraph{The finite endpoint limit.} The convergence of the endpoint
expression follows from its local expansion.
Let $d=1-U$ and let $\mathcal I(U)$ denote the integrand including the
factor $4/\beta$ in \eqref{eq:comp-soft-K-angular-reduced}. Define
\begin{align}
 h_{0*}&=\frac{x^2+2x+4}{2s}-\frac{x^3}{4s^2}+\frac{2(x+2)}x,\\
 B_*&=h_{0*}+2\ln2+\frac{b_1}2\ln s.
\end{align}
The rational coefficients have finite limits at $U=1$, and
$h_1+h_k=4$ there. Also $1-U/k=d/s+O(d^2)$. Hence
\begin{align}
 \mathcal I(1-d)&=\frac{4\operatorname{artanh}\beta}{\beta}
 [-2\ln d+B_*]+O(d\ln d),\\
 \int_{1-d}^1\mathcal I(U)\,\dd U
 &=\frac{4\operatorname{artanh}\beta}{\beta}
 d[-2\ln d+2+B_*]+O(d^2\ln d).
\end{align}
Since $d\ln d\to0$, the omitted tail vanishes. This establishes the
existence of the combined polylogarithmic endpoint limit.

\subsection{Changing the soft energy frame}
A useful independent check is to place the soft-photon energy bound in
the final electron rest frame. Write $z_e=\cosh\eta$ and
$v=\tanh\eta$. In this frame the final electron has energy one, the
initial electron has energy $z_e$, and their interpolated momentum has
\begin{equation}
 E_x=1+x(z_e-1),\qquad
 M_x=1+2(z_e-1)x(1-x),\qquad
 v_x=\frac{x\sqrt{z_e^2-1}}{1+x(z_e-1)}.
\end{equation}
In this paragraph $x$ is a Feynman parameter. The finite cross term
integrates directly after changing from $x$ to $v_x$,
\begin{equation}
 \int_0^1\frac{\operatorname{artanh}v_x}{v_xM_x}\,\dd x
 =\frac{J(v)}{\sqrt{z_e^2-1}},\qquad
 J(v)=\int_0^v\frac{\operatorname{artanh}u}{u(1-u^2)}\,\dd u.
\end{equation}
The integral $J$ is the dilogarithmic expression $J_\beta$ above with
$\beta$ replaced by $v$ and $w$ by $\eta$. Including the two self terms
gives the closed finite function
\begin{equation}
 K_{\mathrm{rest}}(t)=\frac{J(v)}v-\frac{b_e+1}2,
 \qquad b_e=\frac\eta v.
\end{equation}
Thus the difference between a rest-frame energy bound $\sqrt{s}\varepsilon$
and a center-of-mass energy bound $\varepsilon$ is
\begin{equation}
 \frac{\delta_{\mathrm{rest}}-\delta_{\mathrm{CM}}}{\alpha_{\mathrm T}}
 =\frac2\pi\left[(b_e-1)\ln\sqrt{s}+K-K_{\mathrm{rest}}\right].
 \label{eq:comp-soft-frame-change}
\end{equation}
This finite change also follows directly from the angular integral of
the eikonal factor multiplied by $\ln[\sqrt{s}/(p'\cdot n)]$, where
$n=(1,\widehat{\boldsymbol k})$ in the center-of-mass frame.

For example, consider the asymmetric invariant cuts
$2P\cdot k_1>\delta$ and $2p'\cdot k_2>\delta$, with
$P=p+k$ and $\delta=2\sqrt{s}\varepsilon$.
Since $2P\cdot k_1+2p'\cdot k_2=s-1$, their omitted soft regions do not
overlap for sufficiently small $\delta$. The two soft terms have weights
$1/2$ from the identical-photon phase-space factor. Their average differs
from the original center-of-mass soft term by half of
\eqref{eq:comp-soft-frame-change}. In the total coefficient multiplying
$\alpha_{\mathrm T}^3/m^2$, the shift is
\begin{equation}
 \Delta C_{\mathrm{VS}}=
 \frac1{2s}\left[\mathcal C_K-\mathcal C_{K,\mathrm{rest}}
 +\frac{\ln s}{2}\mathcal C_{\mathrm{IR}}\right].
\end{equation}
It cannot be omitted if this asymmetric cut is used.

The rest-frame angular coefficient has a short weight-three expression.
Set $r=e^\eta$ and $z=r-1$. Since
\begin{equation}
 \frac{\dd J}{\dd r}=\ln r\left[-\frac1r+\frac1{r-1}+\frac1{r+1}\right],
 \qquad J(0\text{ in }\eta)=0,
\end{equation}
its exact polylogarithmic form is
\begin{equation}
 J=-G(-1,-1;z)+G(0,-1;z)+G(-2,-1;z).
\end{equation}
Let $\widetilde R_+(z)=R_+(1+z)$, with $R_+$ the rational kernel used
in the integrated infrared coefficient. The algebraic map $\mathfrak T$
then gives
\begin{align}
 \mathcal Q(z)&=-\mathfrak T[\widetilde R_+;-1,-1](z)
 +\mathfrak T[\widetilde R_+;0,-1](z)
 +\mathfrak T[\widetilde R_+;-2,-1](z),\\
 \mathcal C_{K,\mathrm{rest}}(s)
 &=\frac{2s}{(s-1)^2}\mathcal Q(s-1)
 -\frac12\mathcal C_{\mathrm{IR}}(s)-\mathcal C_0(s).
\end{align}
This alternative frame is not needed for the symmetric analytic
regulator described next.

\section{Real-emission integration}
\subsection{Real radiation in invariant energy variables}
\label{sec:comp-supp-real-invariants}
The energy ordering used above is convenient for numerical checks.
For analytic integration it is useful to integrate over both photon
labels and include the identical-particle factor $1/2!$.
All invariants in this section are again measured in units of $m^2$.
In the center-of-mass frame define
\begin{equation}
 H=s-1,\qquad x=2\sqrt{s}\,\omega_1,\qquad
 y=2\sqrt{s}\,\omega_2,\qquad z=x+y-H.
\end{equation}
Momentum conservation and the final electron mass shell give
\begin{equation}
 2k_1\cdot k_2=z,\qquad
 2p'\cdot k_1=H-y,\qquad 2p'\cdot k_2=H-x.
\end{equation}
Since $0\leq1-\widehat{\bm n}_1\cdot\widehat{\bm n}_2\leq2$,
$0\leq z\leq xy/s$. Thus the exact integration region is
\begin{equation}
 0<x<H,\qquad H-x<y<\frac{s(H-x)}{s-x}.
 \label{eq:comp-real-dalitz-region}
\end{equation}
This two-dimensional region in final-state energy invariants is
called the Dalitz region. Its boundaries are rational functions. The limits $x=0$ and $y=0$
are the two soft endpoints, while $z=0$ means that the final photons
are parallel.

To obtain the measure, first integrate the three-momentum delta
function over the final electron momentum. Let
$c_{12}=\widehat{\bm n}_1\cdot\widehat{\bm n}_2$ and let $\phi$
be the azimuth of the second photon around the first. The remaining
energy delta function gives
\begin{align}
 E'&=\sqrt{1+\omega_1^2+\omega_2^2+2\omega_1\omega_2c_{12}},\\
 \dd\Phi_3&=\frac{\omega_1\omega_2}{8(2\pi)^5 E'}
 \delta(\sqrt{s}-\omega_1-\omega_2-E')
 \dd\omega_1\dd\omega_2\dd\Omega_1\dd\phi\dd c_{12}.
\end{align}
The absolute derivative of the delta-function argument with respect
to $c_{12}$ is $\omega_1\omega_2/E'$. Its integral therefore
cancels precisely the energy-dependent prefactor in this measure.
A root exists with $-1\leq c_{12}\leq1$ exactly in the Dalitz region
already derived.

For fixed $x,y$, the relative angle is now fixed. Choose a frame
attached to the two outgoing photon directions, and denote the incoming
photon direction in that frame by $\widehat{\bm n}$. The orientation
measure $\dd\Omega_1\dd\phi$ has total volume $8\pi^2$.
Rotations about the incoming direction leave the unpolarized amplitude
unchanged and integrate to $2\pi$, leaving $\dd\Omega_{\bm n}$.
Finally $\dd\omega_1\dd\omega_2=\dd x\dd y/(4s)$, so
\begin{equation}
 \dd\Phi_3=\frac{\dd x\,\dd y}{128\pi^3s}
                    \frac{\dd\Omega_{\bm n}}{4\pi}.
 \label{eq:comp-real-dalitz-measure}
\end{equation}
For example, its massless limit integrates over a triangle of area
$s^2/2$ and gives $\Phi_3=s/(256\pi^3)$.
With $\mathcal R=\overline{|\mathcal M_{\gamma\gamma}|^2}/e_{\mathrm T}^6$,
the real cross section is consequently
\begin{equation}
 \sigma_{\mathrm R}=\frac{\alpha_{\mathrm T}^3}{8sH m^2}
 \int_0^H\dd x\int_{H-x}^{s(H-x)/(s-x)}\dd y\,
 \langle\mathcal R\rangle_{\bm n},\qquad
 \langle f\rangle_{\bm n}=\int\frac{\dd\Omega_{\bm n}}{4\pi}f.
 \label{eq:comp-real-dalitz-cross-section}
\end{equation}
This formula requires the same infrared subtraction as the earlier
real contribution. Its unsubtracted integral diverges at the soft
endpoints.

\subsection{Analytic direction integrals}
Introduce three positive electron--photon invariants,
\begin{align}
 U&=2p\cdot k_1=\frac{x}{2s}
       [s+1+H\widehat{\bm n}\cdot\widehat{\bm n}_1],\\
 V&=2p\cdot k_2=\frac{y}{2s}
       [s+1+H\widehat{\bm n}\cdot\widehat{\bm n}_2],\\
 W&=2p'\cdot k=U+V-z.
 \label{eq:comp-real-angular-denominators}
\end{align}
The labels $UV$, $UW$, and $VW$ below identify pairs of these
directional denominators. In particular, $UV$ here does not mean
ultraviolet. The six electron-line orderings have denominator products
proportional to
\begin{equation}
 H(H-y),\quad H(H-x),\quad U(H-x),\quad V(H-y),
 \quad UW,\quad VW.
\end{equation}
The signs of the individual propagators follow directly from
$(p-k_i)^2-1=-2p\cdot k_i$ and
$(p'-k)^2-1=-2p'\cdot k$.
Thus all directional denominators of the spin-summed squared
amplitude are powers of $U,V,W$. Its numerator is polynomial in
these invariants, since all remaining scalar products follow from
momentum conservation.

Their affine relation gives a useful partial-fraction identity,
\begin{equation}
 \frac1{UVW}=\frac1z\left(\frac1{UW}+\frac1{VW}-\frac1{UV}\right).
 \label{eq:comp-real-three-poles}
\end{equation}
Repeated application after multiplication by further inverse powers
reduces every term with three denominator types to terms with at
most two. This operation is valid for $z>0$. Any apparent pole at
$z=0$ introduced by this reduction must be combined before taking
that limit.

The basic two-denominator direction integral has an elementary
closed form. Write $A=a_0-\bm a\cdot\widehat{\bm n}$ and
$B=b_0-\bm b\cdot\widehat{\bm n}$, where $a$ and $b$ are future
timelike four-vectors. Define
\begin{equation}
 d=a\cdot b,\qquad \Delta=d^2-a^2b^2.
\end{equation}
The bounds $A\geq a_0-|\bm a|>0$ and
$B\geq b_0-|\bm b|>0$ hold uniformly on the sphere. Their convex
combinations are bounded below by
$\min(a_0-|\bm a|,b_0-|\bm b|)>0$, so the Feynman parameter integral
can be interchanged with the direction average.
Direct polar integration of $[vA+(1-v)B]^{-2}$ gives
$[va+(1-v)b]^{-2}$, where the latter square is a Minkowski square.
Consequently,
\begin{align}
 J(a,b)&=\left\langle\frac1{AB}\right\rangle_{\bm n}
 =\int_0^1\frac{\dd v}{v^2a^2+2v(1-v)d+(1-v)^2b^2}\\
 &=\frac1{2\sqrt\Delta}\ln\frac{d+\sqrt\Delta}{d-\sqrt\Delta}.
 \label{eq:comp-real-angular-master}
\end{align}
For $\Delta>0$, the substitution $u=v/(1-v)$ evaluates the parameter
integral explicitly:
\begin{equation}
 J(a,b)=\int_0^\infty\frac{\dd u}{a^2u^2+2du+b^2}
 =\left[\frac1{2\sqrt\Delta}
 \ln\frac{a^2u+d-\sqrt\Delta}{a^2u+d+\sqrt\Delta}
 \right]_{u=0}^{u=\infty}.
\end{equation}
Future timelike vectors obey $d>0$, $\Delta\geq0$, and
$d>\sqrt\Delta$. Thus both logarithm arguments are positive, and
the primitive tends to zero at infinity. These endpoint values fix
the real logarithm and its constant. At $\Delta=0$ the quadratic
denominator is $(\sqrt{a^2}u+\sqrt{b^2})^2$, whose integral is $1/d$.
The positive bounds on $A$ and $B$ persist in a neighborhood of any
fixed future timelike pair. The integrand and its derivatives of any
finite order are then uniformly bounded on the sphere, which permits
differentiation under the integral.
Holding $\bm a,\bm b$ fixed, all higher denominator powers follow by
differentiation,
\begin{equation}
 \left\langle A^{-r}B^{-t}\right\rangle_{\bm n}
 =\frac{(-\partial_{a_0})^{r-1}(-\partial_{b_0})^{t-1}}
             {(r-1)!(t-1)!}J(a,b),\qquad r,t\geq1.
 \label{eq:comp-real-angular-powers}
\end{equation}
These derivatives are taken before substituting the physical
invariants. They are finite algebraic differentiations of the
logarithm in \eqref{eq:comp-real-angular-master}.
For a single denominator the corresponding formula is
\begin{equation}
 \langle A^{-1}\rangle_{\bm n}
 =\frac1{2|\bm a|}\ln\frac{a_0+|\bm a|}{a_0-|\bm a|},\qquad
 \langle A^{-2}\rangle_{\bm n}=\frac1{a^2}.
\end{equation}
The physical vectors corresponding to $U,V,W$ satisfy
\begin{align}
 a_U^2&=x^2/s,\qquad a_V^2=y^2/s,\qquad a_W^2=H^2/s,\\
 a_U\cdot a_V&=xy/s+H^2z/(2s),\\
 a_U\cdot a_W&=xH/s+zH(H-x)/(2s),\\
 a_V\cdot a_W&=yH/s+zH(H-y)/(2s).
 \label{eq:comp-real-angular-Gram}
\end{align}
In particular they remain timelike throughout the interior of the
Dalitz region. These identities make each scalar angular master a
closed logarithm of $s,x,y$.

The next section carries out this reduction for the full six-diagram
trace and lists the resulting coefficients. The remaining $x,y$
integrations still require the infrared subtraction.

\subsection{Explicit completion of the direction integral}
The preceding reduction can be carried out for the full six-diagram
square. We describe the algebra and give its coefficients so that
the result can be reconstructed without evaluating Dirac matrices.
The ordered photon triples are
\begin{equation}
 (a,b,c)=(0,1,2),(0,2,1),(1,0,2),(1,2,0),(2,0,1),(2,1,0).
\end{equation}
Their signed denominator products are, in this order,
\begin{equation}
 d_i=\bigl(H(H-x),H(H-y),-U(H-x),UW,-V(H-y),VW\bigr)_i.
\end{equation}
Let $T_{ij}$ be the numerator trace for the interference of
orders $i,j$, including the initial average $1/4$ and three
polarization sums $-\eta_{\mu\nu}$. Then
\begin{equation}
 \mathcal R=\sum_{i\leq j}\frac{(2-\delta_{ij})T_{ij}}{d_i d_j}.
\end{equation}
There are 21 independent terms, including all 36 interferences.
The trace reduction follows directly by moving the first contracted
gamma matrix through the intervening string. For
$\Gamma_n=\gamma^{\alpha_1}\cdots\gamma^{\alpha_n}$ in four dimensions,
\begin{equation}
 \gamma_\mu\Gamma_n\gamma^\mu
 =4(-1)^n\Gamma_n+
 2\sum_{j=1}^n(-1)^{j-1}
 \gamma^{\alpha_1}\cdots\widehat{\gamma^{\alpha_j}}\cdots
 \gamma^{\alpha_n}\gamma^{\alpha_j}.
\end{equation}
The hat denotes omission. Apply this relation to the three photon
index pairs, using cyclicity of the trace to bring each pair to
the two ends. The remaining traces contain only slashed momenta
and obey
\begin{equation}
 \operatorname{tr}(\slashed v_1\cdots\slashed v_{2n})
 =\sum_{j=2}^{2n}(-1)^j(v_1\cdot v_j)
 \operatorname{tr}(\slashed v_2\cdots
       \widehat{\slashed v_j}\cdots\slashed v_{2n}),
 \qquad \operatorname{tr}\mathbf1=4.
\end{equation}
All scalar products are fixed by the Gram matrix of $(p,k,k_1,k_2)$,
\begin{equation}
 G=\frac12\begin{pmatrix}
 2&H&U&V\\ H&0&x-U&y-V\\ U&x-U&0&z\\ V&y-V&z&0
 \end{pmatrix}.
\end{equation}
Thus this is a finite polynomial calculation with no integration
or input from a published correction factor.

After applying \eqref{eq:comp-real-three-poles} and cancelling
numerator factors, the result has 25 monomials,
\begin{equation}
 \mathcal R=\sum_{\bm e}C_{\bm e}U^{e_1}V^{e_2}W^{e_3},\qquad
 \boxed{\langle\mathcal R\rangle_{\bm n}
          =\sum_{\bm e}C_{\bm e}\mathfrak M_{\bm e}.}
 \label{eq:comp-real-angle-completed}
\end{equation}
The coefficients are listed below. The moments $\mathfrak M$ in
this formula are the following closed expressions, not unevaluated
angular integrals.

For two denominator types write $A=a_0-\bm a\cdot\widehat{\bm n}$,
$B=b_0-\bm b\cdot\widehat{\bm n}$ and abbreviate
$\alpha=a^2$, $\beta=b^2$, $d=a\cdot b$,
$\Delta=d^2-\alpha\beta$. Here $\alpha,\beta$ are auxiliary
kinematic symbols and are not the electromagnetic coupling.
Let $J$ be the logarithm in \eqref{eq:comp-real-angular-master}.
The required inverse-power moments are
\begin{align}
 \langle A^{-1}B^{-1}\rangle&=J,\\
 \langle A^{-2}B^{-1}\rangle
 &=\frac{(b_0d-a_0\beta)J+a_0d/\alpha-b_0}{\Delta},\\
 \langle A^{-1}B^{-2}\rangle
 &=\frac{(a_0d-b_0\alpha)J+b_0d/\beta-a_0}{\Delta},\\
 \langle A^{-2}B^{-2}\rangle
 &=\frac{Q_JJ+Q_0}{\Delta^2},\\
 Q_J&=2a_0b_0(\alpha\beta+2d^2)
       -d(\Delta+3\alpha b_0^2+3\beta a_0^2),\\
 Q_0&=\Delta+(2\alpha\beta+d^2)
          \left(\frac{a_0^2}{\alpha}+\frac{b_0^2}{\beta}\right)
          -6a_0b_0d.
\end{align}
For positive powers of $A$ and one denominator type $B$, define
\begin{align}
 r^2&=b_0^2-\beta,\qquad
 h=\frac{a_0b_0-d}{r^2},\qquad g=a_0-hb_0,\\
 v_\perp^2&=a_0^2-\alpha-\frac{(a_0b_0-d)^2}{r^2},\qquad
 L_B=\frac1{2r}\ln\frac{b_0+r}{b_0-r},\\
 n_0&=g^2-\frac{\beta v_\perp^2}{2r^2},\qquad
 n_1=2gh+\frac{b_0v_\perp^2}{r^2},\qquad
 n_2=h^2-\frac{v_\perp^2}{2r^2}.
\end{align}
Integration over the azimuth about $\bm b$ first gives the
polynomials $g+hB$ and $n_0+n_1B+n_2B^2$. Hence
\begin{align}
 \langle A/B\rangle&=gL_B+h,\qquad
 \langle A/B^2\rangle=g/\beta+hL_B,\\
 \langle A^2/B\rangle&=n_0L_B+n_1+n_2b_0,\\
 \langle B^{-1}\rangle&=L_B,\qquad
 \langle B^{-2}\rangle=1/\beta,\qquad
 \langle A\rangle=a_0,\qquad\langle1\rangle=1.
\end{align}
These formulas cover every monomial in
\eqref{eq:comp-real-angle-completed}. For the physical substitutions,
\begin{equation}
 (a_{U0},a_{V0},a_{W0})
 =\left(\frac{x(s+1)}{2s},\frac{y(s+1)}{2s},
                    \frac{(s-1)(2s-x-y)}{2s}\right).
\end{equation}
Their squares and scalar products are given in
\eqref{eq:comp-real-angular-Gram}.
When $\Delta$ or $r$ vanishes while the denominator vectors remain
future timelike, the expressions mean their continuous limits.
The same positive lower bounds justify taking these limits under
the direction integral. Apparent inverse powers at these points
are removable. Soft limits in which a denominator vector vanishes
instead require the infrared subtraction.

The direction integration is now complete analytically. The
remaining real contribution is a two-dimensional Dalitz integral
of explicit rational functions and logarithms, with its two soft
endpoints still to be combined with the virtual correction.
% BEGIN EXACT REAL ANGULAR COEFFICIENTS
\subsubsection{Coefficients of the integrated direction moments}
The following coefficients use $y=H-x+z$ to shorten the polynomials.
Every exponent triple not displayed has zero coefficient.
\begin{equation}
C_{-2,0,-2}=-32.
\end{equation}
\begin{equation}
C_{-2,0,-1}=\frac{8 \left(H^{2} - 2 H x - 4 H + x^{2} + 4 x - 8\right)}{- H + x}.
\end{equation}
\begin{equation}
C_{-2,0,0}=- \frac{8 \left(- H z + 4 H + x z - 4 x + 4\right)}{\left(- H + x\right)^{2}}.
\end{equation}
\begin{equation}
C_{-2,1,0}=\frac{8}{- H + x}.
\end{equation}
\begin{equation}
C_{-1,-1,0}=\frac{P_{-1,-1,0}}{z^{2} \left(- H + x\right) \left(x - z\right)}.
\end{equation}
\begin{equation}
\begin{aligned}
P_{-1,-1,0}={}& - 64 x^{2} + 32 z^{3} - 24 x z^{3} - 4 H z^{4} - 4 H^{3} z^{2} + 16 H z^{3} + 16 H^{2} z^{2}\\
& + 24 x^{2} z^{2} + 32 H z^{2} + 32 H^{2} z + 64 H x + 64 x z - 32 H x z^{2} - 8 H^{2} x z + 8 H x^{2} z.
\end{aligned}
\end{equation}
\begin{equation}
C_{-1,0,-2}=\frac{8 \left(z^{2} + 4 z - 8\right)}{z}.
\end{equation}
\begin{equation}
C_{-1,0,-1}=\frac{P_{-1,0,-1}}{H z^{2} \left(- H + x\right) \left(x - z\right)}.
\end{equation}
\begin{equation}
\begin{aligned}
P_{-1,0,-1}={}& - 64 z^{3} + 32 z^{4} - 128 H z^{3} - 64 H^{2} x - 56 H^{2} z^{3} - 32 H^{3} z - 8 x z^{4} - 4 x z^{5}\\
& - 4 H^{2} z^{4} - 4 H^{4} z^{2} - 4 x^{3} z^{3} + 4 H z^{5} + 4 H^{3} z^{3} + 8 x^{2} z^{3} + 32 x z^{3}\\
& + 32 x^{2} z^{2} + 64 H x^{2} + 128 H z^{2} + 128 H^{2} z^{2} - 64 H x z - 64 H x z^{2}\\
& - 56 H x^{2} z^{2} - 12 H^{2} x z^{3} - 12 H^{2} x^{2} z^{2} - 8 H^{2} x^{2} z + 4 H x z^{4}\\
& + 4 H x^{3} z^{2} + 8 H^{3} x z + 12 H x^{2} z^{3} + 12 H^{3} x z^{2} + 48 H^{2} x z^{2} + 64 H x z^{3}.
\end{aligned}
\end{equation}
\begin{equation}
C_{-1,0,0}=\frac{P_{-1,0,0}}{H \left(- H + x\right)^{2} \left(x - z\right)}.
\end{equation}
\begin{equation}
\begin{aligned}
P_{-1,0,0}={}& - 64 z - 16 x^{3} - 8 H^{4} - 4 x^{4} + 32 x^{2} + 64 H + 64 H^{3} + 160 H^{2} - 160 H x - 96 H^{2} x\\
& - 64 H z - 24 H^{2} z - 24 H^{2} x^{2} - 8 H^{2} z^{2} - 4 H^{3} z - 4 x^{2} z^{2} + 16 H x^{3}\\
& + 16 x^{2} z + 20 H^{3} x + 32 x z + 40 H x^{2} + 4 H x^{2} z + 12 H x z^{2} + 16 H x z.
\end{aligned}
\end{equation}
\begin{equation}
C_{-1,1,0}=- \frac{4 \left(2 H^{2} - H x + H z - 4 H + x z + 4 x - 2 z - 8\right)}{H \left(- H + x\right) \left(x - z\right)}.
\end{equation}
\begin{equation}
C_{-1,2,0}=- \frac{4 \left(- 2 H + x\right)}{H \left(- H + x\right) \left(x - z\right)}.
\end{equation}
\begin{equation}
C_{0,-2,-2}=-32.
\end{equation}
\begin{equation}
C_{0,-2,-1}=- \frac{8 \left(x^{2} - 2 x z - 4 x + z^{2} + 4 z - 8\right)}{x - z}.
\end{equation}
\begin{equation}
C_{0,-2,0}=\frac{8 \left(x z - 4 x - z^{2} + 4 z - 4\right)}{\left(x - z\right)^{2}}.
\end{equation}
\begin{equation}
C_{0,-1,-2}=\frac{8 \left(z^{2} + 4 z - 8\right)}{z}.
\end{equation}
\begin{equation}
C_{0,-1,-1}=\frac{P_{0,-1,-1}}{H z^{2} \left(- H + x\right) \left(x - z\right)}.
\end{equation}
\begin{equation}
\begin{aligned}
P_{0,-1,-1}={}& - 64 z^{3} - 8 z^{6} + 96 z^{4} - 96 H z^{3} - 96 x z^{3} - 64 H^{2} x - 56 H^{2} z^{3} - 32 H^{3} z\\
& - 16 H^{3} z^{2} - 12 x^{2} z^{4} - 8 x z^{4} + 4 x^{3} z^{3} + 8 H z^{5} + 8 x^{2} z^{3} + 16 H z^{4}\\
& + 16 x z^{5} + 32 x^{2} z^{2} + 64 H x^{2} + 96 H^{2} z^{2} + 128 H z^{2} - 64 H x z - 56 H x^{2} z^{2}\\
& - 16 H x z^{4} - 8 H^{2} x^{2} z - 4 H x^{3} z^{2} + 8 H^{3} x z + 12 H x^{2} z^{3} + 32 H x z^{3}\\
& + 80 H^{2} x z^{2}.
\end{aligned}
\end{equation}
\begin{equation}
C_{0,-1,0}=\frac{P_{0,-1,0}}{H \left(- H + x\right) \left(x - z\right)^{2}}.
\end{equation}
\begin{equation}
\begin{aligned}
P_{0,-1,0}={}& - 64 H - 64 z^{2} - 32 H^{2} - 32 x^{2} - 16 x^{3} + 4 x^{4} + 8 H^{3} + 8 z^{4} + 64 z - 96 H x\\
& - 64 H^{2} x - 40 H z^{2} - 24 x z^{3} - 16 x z^{2} - 16 x^{3} z - 8 H z^{3} - 8 H^{2} z^{2} - 4 H^{3} z\\
& + 4 H^{3} x + 8 H x^{2} + 28 x^{2} z^{2} + 32 x^{2} z + 56 H^{2} z + 96 x z + 128 H z - 4 H x^{2} z\\
& + 8 H^{2} x z + 12 H x z^{2} + 32 H x z.
\end{aligned}
\end{equation}
\begin{equation}
C_{0,0,-2}=-16.
\end{equation}
\begin{equation}
C_{0,0,-1}=\frac{P_{0,0,-1}}{H \left(- H + x\right) \left(x - z\right)}.
\end{equation}
\begin{equation}
\begin{aligned}
P_{0,0,-1}={}& - 32 z - 8 z^{2} + 4 z^{3} + 16 H^{2} + 32 H + 32 x^{2} - 32 H x - 32 x z - 12 H z^{2} + 4 x z^{2}\\
& + 8 H^{2} z + 24 H z - 4 H x z.
\end{aligned}
\end{equation}
\begin{equation}
C_{0,0,0}=\frac{P_{0,0,0}}{H^{2} \left(- H + x\right)^{2} \left(x - z\right)^{2}}.
\end{equation}
\begin{equation}
\begin{aligned}
P_{0,0,0}={}& - 32 H^{2} - 32 H^{3} - 32 z^{2} - 16 x^{4} - 72 H^{2} z^{2} - 64 H z^{2} - 48 H^{3} x - 32 H^{2} x\\
& - 32 x^{2} z - 24 x^{2} z^{2} - 20 H^{3} z^{2} - 8 H z^{3} - 4 H^{4} z - 4 H^{2} x^{3} + 4 H^{4} x\\
& + 8 x z^{3} + 32 H x^{2} + 32 H x^{3} + 32 x z^{2} + 32 x^{3} z + 32 H^{2} x^{2} + 48 H^{3} z + 64 H z\\
& + 96 H^{2} z - 120 H x^{2} z - 20 H^{2} x^{2} z - 4 H x^{2} z^{2} + 4 H x^{3} z + 20 H^{3} x z\\
& + 24 H^{2} x z^{2} + 40 H^{2} x z + 96 H x z^{2}.
\end{aligned}
\end{equation}
\begin{equation}
C_{0,1,0}=\frac{4 \left(3 H x - 3 H z + 2 H - 2 z\right)}{H \left(- H + x\right) \left(x - z\right)^{2}}.
\end{equation}
\begin{equation}
C_{1,-2,0}=- \frac{8}{x - z}.
\end{equation}
\begin{equation}
C_{1,-1,0}=\frac{4 \left(- H^{2} - H x - H z + x z + 4 x - z^{2} - 2 z + 8\right)}{H \left(- H + x\right) \left(x - z\right)}.
\end{equation}
\begin{equation}
C_{1,0,-1}=- \frac{4 \left(- H + z\right) \left(- H + 2 x - z\right)}{H \left(- H + x\right) \left(x - z\right)}.
\end{equation}
\begin{equation}
C_{1,0,0}=\frac{4 \left(- 3 H^{2} + 3 H x - 2 H + 2 z\right)}{H \left(- H + x\right)^{2} \left(x - z\right)}.
\end{equation}
\begin{equation}
C_{2,-1,0}=\frac{4 \left(H + x - z\right)}{H \left(- H + x\right) \left(x - z\right)}.
\end{equation}
% END EXACT REAL ANGULAR COEFFICIENTS

\subsection{Rationalization of the remaining logarithms}
After the completed angle average is collected, its functional
basis is particularly small,
\begin{equation}
 \langle\mathcal R\rangle=R_0+R_U L_U+R_V L_V+R_W L_W
                  +R_{UV}J_{UV}+R_{UW}J_{UW}+R_{VW}J_{VW}.
\end{equation}
Each $R$ is rational in $H,x,z$. The elementary single-denominator
functions obey $L_U=s\ln s/(xH)$ and $L_V=s\ln s/(yH)$.
The coefficients follow by inserting the moment formulas above in
\eqref{eq:comp-real-angle-completed}. The square roots in the other
four logarithms can each be rationalized before the remaining
energy integrals. This operation is performed separately in each
logarithmic term, with a common original integration domain.

For the $UV$ term, the angular Gram determinant is
\begin{equation}
 \Delta_{UV}=\frac{H^2z(4xy+H^2z)}{4s^2}.
\end{equation}
Its logarithm is a hyperbolic angle. Defining its positive
exponential $v$ fixes the substitution uniquely by $v\geq1$,
\begin{align}
 \frac{H^2z}{xy}&=\frac{(v-1)^2}{v},\qquad
 D_v=H^2v-x(v-1)^2,\\
 z&=\frac{x(H-x)(v-1)^2}{D_v},\qquad
 y=\frac{H^2v(H-x)}{D_v},\\
 \frac{\partial z}{\partial v}
 &=\frac{x(H-x)H^2(v^2-1)}{D_v^2},\qquad
 J_{UV}=\frac{2sv\ln v}{xy(v^2-1)}.
\end{align}
The exact upper Dalitz boundary gives $v=s$. Therefore
\begin{equation}
 0<x<H,\quad1<v<s,\qquad
 J_{UV}\dd z=\frac{2s\ln v}{D_v}\dd v.
 \label{eq:comp-real-UV-rational}
\end{equation}
No square root remains. At fixed $v$, the coefficient multiplying
$\ln v$ is a rational function of $x$.

For the $UW$ term the corresponding determinant and substitution
are
\begin{align}
 \Delta_{UW}&=\frac{H^2z(H-x)[4x+z(H-x)]}{4s^2},\\
 z&=\frac{x(v-1)^2}{(H-x)v},\qquad
 \frac{\partial z}{\partial v}=\frac{x(v^2-1)}{(H-x)v^2},\\
 J_{UW}&=\frac{2sv\ln v}{Hx(v^2-1)},\qquad
 J_{UW}\dd z=\frac{2s\ln v}{H(H-x)v}\dd v.
\end{align}
Its transformed domain is $1<v<s-x$. Exchanging the integration
order gives
\begin{equation}
 1<v<s,\qquad0<x<s-v.
 \label{eq:comp-real-UW-rational}
\end{equation}
Again the inner integration is rational. The $VW$ sector follows
by exchanging $x$ and $y$ before making the same transformation.

Finally the logarithm $L_W$ contains
$\sqrt{(s+1-z)^2-4s}$. Use the positive solution
\begin{equation}
 s+1-z=u+\frac{s}{u},\qquad\sqrt{s}\leq u\leq s.
\end{equation}
The original Dalitz bounds at fixed $z$ become
\begin{equation}
 s-u\leq x\leq s-\frac{s}{u},\qquad
 L_W|\dd z|=\frac{s}{Hu}\ln\frac{u^2}{s}\dd u.
 \label{eq:comp-real-W-rational}
\end{equation}
Thus every term has a rational inner integral with algebraic
boundaries, followed by one logarithmic integration. The square
roots and their signs have been fixed by positive physical energies
throughout.

These separated logarithmic terms must not be assigned unrelated
endpoint cutoffs. For example, a common center-of-mass soft boundary
$\omega_i>\varepsilon$ means $x,y>\delta$, where
$\delta=2\sqrt{s}\,\varepsilon$. In
\eqref{eq:comp-real-UV-rational} it becomes
\begin{equation}
 \delta<x<\frac{H^2v(H-\delta)}{H^2v-\delta(v-1)^2}.
\end{equation}
In \eqref{eq:comp-real-UW-rational} the same boundary is
$x>\delta$ together with
$H-x+x(v-1)^2/[(H-x)v]>\delta$.
The finite terms generated by these transformed boundaries belong
to the cross section. Replacing them by arbitrary small limits
would change the finite answer. Apparent poles at $z=0$ introduced
by the partial fractions must likewise cancel after the rational
and logarithmic terms are combined. The transformations must therefore be combined with the local
Gram subtraction and the common regulator matching developed below.

\subsection{Removal of artificial Gram singularities}
The original angular average is regular when $z=0$ at fixed
$0<x<H$. Individual logarithmic coefficients contain $z^{-2}$
because of the partial-fraction reduction. To remove these
artificial singularities locally, write $a=H-x$ and define
\begin{align}
 T_{UV}(z)&=\frac{s}{xa}-\frac{s(6x+H^2)}{6x^2a^2}\,z,\\
 T_{UW}(z)&=\frac{s}{xH}-\frac{sa}{6Hx^2}\,z,\\
 T_{VW}(z)&=\frac{s}{aH}-\frac{s(x+6)}{6Ha^2}\,z.
\end{align}
These are the constant and linear Taylor terms of the corresponding
$J$ functions, including $y=a+z$. They follow from
\begin{equation}
 \frac{\operatorname{arcosh}(1+q)}{\sqrt{(1+q)^2-1}}
       =1-\frac q3+\mathcal O(q^2).
\end{equation}
Replace each $R_{ij}J_{ij}$ by
$R_{ij}(J_{ij}-T_{ij})$, and add the subtracted pieces to
\begin{equation}
 R_0^{\mathrm{reg}}=R_0+
             R_{UV}T_{UV}+R_{UW}T_{UW}+R_{VW}T_{VW}.
\end{equation}
This is an exact algebraic identity. Each difference $J-T$ is
$\mathcal O(z^2)$. On substituting the explicit coefficients,
the reduced rational denominator of $R_0^{\mathrm{reg}}$ has no
factor $z$. Thus all these contributions can be integrated across
$z=0$ without assigning separate unphysical endpoint cutoffs.
This operation does not remove either physical soft endpoint.

\subsection{Analytic regulation of the soft endpoints}
A Bose-symmetric analytic regulator makes the remaining endpoint
operations independent of frame-dependent energy boundaries.
Multiply the real phase-space integrand by $(xy)^\eta$, initially
with positive real $\eta$ large enough for convergence, and then
continue meromorphically to $\eta=0$. Since $x,y>0$ in the physical
region, the real logarithm defines this factor without a branch
ambiguity. This is an intermediate integration device. The virtual
amplitude and its infrared matching still use the photon mass
$\lambda$.

Define the dimensionless regulated real coefficient by
\begin{equation}
 \mathcal C_{\mathrm R}(\eta)=\frac1{8sH}
 \int_0^H\dd x\int_{H-x}^{s(H-x)/(s-x)}\dd y\,
     (xy)^\eta\langle\mathcal R\rangle
 =\frac{A(s)}{\eta}+F(s)+\mathcal O(\eta).
 \label{eq:comp-real-analytic-regulator}
\end{equation}
The regulator is applied to the complete trace, including the
identical-photon factor already present in its normalization.
At one soft endpoint $x=2\sqrt{s}\,\omega$ while $y\to H$.
Consequently $(xy)^\eta\to(2\sqrt{s}H\omega)^\eta$.
The other endpoint gives the same contribution. Each carries
one half of the usual eikonal result because of $1/2!$, so together
\begin{equation}
 A(s)=\frac1s\int_{-1}^1\dd c\,
             \mathcal B(s,c)[b_e(s,c)-1].
 \label{eq:comp-real-pole-residue}
\end{equation}
The functions $\mathcal B$ and $b_e$ are the Born square and the
relative-velocity factor defined earlier.

This relation also fixes the finite matching constant, rather than
only the pole. Put $C=2\sqrt{s}H$. Integration of the leading soft
term below a center-of-mass energy $\varepsilon$ gives
\begin{equation}
 A(s)C^\eta\int_0^\varepsilon\omega^{\eta-1}\dd\omega
 =\frac{A(s)}\eta+A(s)\ln(C\varepsilon)+\mathcal O(\eta).
\end{equation}
Thus the real coefficient above that energy boundary is
$F(s)-A(s)\ln(C\varepsilon)+o(1)$.
The photon-mass soft term derived previously is
\begin{equation}
 \mathcal C_{\mathrm{soft}}(\varepsilon,\lambda)
 =A(s)\ln\frac{2\varepsilon}{\lambda}
                  -\frac1s\int_{-1}^1\mathcal B K\,\dd c.
\end{equation}
Adding it to the virtual contribution
$\mathcal C_{\mathrm{virt}}=(4s)^{-1}\int\operatorname{Re}\mathcal T\,\dd c$
therefore gives
\begin{equation}
 \mathcal C_{\mathrm{total}}
 =\mathcal C_{\mathrm{virt}}+F(s)
   +A(s)\ln\frac1{\sqrt{s}H\lambda}
   -\frac1s\int_{-1}^1\mathcal B K\,\dd c.
 \label{eq:comp-real-regulator-matching}
\end{equation}
The coefficient of $\ln\lambda$ cancels the virtual coefficient
$A(s)\ln\lambda$. Both the analytic pole and the temporary energy
boundary have disappeared. Equation~\eqref{eq:comp-real-regulator-matching}
fixes how the finite part $F(s)$ must be combined with the earlier
photon-mass calculation. The following analytic primitives and endpoint formulas evaluate
this finite part without phase-space quadrature.

\subsection{A complete logarithmic-sector primitive}
The part proportional to $L_W$ can be integrated further without
any phase-space quadrature. Its rational coefficient has the simple
partial-fraction form
\begin{equation}
 R_W=C_0(z)+C_1(z)\left(\frac1{x-z}-\frac1{x-H}\right),
\end{equation}
where, with $D_z=(H-z)^2-4z$,
\begin{align}
 C_0(z)&=-\frac8{HD_z}
       [H^2z-4H^2-2Hz^2+12Hz+z^3-8z^2+16z],\\
 C_1(z)&=\frac4{HD_z}\bigl[
       -H^3z-4H^3+3H^2z^2+8H^2z-8H^2\\
 &\hspace{25mm}{}-3Hz^3+2Hz^2+32Hz+z^4-6z^3+32z\bigr].
\end{align}
At fixed $z$, the exact limits are
$x_\pm=(H+z\pm\sqrt{D_z})/2$. Their substitution into the
rational primitive gives
\begin{equation}
 \int_{x_-}^{x_+}R_W\dd x
 =C_0(z)\sqrt{D_z}+2C_1(z)
       \ln\frac{H-z+\sqrt{D_z}}{H-z-\sqrt{D_z}}.
\end{equation}
Use $z=s+1-u-s/u$ as above. Define two rational functions
\begin{equation}
 A(u)=\frac{s}{Hu}C_0(s+1-u-s/u)(u-s/u),\qquad
 B(u)=\frac{2s}{Hu}C_1(s+1-u-s/u).
\end{equation}
The remaining integrand of this sector is exactly
\begin{equation}
 A(u)\ln\frac{u^2}{s}
 +B(u)\ln\frac{u^2}{s}\ln\frac{u(u-1)}{s-u},\qquad
 \sqrt{s}<u<s.
 \label{eq:comp-real-W-last-integrand}
\end{equation}
Only simple real logarithms occur. The real branches are fixed by
$u>1$ and $s-u>0$.

For a compact analytic primitive put $r=\sqrt{s}>1$. Denote by
$G_r(a_1,\ldots,a_n;u)$ the generalized polylogarithm based at $r$,
with derivative
\begin{equation}
 \frac{\dd}{\dd u}G_r(a_1,\ldots,a_n;u)
 =\frac{G_r(a_2,\ldots,a_n;u)}{u-a_1},\qquad G_r(;u)=1.
\end{equation}
The integration path is the real interval from $r$ to $u$. At a
basepoint letter $r$ use the continuous or combined tangential
limit. All other letters below are $0,1,-r,s$, which lie outside
the open interval. Thus the functions that occur are real there.
In particular, suppressing $r,u$ in the notation,
\begin{equation}
 \ln(u^2/s)=2G(0),\qquad
 \ln[u(u-1)/(s-u)]=G(0)+G(1)-G(s).
\end{equation}
Shuffle multiplication therefore reduces
\eqref{eq:comp-real-W-last-integrand} to a finite weight-two
integrand. Its primitive has weight at most three and is
\begin{align}
 \mathcal F_W={}&-\frac{16(r+1)^2}{r-1}
       [G(-r,0)+G(-r,1)-G(-r,s)]\\
 &+\frac{16(r-1)^2}{r+1}
       [G(r,0)+G(r,1)-G(r,s)]\\
 &-\frac{32r^2(2r^2+1)}{(r^2-1)^2}
       [2G(0,0,0)+G(0,0,1)+G(0,1,0)-G(0,0,s)-G(0,s,0)]\\
 &+A_{00}G(0,0)+A_{01}[G(0,1)-G(0,s)]
        +A_{10}G(1,0)+A_{s0}G(s,0)\\
 &+A_0G(0)+A_1[G(1)-G(s)]+A_{\mathrm c}.
 \label{eq:comp-real-W-GPL-primitive}
\end{align}
The rational coefficients below are obtained by partial fractions
and integration by parts. They contain no fitted constants.
Differentiating this finite expression yields
\eqref{eq:comp-real-W-last-integrand}.
The upper soft endpoint must still be combined with the other
sectors in the common analytic regulator. This primitive alone
is not the finite inclusive real correction.
% BEGIN REAL W GPL COEFFICIENTS
\begin{equation}
A_{00}=\frac{N_{00}}{u^{2} \left(r - 1\right)^{2} \left(r + 1\right)^{2} \left(r - u\right) \left(r + u\right)}.
\end{equation}
\begin{equation}
\begin{aligned}
N_{00}={}& - 16 r^{8} + 64 u^{4} - 192 r^{2} u^{3} - 112 r^{4} u^{4} - 48 r^{6} u^{2} - 16 r^{2} u^{6}\\
& - 8 r^{6} u^{4} + 8 r^{8} u^{2} + 32 r^{6} u + 32 r^{8} u + 32 r^{2} u^{5} + 32 r^{4} u^{5}\\
& + 56 r^{4} u^{2} + 64 r^{4} u^{3} + 72 r^{2} u^{4}.
\end{aligned}
\end{equation}
\begin{equation}
A_{01}=\frac{N_{01}}{u^{2} \left(r - 1\right)^{2} \left(r + 1\right)^{2} \left(r - u\right) \left(r + u\right)}.
\end{equation}
\begin{equation}
\begin{aligned}
N_{01}={}& - 8 r^{8} + 32 u^{4} - 96 r^{2} u^{3} - 88 r^{4} u^{4} - 8 r^{2} u^{6} + 8 r^{6} u^{2} + 16 r^{6} u\\
& + 16 r^{8} u + 16 r^{2} u^{5} + 16 r^{4} u^{5} + 32 r^{4} u^{3} + 64 r^{2} u^{4}.
\end{aligned}
\end{equation}
\begin{equation}
A_{10}=\frac{N_{10}}{u^{2} \left(r - 1\right)^{2} \left(r + 1\right)^{2} \left(r - u\right) \left(r + u\right)}.
\end{equation}
\begin{equation}
\begin{aligned}
N_{10}={}& - 8 r^{8} - 96 r^{2} u^{3} - 56 r^{4} u^{4} - 24 r^{6} u^{2} - 8 r^{2} u^{6} - 8 r^{8} u^{2}\\
& + 8 r^{4} u^{2} + 8 r^{6} u^{4} + 16 r^{6} u + 16 r^{8} u + 16 r^{2} u^{5} + 16 r^{4} u^{5}\\
& + 32 r^{2} u^{2} + 32 r^{4} u^{3} + 56 r^{2} u^{4}.
\end{aligned}
\end{equation}
\begin{equation}
A_{s0}=\frac{N_{s0}}{u^{2} \left(r - 1\right)^{2} \left(r + 1\right)^{2} \left(r - u\right) \left(r + u\right)}.
\end{equation}
\begin{equation}
\begin{aligned}
N_{s0}={}& - 32 u^{4} + 8 r^{8} - 56 r^{4} u^{2} - 32 r^{4} u^{3} - 16 r^{6} u - 16 r^{8} u - 16 r^{2} u^{5}\\
& - 16 r^{4} u^{5} - 8 r^{2} u^{4} - 8 r^{8} u^{2} + 8 r^{2} u^{6} + 8 r^{6} u^{4} + 24 r^{4} u^{4}\\
& + 56 r^{6} u^{2} + 96 r^{2} u^{3}.
\end{aligned}
\end{equation}
\begin{equation}
A_{0}=\frac{12 r^{2} \left(r^{4} u + 2 r^{4} + 2 r^{2} u^{2} + 6 r^{2} + 6 u^{2} - u\right)}{u \left(r - 1\right)^{2} \left(r + 1\right)^{2}}.
\end{equation}
\begin{equation}
A_{1}=\frac{N_{1}}{u^{2} \left(r - 1\right)^{2} \left(r + 1\right)^{2}}.
\end{equation}
\begin{equation}
\begin{aligned}
N_{1}={}& - 4 r^{6} - 32 r^{4} u^{2} - 12 r^{2} u^{2} - 12 r^{6} u^{2} - 4 r^{2} u^{4} + 16 r^{4} u + 16 r^{6} u\\
& + 16 r^{2} u^{3} + 16 r^{4} u^{3}.
\end{aligned}
\end{equation}
\begin{equation}
A_{\mathrm c}=\frac{4 r^{2} \left(r - u\right) \left(r + u\right) \left(7 r^{2} + 17\right)}{u \left(r - 1\right)^{2} \left(r + 1\right)^{2}}.
\end{equation}
% END REAL W GPL COEFFICIENTS

\subsection{The sector proportional to the energy logarithm}
The single denominators $U,V$ give another sector that can be
integrated completely in the interior. Substituting
$L_U=s\ln s/(xH)$ and $L_V=s\ln s/(yH)$ first combines their
coefficients into a rational function. At fixed $z$ its poles in
$x$ are only $0,z,H,H+z$, with orders at most three. The exact
limits $x_\pm$ therefore produce rational functions and the two
logarithms
\begin{equation}
 \ln\frac{x_+}{x_-}
   =\ln\frac{s(u-1)}{u(s-u)},\qquad
 \ln\frac{x_+-z}{x_--z}
   =\ln\frac{u(u-1)}{s-u}.
\end{equation}
After $z=s+1-u-s/u$, no square root remains. The final primitive
can be written with the same basepoint $r=\sqrt{s}$ as above.
Define the conjugate letters
\begin{equation}
 a_\pm=1\pm\ii\sqrt H,\qquad H=s-1>0.
\end{equation}
They are the roots of $u^2-2u+s$. Both lie off the real path,
and complex conjugation exchanges their GPLs. Thus their sum
below is manifestly real, with no choice based on numerical values.
The primitive of this entire logarithmic sector is $\ln s$ times
\begin{align}
 \mathcal F_{\ln s}={}&\frac{96s^2}{H^2}G(0,0)
 +\frac{32s(H^2-4H-6)}{H^2}[G(1,0)+G(s,0)]\\
 &+\frac{32s(H^2-H-3)}{H^2}\biggl[
 G(0,1)-G(0,s)\\
 &\hspace{18mm}{}-\sum_{a=a_\pm}
       \{G(a,0)+G(a,1)-G(a,s)\}\biggr]\\
 &+P_0G(0)+P_1[G(1)-G(s)]+P_{\mathrm c}.
 \label{eq:comp-real-log-s-primitive}
\end{align}
There are fourteen GPL words, all of weight at most two.
The rational coefficients are displayed below. Exact symbolic
differentiation of every GPL coefficient recovers the rational
and logarithmic integrand after the $x$ integration. This is an
independent check of the primitive, rather than a fit to a total
cross section. As in the previous sector, its soft endpoint is to
be combined in the common regulator before the finite part is taken.
% BEGIN REAL LOGS GPL COEFFICIENTS
\begin{equation}
P_{0}=\frac{Q_{0}}{H^{2} u^{2} \left(H + 1\right)}.
\end{equation}
\begin{equation}
\begin{aligned}
Q_{0}={}&12 - 48 u - 48 u^{3} - 32 u^{2} + 12 H^{4} + 12 u^{4} + 48 H + 48 H^{3} + 72 H^{2}\\
& - 136 H u - 120 H^{2} u - 88 H u^{3} - 56 H u^{2} - 32 H^{2} u^{3} - 24 H^{3} u\\
& + 8 H^{4} u + 8 H^{3} u^{3} + 12 H^{2} u^{4} + 24 H u^{4} + 28 H^{2} u^{2}\\
& + 36 H^{4} u^{2} + 80 H^{3} u^{2}.
\end{aligned}
\end{equation}
\begin{equation}
P_{1}=\frac{Q_{1}}{H^{4} u^{2} \left(H + 1\right)}.
\end{equation}
\begin{equation}
\begin{aligned}
Q_{1}={}& -112 - 1696 u^{2} - 796 H^{2} - 624 H^{3} - 480 H - 216 H^{4} - 112 u^{4} - 16 H^{5}\\
& + 4 H^{6} + 960 u + 960 u^{3} - 7448 H^{2} u^{2} - 5856 H u^{2} - 4152 H^{3} u^{2}\\
& - 932 H^{4} u^{2} - 256 H u^{4} - 172 H^{2} u^{4} - 112 H^{5} u^{2} - 36 H^{6} u^{2}\\
& - 24 H^{3} u^{4} + 4 H^{4} u^{4} + 24 H^{6} u + 24 H^{5} u^{3} + 224 H^{4} u^{3}\\
& + 248 H^{5} u + 1272 H^{3} u^{3} + 1496 H^{4} u + 2816 H u^{3} + 2928 H^{2} u^{3}\\
& + 3776 H u + 4200 H^{3} u + 5744 H^{2} u.
\end{aligned}
\end{equation}
\begin{equation}
P_{\mathrm c}=\frac{Q_{\mathrm c}}{H^{3} u^{2} \left(H + 1\right)}.
\end{equation}
\begin{equation}
\begin{aligned}
Q_{\mathrm c}={}& -56 - 368 u^{3} - 272 H^{2} - 206 H - 144 H^{3} - 16 H^{4} + 6 H^{5} + 56 u^{4} + 368 u\\
& - 880 H u^{3} - 660 H^{2} u^{3} - 136 H^{3} u^{3} - 6 H^{3} u^{4} - 4 H^{5} u\\
& + 4 H^{4} u^{3} + 28 H^{2} u^{4} + 94 H u^{4} + 132 H^{4} u + 796 H^{3} u + 1248 H u\\
& + 1540 H^{2} u.
\end{aligned}
\end{equation}
% END REAL LOGS GPL COEFFICIENTS



\subsection{Factorization of both soft endpoints in the pair sector}
The $UW$ sector admits a rectangular domain in which the analytic
regulator factorizes at both soft endpoints. Starting from
\eqref{eq:comp-real-UW-rational}, put
\begin{align}
 a&=H-x,\qquad v=1+at,\qquad0<a<H,\quad0<t<1,\\
 z&=\frac{a(H-a)t^2}{1+at},\qquad
 y=\frac{a[1+Ht^2+at(1-t)]}{1+at},\qquad
 \left|\frac{\partial(x,v)}{\partial(a,t)}\right|=a.
\end{align}
The product that regulates the two soft endpoints is therefore
\begin{equation}
 xy=a(H-a)f(a,t),\qquad
 f(a,t)=\frac{1+Ht^2+at(1-t)}{1+at}>0.
 \label{eq:comp-real-UW-reg-factor}
\end{equation}
Both endpoints are now ordinary powers of a single variable.
The regulator introduces no new physical branch.

Define two rational functions by the previously specified
coefficient $R_{UW}$ and Taylor subtraction $T_{UW}$,
\begin{align}
 L(a,t)&=a\,R_{UW}(x,z)\frac{2s}{H(H-x)v}
                  \bigg|_{x=H-a,\,v=1+at},\\
 T(a,t)&=a\,R_{UW}(x,z)T_{UW}(z)
                   \frac{\partial z}{\partial v}
                  \bigg|_{x=H-a,\,v=1+at},\\
 F(a,t)&=L(a,t)\ln(1+at)-T(a,t).
\end{align}
The subtraction sign in $F$ is fixed by the local removal of the
artificial Gram singularity. The complementary rational term is
already included in $R_0^{\mathrm{reg}}$.
The only poles of $L,T$ in the complex $a$ plane are
\begin{equation}
 a=0,\quad H,\quad-1/t,\quad-2/t,\quad-1/[t(1-t)].
\end{equation}
For $0<t<1$, the last three are negative and lie outside the
integration interval. Their orders are finite, at most five.
There is no polynomial denominator of degree higher than one
after partial fractions in $a$.

For clarity, the endpoint conversion can be written without a
hypergeometric expansion. Suppose
\begin{equation}
 F(a,t)=\frac{f_1(t)}a+\mathcal O(1),\qquad
 F(H-b,t)=\frac{g_2(t)}{b^2}+\frac{g_1(t)}b+\mathcal O(1).
\end{equation}
For the explicit coefficients above, $f_1(t)=0$ identically.
The Taylor subtraction has therefore also made the $a=0$ soft
corner finite in this particular sector. The other endpoint has
the stated maximum pole orders.
Let $\operatorname{FP}_a\int_0^H F\dd a$ denote the constant term
of the integral from $\delta$ to $H-\delta$ after all powers and
logarithms of $\delta$ are subtracted. Expansion of
\eqref{eq:comp-real-UW-reg-factor} at each endpoint gives
\begin{align}
 \int_0^H[xy]^\eta F(a,t)\dd a
 ={}&\frac{f_1(t)+g_1(t)}\eta
       +\operatorname{FP}_a\int_0^H F(a,t)\dd a\\
 &+f_1(t)\ln[H(1+Ht^2)]+g_1(t)\ln H\\
 &+g_2(t)\left[-\frac1H+\frac{t^2}{1+Ht}\right]
       +\mathcal O(\eta).
 \label{eq:comp-real-UW-eta-finite}
\end{align}
For example, the last term follows because the term linear in
$b$ in the logarithm of the regulator's smooth factor is
$-1/H+t^2/(1+Ht)$. Its multiplication by $\eta$ and a double
endpoint pole produces a finite term. Dropping that term would
change the answer even though the infrared pole remained correct.

Every primitive required on the right-hand side is elementary or
a real dilogarithm. To see this explicitly, define
\begin{equation}
 E_n(a;c)=\begin{cases}
 \ln|a-c|,&n=1,\\
 -(a-c)^{1-n}/(n-1),&n\geq2.
 \end{cases}
\end{equation}
For a simple logarithmic pole, let $w=1+at$ and $d=1+ct$.
When $c=-2/t$ or $c=-1/[t(1-t)]$, one has $d<0$, and a primitive is
\begin{equation}
 I_1(a;c)=\ln w\ln(1-w/d)+\operatorname{Li}_2(w/d).
\end{equation}
Both logarithms and the dilogarithm are real. For $c=H$ the same
formula applies with $0<w/d<1$ in the interior. The exceptional
values are
\begin{equation}
 I_1(a;0)=-\operatorname{Li}_2(-at),\qquad
 I_1(a;-1/t)=\tfrac12\ln^2(1+at).
\end{equation}
In each case $\partial_a I_1=\ln(1+at)/(a-c)$.
Higher poles follow by one integration by parts. With $b=-1/t$
and $n\geq2$,
\begin{align}
 I_n(a;c)&=-\frac{\ln(1+at)}{(n-1)(a-c)^{n-1}}
                  +\frac{Q_{n-1}(a;c,b)}{n-1},\\
 Q_m(a;c,b)&=\frac{E_1(a;b)}{(b-c)^m}
            -\sum_{j=1}^m\frac{E_j(a;c)}{(b-c)^{m-j+1}}
                         \quad(c\ne b),\\
 Q_m(a;b,b)&=E_{m+1}(a;b).
\end{align}
Thus partial fractions of the explicitly defined rational functions
$L,T$ give the finite part in
\eqref{eq:comp-real-UW-eta-finite} as a finite sum of boundary values
of $I_n$ and $E_n$. No numerical integration enters this inner
energy integration.

The nontrivial negative dilogarithm argument at $a=H$ is
$-(1-t)(1+Ht)/t$. Its complementary argument contains
\begin{equation}
 1+Ht-Ht^2,
 \qquad t_\pm=\frac12\left(1\pm\sqrt{1+4/H}\right).
\end{equation}
Here $t_-<0$ and $t_+>1$. The square root $\sqrt{H(H+4)}$
therefore follows directly from the remaining integration alphabet.
It has not been inferred from a published expression.
The inner integration leaves one angular variable $t$. We first
write its finite integrand explicitly, including its regulated
endpoint terms, so that the remaining primitive can be constructed
from its actual poles and logarithms.

\paragraph{Closed inner finite part.} The following boundary formulas
make the inner result a finite expression of logarithms and
dilogarithms, with no integration operation left. Put
$\ell=\ln(1+Ht)$ and write the partial fractions
\begin{equation}
 L=l_0+l_1a+\sum_{c,n}\frac{L_{cn}}{(a-c)^n},\qquad
 T=\sum_{j=0}^3t_j a^j+\sum_{c,n}\frac{T_{cn}}{(a-c)^n}.
\end{equation}
The pole set is the five-element set displayed above.
For a pole of order $m_c$, its coefficient is fixed algebraically,
\begin{equation}
 L_{cn}=\left.\frac1{(m_c-n)!}
  \left(\frac{\partial}{\partial a}\right)^{m_c-n}
       [(a-c)^{m_c}L(a,t)]\right|_{a=c},
\end{equation}
with the same formula for $T_{cn}$. The polynomial parts follow by
ordinary polynomial division. Define $\mathcal I_n(c)$ and
$\mathcal E_n(c)$ to be the finite boundary differences of the
explicit primitives $I_n,E_n$ above. At a negative pole they are
ordinary differences between $a=H$ and $a=0$. At the two physical
endpoints the required values are
\begin{align}
 \mathcal E_1(0)&=\ln H,\qquad
 \mathcal E_1(H)=-\ln H,\qquad
 \mathcal E_2(0)=\mathcal E_2(H)=-H^{-1},\\
 \mathcal I_1(0)&=-\operatorname{Li}_2(-Ht),\\
 \mathcal I_2(0)&=-\ell/H+t\ln H-t\ell+t,\\
 \mathcal I_1(H)&=-\ell\ln H-\operatorname{Li}_2(-Ht)-\ell^2/2,\\
 \mathcal I_2(H)&=-\frac{t}{1+Ht}(\ln H+\ell+1).
\end{align}
The two polynomial logarithmic moments are
\begin{align}
 P_0&=\frac{(1+Ht)\ell-Ht}{t},\\
 P_1&=\frac{[(1+Ht)^2/2-(1+Ht)]\ell
              -(1+Ht)^2/4+(1+Ht)-3/4}{t^2}.
\end{align}
The exact inner coefficient of $\eta^0$ is consequently
\begin{align}
 F_{UW}(t)={}&l_0P_0+l_1P_1
       -\sum_{j=0}^3\frac{t_jH^{j+1}}{j+1}
       +\sum_{c,n}[L_{cn}\mathcal I_n(c)-T_{cn}\mathcal E_n(c)]\\
 &+g_1(t)\ln H+g_2(t)\left[-H^{-1}+\frac{t^2}{1+Ht}\right].
 \label{eq:comp-real-UW-inner-closed}
\end{align}
This formula uses only finite algebraic differentiations and the
real functions already displayed. The corresponding $\eta^{-1}$
coefficient is $g_1(t)$.

For the last integration, every dilogarithm in this result has a
fixed-letter GPL representation. Put $b=-1/H$, $d=-2/H$, and
let $t_\pm$ be the two real roots above. With GPLs now based at zero,
\begin{align}
 \operatorname{Li}_2(-Ht)&=-G(0,b;t),\\
 \operatorname{Li}_2(-1-Ht)+\frac{\pi^2}{12}
       &=-\ln2\,G(b;t)-G(b,d;t).
\end{align}
\begin{equation}
\begin{aligned}
 &\operatorname{Li}_2\left(-\frac{(1-t)(1+Ht)}t\right)
 -\operatorname{Li}_2\left(-\frac{1-t}t\right)\\
 &\qquad=G(b,0;t)
          -\sum_{c=t_\pm}[G(1,c;t)+G(b,c;t)-G(0,c;t)].
\end{aligned}
\end{equation}
The last identity follows by differentiating both sides and using
their common zero limit as $t\to0^+$. It fixes its integration
constant as well as its branches. This gives a weight-two GPL
integrand with rational coefficients for the final energy integral.

\subsection{Closed construction for the regular rational sector}
The rational term $R_0^{\mathrm{reg}}$ also has an analytic primitive
after both Dalitz integrations. The compact construction below
specifies its coefficients by finite algebraic operations, without
introducing a new integral. At fixed $z$, its partial fractions are
\begin{equation}
 R_0^{\mathrm{reg}}=p_0(z)+
          \sum_{\rho\in\mathcal P}\sum_{n=1}^{m_\rho}
               \frac{c_{\rho n}(z)}{(x-\rho)^n},
\end{equation}
where there are sixteen nonzero pole terms and
\begin{equation}
 \mathcal P=\left\{0,z,H,H+z,
              \frac{Hz}{z-4},\ H+z-\frac{Hz}{z-4}\right\}.
\end{equation}
Every order $m_\rho$ is at most three. A uniform formula, including
vanishing coefficients when the actual order is smaller, is
\begin{align}
 c_{\rho n}(z)&=\left.\frac{1}{(3-n)!}
   \frac{\partial^{3-n}}{\partial x^{3-n}}
       [(x-\rho)^3R_0^{\mathrm{reg}}(x,z)]\right|_{x=\rho},
       \quad n=1,2,3,\\
 p_0(z)&=\lim_{x\to\infty}R_0^{\mathrm{reg}}(x,z).
\end{align}
The limit exists because its polynomial part is constant.
This supplies an explicit algebraic definition from the twenty-five
trace coefficients and the displayed Taylor subtractions.

Let $z=s+1-u-s/u$, $x_-=s-u$, $x_+=s-s/u$, and
$j=1-s/u^2$. Define four rational functions by
\begin{align}
 R(u)&=j\left[p_0(x_+-x_-)
  +\sum_{\rho}\sum_{n=2}^{m_\rho}
      \frac{c_{\rho n}}{1-n}
       \{(x_+-\rho)^{1-n}-(x_--\rho)^{1-n}\}\right],\\
 K_1(u)&=j[-c_{0,1}+c_{H+z,1}],\\
 K_2(u)&=j[-c_{H,1}+c_{0,1}+c_{z,1}-c_{H+z,1}],\\
 K_3(u)&=j[c_{Hz/(z-4),1}-c_{H+z-Hz/(z-4),1}].
\end{align}
All quantities on the right are evaluated at $z=s+1-u-s/u$.
The exact $x$ boundary differences then give the final integrand
\begin{equation}
 R(u)+K_1(u)L_1(u)+K_2(u)L_2(u)+K_3(u)L_3(u),
\end{equation}
where
\begin{align}
 L_1&=\ln(u^2/s),\qquad
 L_2=\ln[u(u-1)/(s-u)],\\
 L_3&=\ln\frac{(u-1)(u+s)^2}{u(u+1)^2(s-u)}.
\end{align}
Each logarithm vanishes at $u=\sqrt{s}$. In terms of GPLs based
there, the additional logarithm is
\begin{equation}
 L_3=G(1)+2G(-s)-G(0)-2G(-1)-G(s).
\end{equation}

For precision, define the finite algebraic operation
$\mathcal P[R;w]$ on a rational function $R$ and an empty or
one-letter word $w$. Decompose $R$ into polynomial and pole terms.
On a simple pole its value is
\begin{equation}
 \mathcal P[c/(u-a);w]=cG(a,w;u).
\end{equation}
For a pole of order $n\geq2$ and an empty word it is
$-c/[(n-1)(u-a)^{n-1}]$. For a one-letter word $(b)$ it is
\begin{equation}
 \mathcal P[c/(u-a)^n;(b)]
 =-\frac{cG(b;u)}{(n-1)(u-a)^{n-1}}
                  +\frac{cQ_{n-1}(u;a,b)}{n-1},
\end{equation}
where $Q_m$ is the finite elementary combination already given
above, with logarithms interpreted as the same basepoint GPLs.
Polynomial terms are fixed explicitly, for $k\geq0$, by
\begin{align}
 \mathcal P[u^k;()]&=\frac{u^{k+1}}{k+1},\\
 \mathcal P[u^k;(b)]&=\frac{u^{k+1}-b^{k+1}}{k+1}G(b;u)
  -\frac1{k+1}\sum_{j=0}^k
        \frac{b^j u^{k+1-j}}{k+1-j}.
\end{align}
Linearity completes the definition for any rational coefficient.
These operations terminate after one
step because $w$ has length at most one. Thus $\mathcal P$ denotes
a finite sum of specified GPLs and rational functions.

The complete interior primitive of the regular rational sector is
\begin{align}
 \mathcal F_{\mathrm{reg}}={}&\mathcal P[R;()]
 +\mathcal P[2K_1+K_2-K_3;(0)]
 +\mathcal P[K_2+K_3;(1)]\\
 &-\mathcal P[K_2+K_3;(s)]
 +2\mathcal P[K_3;(-s)]-2\mathcal P[K_3;(-1)].
 \label{eq:comp-real-rational-primitive}
\end{align}
Its expansion contains forty-eight GPL words of weight at most two.
This form keeps the relatively long rational coefficient table
specified by exact derivatives instead of reprinting it.

A possible internal letter in this representation is removable.
The quadratic $u^2-(s-3)u+s$ corresponds to $z=4$. The original
angular denominator $4x+z(H-x)$ is strictly positive in the physical
interior, so $z=4$ is not a singularity of the amplitude. The pole
$x=Hz/(z-4)$ merely moves through infinity at that point.
For $z<4$, the pole $Hz/(z-4)$ is negative and its partner
$H+z-Hz/(z-4)$ exceeds $H+z$. For $z>4$, the first exceeds
$H$ and the partner is smaller than $z$. Both therefore stay outside
$z<x_-\leq x\leq x_+<H$ on either side of $z=4$.
The original finite inner integral is analytic at this point.
At either root of the quadratic, $L_3=0$. Consequently the
combination $K_3L_3$ has a removable limit there, even when one root
lies inside the real $u$ interval. The GPL combination must be
formed before that limit is taken. Passing above or below the
individual artificial letter gives the same complete expression,
because its would-be discontinuity is proportional to this zero.
This argument fixes the continuation from the original real
integrand and does not add a physical threshold at $s=9$.

\paragraph{Exchange symmetry after the local subtraction.}
The full angular expression is invariant under the exchange of the
identical final photons. This exchange sends $x$ to $y=H-x+z$ and
leaves $z$ unchanged. The Taylor polynomials used to remove the Gram
zero were instead expanded at fixed $x$. These operations do not
commute beyond the order of the subtraction. In particular, direct
substitution gives
\begin{equation}
 \begin{split}
 T_{UW}(y,z)-T_{VW}(x,z)
 ={}&-\frac{s z^2}{6H(x-H)^2(x-H-z)^2}\\
 &\times[-H^2-6H+x^2-xz+6x-6z].
 \end{split}
\end{equation}
Using the coefficient $R_{VW}$ of $J_{VW}$, define the rational
exchange correction
\begin{equation}
 D(x,z)=R_{VW}(x,z)[T_{UW}(H-x+z,z)-T_{VW}(x,z)].
\end{equation}
The factor $z^2$ cancels the apparent Gram pole of $R_{VW}$.
The symmetric regulator $(xy)^\eta$, the Dalitz domain, and its
measure are all invariant under photon exchange. Changing variables
in the subtracted $VW$ integral therefore proves
\begin{equation}
 I_{VW}^{\mathrm{sub}}(\eta)
 =I_{UW}^{\mathrm{sub}}(\eta)+I_D(\eta).
\end{equation}
This identity holds first where the regulated integrals converge
and then for their meromorphic continuations. It applies separately
to the pole and finite Laurent coefficients. Consequently the real
integral is assembled from the regular rational term, the two
single-denominator logarithms, the $W$ logarithm, the $UV$ term,
twice the $UW$ term, and $D$. Omitting $D$ would change both the
infrared pole and the finite answer.

\paragraph{Endpoint values of the primitives.}
The preceding formulas specify finite combinations of GPLs rather
than unevaluated phase-space integrals. Their endpoint subtraction
can also be stated algebraically. Write $b=s-u$, let $f(u)$ denote
one complete sector integrand, and let $\mathcal F'(u)=f(u)$.
After combining all terms of that sector its expansion is
\begin{equation}
 f(s-b)=\frac{A}{b}
   +\sum_{n=0}^{\infty}(c_n+d_n\ln b)b^n.
 \label{eq:comp-real-endpoint-series}
\end{equation}
There are no higher negative powers or $\ln b/b$ terms in these
three sectors. This follows by expanding the displayed rational
coefficients and the logarithms before taking the endpoint. The
coefficient $A$ vanishes for the $W$ sector. All series coefficients
are obtained by polynomial division and the Taylor series of
$\ln(1-b/a)$.

Let $\mathcal F_0$ be the finite limit of the complete primitive at
$u=\sqrt{s}$. It includes constant contributions from products of a
vanishing GPL and a rational coefficient with a pole. They are
obtained by multiplying their local series before setting
$u=\sqrt{s}$.

The lower endpoint constants simplify considerably. Direct local
expansion of the finite GPL combinations gives
\begin{align}
 \mathcal F_{0,W}&=\mathcal F_{0,\ln s}=0,\\
 \mathcal F_{0,\mathrm{reg}}
 &=\frac{2s[-s^5-5s^4+18s^3+122s^2+151s-45]}
              {3(s^2+2s-3)}.
\end{align}
For the $W$ primitive, the coefficients that have a simple pole at
the basepoint multiply weight-two words. Those words vanish
quadratically, so they do not contribute a constant. For the
regular rational primitive the nonempty-word coefficients are
regular there. Its lower endpoint constant is therefore the
rational empty-word coefficient evaluated at $u=\sqrt{s}$.

Since
\begin{equation}
 z=\frac{H}{s}b+O(b^2),
\end{equation}
the finite part defined by a cut on $z$ is exactly
\begin{align}
 F_z={}&\mathcal F(s-\delta)-\mathcal F_0
       +A\ln\delta+A\ln(H/s)\\
 &+\sum_{n=0}^{\infty}\delta^{n+1}
  \left[\frac{c_n}{n+1}
       +d_n\left(\frac{\ln\delta}{n+1}
                         -\frac{1}{(n+1)^2}\right)\right].
 \label{eq:comp-real-gpl-endpoint-finite}
\end{align}
Here $\delta>0$ lies inside the convergence disk of the endpoint
series. The series is simply the analytic tail of the primitive.
It is not a numerical phase-space integration. Equivalently,
$F_z=\lim_{\epsilon\to0}[\int_{z>\epsilon}f\,du+A\ln\epsilon]$.
The subsequent conversion to $(xy)^\eta$ must still be included.

The artificial interior GPL letters can be bypassed during series
continuation without choosing a new physical branch. For the
regular rational sector the original finite $x$ integral is analytic
at $z=4$: its denominators have no zero there for interior $x$.
Differentiation of the specified primitive gives this same analytic
function, including at the removable limit. Its integral around a
small closed contour surrounding an artificial letter is therefore
zero. Continuations above and below it agree when the entire GPL
combination is formed. The starting value is fixed by
$\mathcal F_0$ on the real interval. This proves path independence
of the complete expression without dropping an imaginary part of
an individual GPL.

For $1<s<9$ the additional roots are complex and do not meet the
integration interval. For $s>9$ the only root that can meet it is
\begin{equation}
 q_+=\frac{s-3+\sqrt{(s-1)(s-9)}}{2}.
\end{equation}
The root at $s=9$ coincides with the lower endpoint. The original
$x$ integral remains regular there, and the coincident GPL words
are combined before their basepoint series is evaluated. Thus the
same analytic continuation covers both sides of $s=9$.

Exact symbolic differentiation verifies the displayed GPL
primitives for all three sectors.

\subsection{Analytic soft-endpoint terms in the inner energy integral}
The $UV$ logarithmic sector illustrates why the regulator must be
kept until both endpoints have been integrated. Put
\begin{equation}
 k=\frac{(v-1)^2}{Hv},\qquad u=\frac{x}{H},\qquad
 y=\frac{H(1-u)}{1-ku},\qquad
 (xy)^\eta=H^{2\eta}u^\eta(1-u)^\eta(1-ku)^{-\eta}.
\end{equation}
For $1<v<s=1+H$, $0<k<H/s<1$. Thus the third pole
$x=H/k$ is outside the entire inner interval $0<x<H$.
The rational coefficient of $\ln v$ has precisely the form
\begin{equation}
 R(x,v)=C_0+\frac{C_{01}}x+\frac{C_{02}}{x^2}
 +\frac{C_{H1}}{x-H}+\frac{C_{H2}}{(x-H)^2}
 +\frac{C_{a1}}{x-H/k}+\frac{C_{a2}}{(x-H/k)^2}.
\end{equation}
The coefficients are obtained by ordinary partial fractions of the
explicit $R_{UV}$ above after the change of variables.
The script \nolinkurl{compton_real_inner_analytic.py} performs this
algebra and verifies exact reconstruction of the original rational
function. No kinematic fitting is involved.

Each integral is an Euler beta function times a Gauss
hypergeometric function. For example, for sufficiently large
$\operatorname{Re}\eta$,
\begin{align}
 \int_0^H\frac{(xy)^\eta}{x^n}\dd x
 &=H^{1-n+2\eta}B(1+\eta-n,1+\eta)
 {}_2F_1(\eta,1+\eta-n;2+2\eta-n;k),\\
 \int_0^H\frac{(xy)^\eta}{(x-H)^n}\dd x
 &=(-1)^nH^{1-n+2\eta}B(1+\eta,1+\eta-n)
 {}_2F_1(\eta,1+\eta;2+2\eta-n;k),\\
 \int_0^H\frac{(xy)^\eta}{(x-H/k)^n}\dd x
 &=(-H/k)^{-n}H^{1+2\eta}B(1+\eta,1+\eta)
 {}_2F_1(n+\eta,1+\eta;2+2\eta;k).
\end{align}
Here $B(a,b)=\Gamma(a)\Gamma(b)/\Gamma(a+b)$. The Gauss function is
\begin{equation}
 {}_2F_1(a,b;c;k)=\sum_{n=0}^{\infty}
 \frac{(a)_n(b)_n}{(c)_n n!}k^n,\qquad |k|<1,
 \qquad (a)_n=\frac{\Gamma(a+n)}{\Gamma(a)},\quad (a)_0=1.
\end{equation}
Expanding $(1-ku)^{-a}$ and integrating each power against
$u^{b-1}(1-u)^{c-b-1}$ gives the displayed beta integrals in their
convergence domain. These expressions
then uniquely define the continuation to $\eta=0$. Their Laurent
expansion gives, at fixed $v$ in the interior,
\begin{align}
 \int_0^H (xy)^\eta R(x,v)\dd x
 &=\frac{C_{01}-C_{H1}}\eta+F_{UV}(v)+\Order(\eta),\\
 F_{UV}(v)&=H C_0+\frac{k-2}{H}C_{02}
 -\frac{2-k}{H(1-k)}C_{H2}
 +2(C_{01}-C_{H1})\ln H\\
 &\quad +(C_{H1}+C_{a1})\ln(1-k)
 +\frac{k^2}{H(1-k)}C_{a2}.
 \label{eq:comp-real-UV-inner-finite}
\end{align}
The double-pole terms have finite limits but no $1/\eta$ pole.
Their finite constants cannot be obtained by simply setting
$\eta=0$ in the integrand. For example,
\begin{equation}
 (1-u)^\eta(1-ku)^{-\eta}
 =1+\eta(k-1)u+\Order(u^2),\qquad
 \int_0^1u^{\eta-2}\dd u=\frac1{\eta-1}.
\end{equation}
The linear term contributes $\eta(k-1)/\eta=k-1$, which combines
with $-1$ from the first term to give $k-2$.
This explicitly accounts for the finite endpoint contribution.

Equation~\eqref{eq:comp-real-UV-inner-finite} still has to be
combined with the local $z=0$ subtraction and integrated over $v$.
It must not be integrated as an isolated term across a spurious
Gram pole before those cancellations have been made.

\section{Matching the real-energy endpoints}
\subsection{Soft endpoints of the real-energy integrals}
The remaining rational and single-logarithm sectors have elementary
soft residues. Put $H=s-1$, and use $e=x$ at the first soft endpoint
and $e=y$ at the second. Here $e$ is an energy invariant, not the
charge. With $z=er$, the exact phase-space boundary is
\begin{equation}
 0<e<H,\qquad 0<r<\rho(e)=\frac{H-e}{H+1-e},\qquad
 xy=e[H+e(r-1)],\qquad \dd x\,\dd z=e\,\dd e\,\dd r.
 \label{eq:comp-real-soft-domain}
\end{equation}
To obtain this domain from $0<z<x(H-x)/(s-x)$, set
$x=e$ and divide the inequality by $e>0$. At the other soft
endpoint set $y=e$ and use $x=H-e+z$ in the same inequality.
Both give $r<\rho(e)$. The two changes of variables have Jacobian
magnitude $e$. A disjoint partition near the two soft corners
allows their residues to be added without double counting. Away
from both corners the regulators converge to one on compact
subsets, so their difference has zero limit there.

Let $q=\rho(0)=H/(H+1)$. For a sector with a simple soft
singularity, define its two residues by
\begin{equation}
 f_X(r)=\lim_{e\to0}e^2R(H,e,er),\qquad
 f_Y(r)=\lim_{e\to0}e^2R(H,H+e(r-1),er).
\end{equation}
The rational sector and the two terms multiplying $\ln s$ have
this behavior. The coefficient of $L_W$ is only of order $e^{-1}$,
and $L_W(0)$ is finite, so this sector needs no soft matching term.
All coefficients in this section precede the common normalization
$1/(8sH)$ of the real cross section.

Write $f=f_X+f_Y$. The regulator $(xy)^\eta$ gives a pole
$A/\eta$ and finite part $F_\eta$. If the energy primitives are
instead evaluated with $z>\zeta$, define their finite part by
subtracting $-A\ln\zeta$. The symbol $\operatorname{FP}$ denotes
the constant term after that subtraction. Then
\begin{align}
 A&=\int_0^q f(r)\,\dd r,\\
 F_\eta&=\operatorname{FP}_{\zeta\to0}
                \int_{z>\zeta}R\,\dd x\,\dd z
       +A\ln H-\int_0^q f(r)\ln r\,\dd r.
 \label{eq:comp-real-simple-matching}
\end{align}
Indeed, the soft measure is $f(r)\dd e/e$, while $z>\zeta$
sets its lower endpoint to $e=\zeta/r$. The analytic regulator
instead multiplies this measure by $(He)^\eta$ at leading order.
The difference produces the two logarithmic terms above. The
$O(e)$ change in $\rho(e)$ is integrable for a simple soft pole.

These residue integrals do not require numerical quadrature.
Partial fractions give a polynomial and poles at $r=1$ and
$r=-4/H$, both outside $0\leq r\leq q$. The necessary moments are
\begin{align}
 \int_0^q r^n\,\dd r&=\frac{q^{n+1}}{n+1},\\
 \int_0^q r^n\ln r\,\dd r
 &=q^{n+1}\left[\frac{\ln q}{n+1}-\frac1{(n+1)^2}\right],\\
 \int_0^q\frac{\dd r}{r-a}&=\ln(1-q/a),\\
 \int_0^q\frac{\ln r}{r-a}\,\dd r
 &=\ln q\ln(1-q/a)+\operatorname{Li}_2(q/a).
\end{align}
Here $n\geq0$ and $a$ is a pole outside the interval.
Repeated poles follow by applying
$\partial_a^{j-1}/(j-1)!$. Every logarithm above has a positive
argument, and $q/a<1$. These formulas therefore specify real
branches directly from the physical integration interval.
The independently generated residue coefficients and their exact
moments are recorded in
\nolinkurl{compton_real_soft_matching_integrated.json}.


\paragraph{Bose symmetry of the subtractions.} The two mixed angular
kernels $J_{UW}$ and $J_{VW}$ are exchanged by $x\leftrightarrow y$.
Their local Taylor subtractions, however, were expanded at fixed
$x$. Thus $T_{VW}(H,x,z)$ is not exactly
$T_{UW}(H,H-x+z,z)$. To replace the subtracted $VW$ sector by a
second copy of $UW$, the rational sector must receive
\begin{equation}
 D(H,x,z)=R_{VW}(H,x,z)
 \big[T_{UW}(H,H-x+z,z)-T_{VW}(H,x,z)\big].
 \label{eq:comp-real-bose-restoration}
\end{equation}
The bracket starts at $z^2$, so it introduces no artificial Gram
pole. It can nevertheless change the soft residue because its
expansion at fixed $x$ is not uniform as $y\to0$.

At this endpoint write
$D(H,H+e(r-1),er)=a(r)e^{-3}+b(r)e^{-2}+O(e^{-1})$.
The coefficients needed in this endpoint expansion are explicit:
\begin{align}
 a(r)&=\frac{16(H+1)(Hr+2)(Hr-2H-6)}{3H(r-1)^3},\\
 b(r)&=\frac{2(H+1)}{3H^2(r-1)^3(Hr+4)^2}
                    \sum_{j=0}^5 B_j(H)r^j.
\end{align}
The six polynomials are
\begin{align}
 B_0(H)&=256 H,\\
 B_1(H)&=64 \left(H^{3} + 6 H^{2} - 11 H - 12\right),\\
 B_2(H)&=16 \left(H^{4} + 2 H^{3} + 2 H^{2} + 128 H + 108\right),\\
 B_3(H)&=- 8 \left(H^{5} + 8 H^{4} - 19 H^{3} - 160 H^{2} - 192 H - 72\right),\\
 B_4(H)&=- 2 H \left(H^{5} + 3 H^{4} - 36 H^{3} - 108 H^{2} - 152 H - 96\right),\\
 B_5(H)&=H^{2} \left(H + 6\right) \left(H^{3} - 4 H^{2} - 12 H - 8\right).
\end{align}
These follow by substituting the exact soft coordinates into
\eqref{eq:comp-real-bose-restoration} and expanding the resulting
rational numerator and denominator through the first two powers.

The moving angular boundary now contributes to the logarithmic
soft coefficient. Since $\rho'(0)=-1/(H+1)^2$, it is
\begin{equation}
 A_D=\int_0^q b(r)\,\dd r-\frac{a(q)}{(H+1)^2}.
 \label{eq:comp-real-bose-pole}
\end{equation}
More explicitly, the radial density after the angular integration is
\begin{align}
 \int_0^{\rho(e)} eD\,\dd r
 &=\frac{1}{e^2}\int_0^q a(r)\,\dd r
   +\frac1e\left[\int_0^q b(r)\,\dd r+\rho'(0)a(q)\right]
   +O(1).
\end{align}
The first term is a power singularity, whose analytic continuation
has a pole at $\eta=1$. The coefficient of $1/e$ gives the pole
at $\eta=0$. This explains why expanding at a fixed angular upper
limit would miss the second term in
\eqref{eq:comp-real-bose-pole}. The terms generated by the
$O(\eta e)$ part of $(xy)^\eta$ affect its finite part and must
not be discarded.

Its finite contribution must also be included when combining the
real sectors. Omitting $D$ would violate the integrated soft
coefficient even though the original amplitude is Bose symmetric.

\paragraph{Integrated infrared check.} Denote the poles of the
rational, $\ln s$, $UV$, and $UW$ sectors by
$A_{\mathrm{rat}},A_{\log},A_{UV},A_{UW}$, respectively.
The required sum, before the common factor $1/(8sH)$, is
\begin{equation}
 A_{\mathrm{rat}}+A_{\log}+A_{UV}+2A_{UW}+A_D
 =8H\int_{-1}^1\mathcal B(c)\,[b_e(c)-1]\,\dd c.
 \label{eq:comp-real-integrated-pole-check}
\end{equation}
The right-hand side follows independently from the soft current
and the Born trace. It equals the coefficient required to cancel
the photon-mass logarithm of the virtual amplitude after the
regulator matching.

\subsection{Matching a moving soft boundary}
\label{sec:compton-soft-matching-check}

The separate rational subtraction terms can have stronger soft poles than
 their sum.  Their finite parts therefore require the exact phase-space
boundary and the first subleading term of the regulator.  This section derives
that matching directly from the integration region.

Let $H>0$, $s=H+1$, and $y=H-x+z$.  The dimensionless Dalitz region is
\begin{equation}
 0<x<H,\qquad 0<z<\frac{x(H-x)}{s-x}.
\end{equation}
At either soft corner introduce a variable $e$, which here denotes a
kinematic variable rather than the electric charge, and $r=z/e$.  The two
choices are $e=x$ and $e=y$.  In both cases the Jacobian is $e$ and
\begin{align}
 xy&=e[H+e(r-1)],\qquad
 0<e<H,\qquad 0<r<R(e)=\frac{H-e}{s-e},\\
 R(e)&=q+q_1e+O(e^2),\qquad
 q=\frac{H}{s},\qquad q_1=-\frac{1}{s^2}.
\end{align}
Consider a kernel whose expansion at this corner is
\begin{equation}
 K(e,er)=\frac{a(r)}{e^3}+\frac{b(r)}{e^2}+O(e^{-1}),
 \label{eq:compton-matching-kernel}
\end{equation}
where $a$ and $b$ are smooth on $[0,q]$.  Terms with $a=0$ include the usual
logarithmic soft singularity.  Define the four moments
\begin{align}
 P&=\int_0^q r a(r)\,\mathrm{d}r,\\
 A&=\int_0^q b(r)\,\mathrm{d}r+q_1a(q),\\
 B&=\int_0^q b(r)\ln r\,\mathrm{d}r+q_1a(q)\ln q,\\
 C&=\frac{1}{H}\int_0^q(1-r)a(r)\,\mathrm{d}r.
\end{align}
The moment $A$ includes the displacement of the upper boundary.  Indeed,
after multiplication by the Jacobian, the singular part can be written,
in the sense of integration against smooth functions of $r$, as
\begin{equation}
 eK\,\theta(R(e)-r)
 =\frac{a(r)\theta(q-r)}{e^2}
 +\frac{b(r)\theta(q-r)+q_1a(q)\delta(r-q)}{e}
 +O(1).
 \label{eq:compton-moving-boundary-distribution}
\end{equation}
Here $\theta$ is the Heaviside step function and $\delta$ is the
Dirac delta distribution. The last term includes distributions at $r=q$ with integrable coefficients
in $e$.  It contributes the same finite limit in the two regularizations.

Let $I_z(\zeta)$ be the integral with $z>\zeta>0$.  Let $I_\eta$ be the
integral multiplied by $(xy)^\eta$, initially at sufficiently large
$\operatorname{Re}\eta$ and then continued to zero.  Define their finite
parts by
\begin{align}
 I_z(\zeta)&=\frac{P}{\zeta}-A\ln\zeta+F_z+o(1),\\
 I_\eta&=\frac{A}{\eta}+F_\eta+O(\eta).
\end{align}
The relation between them is
\begin{equation}
 \boxed{F_\eta=F_z+A\ln H-B-C.}
 \label{eq:compton-general-soft-matching}
\end{equation}
Contributions from distinct corners are added.  All quantities here refer
to the unnormalized integral, before the common cross-section prefactor.

To prove the relation, first replace the regulator by $(He)^\eta$.
Integration of the $e^{-1}$ term in
\eqref{eq:compton-moving-boundary-distribution} gives $A/\eta+A\ln H$.
Changing the lower boundary from $e>\zeta/r$ to an $r$-independent lower
boundary gives the logarithmic moment $B$.  The $e^{-2}$ term gives
$P/\zeta$, with no additional finite shift for a fixed upper boundary.
The remaining regulator dependence follows from the exact identity for
$xy$:
\begin{equation}
 (xy)^\eta=(He)^\eta
 \left[1-\frac{\eta e(1-r)}{H}
 +O(\eta e^2)+O(\eta^2e^2)\right].
\end{equation}
Its linear term multiplies $a/e^2$.  The resulting power-regulated integral contains
$\eta\int e^{\eta-1}\mathrm{d}e$, whose limit is one.  This gives $-C$.
Terms quadratic in $e$ are integrable at $\eta=0$ and give no finite
regulator shift.  Omitting the linear term would thus give an incorrect
constant even if the residue $A$ were correct.

The corner of the moving boundary does not add another constant.
The lower intersection with $z=\zeta$ satisfies
\begin{equation}
 e_{\min}=\frac{\zeta}{q}-\frac{q_1\zeta^2}{q^3}
 +O(\zeta^3).
\end{equation}
The interval displaced from $\zeta/q$ has width $O(\zeta^2)$.
After the leading angular interval has closed, its integrated density
is at most $O(\zeta^{-1})$, so its contribution is $O(\zeta)$.
Smoothness at $r=0$ likewise makes the excluded strip contribute a
vanishing remainder after the displayed singular terms have been removed.
This establishes that the boundary term already present in $A$ and $B$
is sufficient.

\paragraph{A logarithmic test kernel.}
For $K=x^{-2}$, $a=0$, $b=1$, and $A=q$.  The exact analytically regulated integral is
\begin{equation}
 I_\eta^{(2)}=
 H^{2\eta}\frac{\Gamma(\eta-1)\Gamma(\eta+1)}{\Gamma(2\eta+1)}
 \left[{}_2F_1(\eta+1,\eta-1;2\eta+1;q)-1\right].
\end{equation}
This follows by first integrating $y^\eta$ between
$H-x$ and $s(H-x)/(s-x)$ and then using the Euler beta integral.
Let
\begin{equation}
 x_\pm=\frac{H+\zeta\pm\sqrt{(H-\zeta)^2-4\zeta}}{2}.
\end{equation}
The independent cutoff integral and the two finite parts are
\begin{align}
 I_z^{(2)}&=q\ln\frac{x_+}{x_-}
 +\frac{1}{s}\ln\frac{s-x_+}{s-x_-}
 +\zeta\left(\frac{1}{x_+}-\frac{1}{x_-}\right),\\
 F_z^{(2)}&=2q\ln H-\ln s-q,\\
 F_\eta^{(2)}&=2q\ln H-\frac{\ln s}{s}.
\end{align}
Since $B=q(\ln q-1)$, these expressions obey
\eqref{eq:compton-general-soft-matching}.  In particular, the finite
$-q$ in the cutoff integral cannot be discarded.

\paragraph{A power-divergent test kernel.}
For $K=x^{-3}$, $a=1$, $b=0$, and the residue comes entirely from the
moving boundary: $A=-1/s^2$.  Direct integration gives
\begin{align}
 I_\eta^{(3)}&=
 H^{2\eta-1}\frac{\Gamma(\eta-2)\Gamma(\eta+1)}{\Gamma(2\eta)}
 \left[{}_2F_1(\eta+1,\eta-2;2\eta;q)-1\right],\\
 I_z^{(3)}&=J(x_+)-J(x_-),\\
 J(x)&=-\frac{q}{x}-\frac{\ln x}{s^2}
 +\frac{\ln(s-x)}{s^2}+\frac{\zeta}{2x^2}.
\end{align}
Their expansions yield
\begin{align}
 P&=\frac{q^2}{2},\qquad B=-\frac{\ln q}{s^2},\qquad
 C=\frac{q-q^2/2}{H},\\
 F_z^{(3)}&=-\frac{1}{s}-\frac{2\ln H}{s^2},\\
 F_\eta^{(3)}&=-\frac{1}{s}-\frac{2\ln H+\ln s}{s^2}
 -\frac{q-q^2/2}{H}.
\end{align}
Both the boundary term and $C$ are necessary for agreement.

\paragraph{The Bose subtraction correction.}
For the correction $D=R_{VW}[T_{UW}(y,z)-T_{VW}(x,z)]$ constructed from the
rational coefficients and their local subtraction terms, the $x$ corner
is integrable.  At the $y$ corner its leading coefficient is
\begin{equation}
 a(r)=\frac{16(H+1)(Hr+2)(Hr-2H-6)}{3H(r-1)^3}.
\end{equation}
The coefficient $b(r)$ is obtained by differentiating $e^3D$ at $e=0$,
after substituting $y=e$, $z=er$, and $x=H+e(r-1)$.
This is a finite algebraic operation on the exact rational subtraction.
The only poles of these coefficients in $r$ are at $1$ and $-4/H$,
outside $[0,q]$.  Their moments are therefore closed logarithms and
dilogarithms.  For example, every repeated pole follows by differentiation
with respect to its location $u$:
\begin{align}
 \int_0^q\frac{\mathrm{d}r}{(r-u)^n}
 &=\frac{\partial_u^{n-1}}{(n-1)!}\ln(1-q/u),\\
 \int_0^q\frac{\ln r\,\mathrm{d}r}{(r-u)^n}
 &=\frac{\partial_u^{n-1}}{(n-1)!}
 \left[\ln q\ln(1-q/u)+\operatorname{Li}_2(q/u)\right].
\end{align}
The arguments make both expressions real for these pole locations.
\paragraph{Integrating the rational correction.}
Write the exact partial fraction decomposition of $D$ in $x$ as
\begin{equation}
 D(x,z)=\sum_{a,n}\frac{c_{a,n}(z)}{(x-a(z))^n}.
\end{equation}
There are nine terms, with $(a,n)$ equal to
\begin{gather}
 (H+z,1),\ (H+z,2),\ (H,1),\ (H,2),\ (H,3),\
 (z,1),\ (z,2),\ (a_*,1),\ (a_*,2),\\
 a_* = \frac{-4H+z^2-4z}{z-4}.
\end{gather}
The coefficients are fixed algebraically by the residues of $D$, including
 derivatives for repeated poles.  Thus no fit is involved.  Rationalize
the physical boundary by
\begin{equation}
 z=s+1-u-\frac{s}{u},\qquad
 x_-=s-u,\qquad x_+=s-\frac{s}{u},\qquad \sqrt{s}<u<s.
\end{equation}
The inner integral is explicitly
\begin{align}
 \int_{x_-}^{x_+}D\,\mathrm{d}x
 &=\sum_a c_{a,1}(z)\ln\frac{x_+-a}{x_--a}\\
 &\quad+\sum_{a,n>1}\frac{c_{a,n}(z)}{1-n}
 \left[(x_+-a)^{1-n}-(x_--a)^{1-n}\right].
\end{align}
The ratios are positive in the physical interval, as follows directly
from the positions of the poles relative to $[x_-,x_+]$.
After multiplication by $1-s/u^2$, the outer integrand consists of rational
functions multiplying the following three real logarithms:
\begin{align}
 L_1&=\ln\frac{u^2}{s},\\
 L_2&=\ln\frac{u(u-1)}{s-u},\\
 L_3&=\ln\frac{(u-1)(u+s)^2}{u(u+1)^2(s-u)}.
\end{align}
Each logarithm vanishes at $u=\sqrt{s}$.  Partial fractions followed by
integration by parts therefore gives a finite sum of generalized
polylogarithms of weight at most two, based at $\sqrt{s}$, with rational
coefficients.  The letters are the zeros of the rational denominators and
of the logarithm arguments.  Differentiating that finite sum checks the
original outer integrand term by term.

The power divergence also affects the conversion of the outer endpoint.
For $b=s-u$, the exact relation gives
\begin{equation}
 \zeta=\frac{H}{s}b-\frac{b^2}{s^2}+O(b^3),\qquad
 \frac{P}{\zeta}=\frac{Ps}{Hb}+\frac{P}{H^2}+O(b).
\end{equation}
Consequently, if $F_b$ is the finite part after subtracting the
$b^{-1}$ and $\ln b$ terms, then
\begin{equation}
 F_z=F_b+A\ln(H/s)-\frac{P}{H^2}.
\end{equation}
The last constant is required before applying
\eqref{eq:compton-general-soft-matching}.  Convergent Taylor series for the
endpoint remainder and for the generalized polylogarithms evaluate this
closed result without quadrature.

\subsection{Cancelling a quadratic pole family before integration}
\label{sec:comp-real-kernel-collection}
The separate primitive constructions above also provide a way to check
individual sectors. For the total, a shorter route combines their inner
integrals before constructing the outer primitives. The four contributions
use the same coordinate $u$, the same interval $\sqrt{s}<u<s$, and the
same relation $z=s+1-u-s/u$ to the original endpoint. They can therefore
be added at any fixed positive cutoff in $z$. This observation permits
cancellations before either a generalized polylogarithm or an endpoint
constant is evaluated.

For $A=R_0,D$, denote the rational coefficients of the corresponding
inner integral, including the Jacobian $1-s/u^2$, by $K_{A,j}$, with
$j=0,1,2,3$. Thus its outer integrand is
$K_{A,0}+\sum_{j=1}^3K_{A,j}L_j$, with the real logarithms $L_j$
already defined above. Write the energy-logarithm sector as
$\ell(K_{\log,0}+K_{\log,1}L_1+K_{\log,2}L_2)$, where
$\ell=\ln s$, and the single-$W$ sector as $A_WL_1+B_WL_1L_2$.
Here $A_W$ and $B_W$ are its two rational coefficients after the same
change of variable. The combined integrand is exactly
\begin{align}
 f_4(u)&=\overline K_0+\overline K_1L_1+
             \overline K_2L_2+\overline K_3L_3+B_WL_1L_2,\\
 \overline K_0&=K_{R_0,0}+K_{D,0}+\ell K_{\log,0},\\
 \overline K_1&=K_{R_0,1}+K_{D,1}+\ell K_{\log,1}+A_W,\\
 \overline K_2&=K_{R_0,2}+K_{D,2}+\ell K_{\log,2},\\
 \overline K_3&=K_{R_0,3}+K_{D,3}.
 \label{eq:comp-real-combined-kernel}
\end{align}
All these coefficients follow from the preceding rational partial
fractions and inner boundary differences.

Consider the polynomial
\begin{equation}
 Q_b(u)=u^2+2u+s=(u+1)^2+s-1.
\end{equation}
For $s>1$ it is irreducible over the real numbers. Its two roots
$b_\pm=-1\pm\ii\sqrt{s-1}$ would introduce an additional pair of GPL
letters if the sectors were integrated separately. Polynomial division
shows that their entire pole family cancels. To specify this statement
precisely, let $\nu_b(K)$ be the multiplicity of $Q_b$ in the reduced
denominator of a rational coefficient $K$. The inner integrals give
\begin{equation}
\begin{array}{c|ccc}
 j&\nu_b(K_{R_0,j})&\nu_b(K_{D,j})&
                 \nu_b(K_{R_0,j}+K_{D,j})\\ \hline
 0&2&2&0\\
 1&0&0&0\\
 2&3&3&0\\
 3&3&3&0
\end{array}.
 \label{eq:comp-real-b-pole-orders}
\end{equation}
The cancellation can be checked without factoring $Q_b$ into its roots.
For each row, put the two coefficients over a common denominator
$Q_b^n d$, where $d$ has no factor $Q_b$, and write their numerator
sum as $N$. Euclidean division gives
\begin{equation}
 N=Q_b^n M+R,\qquad \deg_uR<2n,\qquad R=0.
\end{equation}
Cancelling $Q_b^n$ then gives $M/d$. This is an exact identity over
the rational functions of $H=s-1$. The energy-logarithm and single-$W$
coefficients have no $Q_b$ factor, so none is reintroduced in
\eqref{eq:comp-real-combined-kernel}.

The preferred outer integration applies partial fractions to these
combined coefficients. Products of logarithms are converted to GPL
words by the shuffle identity, and repeated poles are reduced by the
integration-by-parts operations already given. The letters $b_\pm$
need never be introduced, and their rational pole terms and six
weight-two functions need never be evaluated separately. This reduction
also leaves no boundary constant to determine for the removed group:
its integrand vanishes identically at every common finite cutoff.
Linearity then preserves the identity when the endpoint finite parts
and their previously derived cutoff conversions are taken. The exact
coefficient construction and the polynomial divisions are implemented
in \nolinkurl{compton_real_combined_kernel.py} using only the inner
rational kernels.

\section{Completion of the pair sectors}
\subsection{The final coordinate of the UV difference}
\label{sec:comp-uv-final-coordinate}
Throughout this calculation the electron mass is the unit of energy,
$s=P^2>1$, $P=p+k_{\mathrm{in}}$ is the incoming total momentum,
with $k_{\mathrm{in}}$ the incident photon momentum, and $H=s-1$ is
fixed and positive. The original Dalitz variables are
\begin{equation}
 x=2P\cdot k_1,\qquad y=2P\cdot k_2,\qquad
 z=x+y-H=2k_1\cdot k_2.
\end{equation}
The final electron is on shell. Positive final energies and the
three-body kinematic boundary give
\begin{equation}
 0<x<H,\qquad 0<z<\frac{x(H-x)}{s-x},\qquad y=H-x+z.
\end{equation}
The hyperbolic angle in the $UV$ angular function is rationalized by
its positive exponential $v$. Explicitly,
\begin{align}
 D_v&=H^2v-x(v-1)^2,\qquad
 z=\frac{x(H-x)(v-1)^2}{D_v},\\
 y&=\frac{H^2v(H-x)}{D_v},\qquad
 \left|\frac{\partial(x,z)}{\partial(x,v)}\right|
 =\frac{x(H-x)H^2(v^2-1)}{D_v^2}.
\end{align}
At fixed $x$, the lower and upper Dalitz boundaries map to $v=1$ and
$v=s$. The Jacobian is positive in the interior. Thus the new domain
is the rectangle $0<x<H$, $1<v<s$. This is why the inner energy
integration can be completed before the final $v$ integration.
In particular,
\begin{equation}
 xy=\frac{x(H-x)}{1-kx/H},\qquad k=\frac{(v-1)^2}{Hv},
 \qquad 0\leq k\leq\frac Hs<1.
\end{equation}
The analytic regulator therefore has the positive real value
$[x(H-x)]^\eta(1-kx/H)^{-\eta}$. The denominator is at least $1/s$
on the closed rectangle. No logarithmic branch is inferred from a
numerical result.

After the inner energy integration, the $UV$ angular difference sector
is a finite sum of terms
\begin{equation}
 R_{ab}(H,v)\ln^a v\ln^b(1-k),\qquad
 k=\frac{(v-1)^2}{Hv},\qquad 1\leq v\leq H+1.
\end{equation}
Here $H=s-1>0$, and the $R_{ab}$ are the rational functions obtained
from the real-emission amplitude. Their coefficients may include
$\ln H$, which is constant in this integration. The pole and finite
parts of the analytic regulator are treated separately.

\paragraph{The endpoint Laurent coefficients.}
The separated rational energy kernels have poles of at most second
order at $x=0,H$. Other poles are at $x=H/k>H$ and hence outside the
integration path. Start with $\operatorname{Re}\eta>1$, where these
endpoint integrals converge absolutely. Subtract the constant and linear
Taylor terms of the smooth endpoint factor before analytically continuing
to a neighborhood of $\eta=0$. This operation gives an ordinary convergent
remainder and explicit Mellin integrals for the removed terms.
For example, near $x=0$ the smooth factor obeys
\begin{equation}
 (H-x)^\eta(1-kx/H)^{-\eta}
 =H^\eta\left[1+\eta\frac{k-1}{H}x+O(x^2)\right].
\end{equation}
The kernel $x^{-2}$ consequently has the finite coefficient
$-1/H+(k-1)/H=(k-2)/H$. The first term is the finite part of the
unregulated double pole. The second follows from $\eta x$ multiplying
$x^{-2+\eta}$. Although that term vanishes pointwise as $\eta\to0$,
its endpoint integral has a nonzero limit.
At $x=H$, write $b=H-x$. The corresponding expansion is
\begin{equation}
 x^\eta(1-kx/H)^{-\eta}
 =\left(\frac H{1-k}\right)^\eta
 \left[1-\frac{\eta b}{H(1-k)}+O(b^2)\right].
\end{equation}
Including the signs in $(x-H)^{-1}=-b^{-1}$ gives the complete rules
\begin{align}
 \int_0^H\frac{[x(H-x)]^\eta}{x(1-kx/H)^\eta}\,\dd x
 &=\frac1\eta+2\ln H+O(\eta),\\
 \int_0^H\frac{[x(H-x)]^\eta}{x^2(1-kx/H)^\eta}\,\dd x
 &=\frac{k-2}{H}+O(\eta),\\
 \int_0^H\frac{[x(H-x)]^\eta}{(x-H)(1-kx/H)^\eta}\,\dd x
 &=-\frac1\eta+\ln(1-k)-2\ln H+O(\eta),\\
 \int_0^H\frac{[x(H-x)]^\eta}{(x-H)^2(1-kx/H)^\eta}\,\dd x
 &=-\frac{2-k}{H(1-k)}+O(\eta).
\end{align}
The displayed integrals near $\eta=0$ mean their analytic continuations
from the convergence half-plane. The other energy kernels are bounded
on $0\leq x\leq H$ and can be integrated at $\eta=0$ for their finite
coefficient. These rules explain every regulator-generated constant used
in the inner integration. They do not discard power-divergent endpoint
terms by hand.

Put $\zeta=v-1$ for the final polylogarithm coordinate. This $\zeta$ is
distinct from the original Dalitz variable $z$. The two roots of $1-k$ are
\begin{equation}
 r_\pm=\frac{H+2\pm\sqrt{H(H+4)}}2,\qquad
 r_+r_-=1,\qquad 0<r_-<1<H+1<r_+.
\end{equation}
Factorization, together with the value at $\zeta=0$, gives
\begin{align}
 \ln v&=G(-1;\zeta),\\
 \ln(1-k)&=G(r_+-1;\zeta)+G(r_--1;\zeta)-G(-1;\zeta).
\end{align}
The normalization of the multiple polylogarithms is
\begin{align}
 G(\varnothing;\zeta)&=1,\qquad
 G(a,\boldsymbol b;\zeta)
 =\int_0^\zeta\frac{\dd u}{u-a}G(\boldsymbol b;u),\\
 G(\underbrace{0,\ldots,0}_{n};\zeta)&=\frac{\ln^n\zeta}{n!}.
\end{align}
For words with a convergent lower endpoint, the path is the real segment.
The last line fixes the normalization for logarithmic lower endpoints.
All logarithms are real on $0\leq \zeta\leq H$. Their letters lie outside
this interval. Products are expanded by the shuffle identity, and the
rational factors are integrated by partial fractions and the finite
polylogarithmic reduction already described for the soft angular term.

\paragraph{A removable origin.}
Individual rational coefficients can have poles at $\zeta=0$, even though
their sum multiplied by the logarithms is regular. An origin pole must
therefore be handled before setting the lower endpoint to zero. For an
analytic word write
\begin{equation}
 G(\boldsymbol a;\zeta)=\sum_{j\geq1}g_j(\boldsymbol a)\zeta^j.
\end{equation}
Normalize each primitive to have zero finite part at $\zeta=0$, meaning
zero constant term after separating its negative powers and logarithms.
For $\boldsymbol a=(b,\boldsymbol c)$ and $n\geq2$, integration by
parts then gives
\begin{align}
 \mathfrak T[\zeta^{-n};b,\boldsymbol c](\zeta)
 &=-\frac{G(b,\boldsymbol c;\zeta)}{(n-1)\zeta^{n-1}}
 +\frac{g_{n-1}(b,\boldsymbol c)}{n-1}\\
 &\quad+\frac1{n-1}\mathfrak T\!\left[
 \frac1{\zeta^{n-1}(\zeta-b)};\boldsymbol c\right](\zeta).
\end{align}
The displayed constant is essential. It subtracts the finite part of
the boundary term and is absent in the simpler case of a nonzero pole.
For the empty word the origin rules are
\begin{equation}
 \mathfrak T[\zeta^{-1};\varnothing](\zeta)=\ln \zeta,\qquad
 \mathfrak T[\zeta^{-n};\varnothing](\zeta)=\frac{\zeta^{1-n}}{1-n}
 \quad(n\geq2).
\end{equation}
Linearity then makes the normalized primitive of the full regular
integrand vanish at zero. Negative powers and $\ln \zeta$ must cancel in
the complete expression. This is a check on the independently derived
rational coefficients, not a choice of a finite subtraction constant.

An elementary example displays the role of the constant without any
numerical evaluation,
\begin{equation}
 \int_0^Z\frac{\ln(1+\zeta)-\zeta+\zeta^2/2}{\zeta^3}\,\dd\zeta
 =\left(\frac12-\frac1{2Z^2}\right)\ln(1+Z)
 +\frac1{2Z}-\frac14.
\end{equation}
The right-hand side differentiates to the integrand and tends to zero
as $Z\to0^+$. The $-1/4$ is the finite-part normalization constant.

The script \texttt{compton\_soft\_uv\_check.py} implements this reduction.
It reads the two independently generated inner-energy coefficient files,
forms their difference, and evaluates the resulting ordinary multiple
polylogarithms by convergent power series and analytic continuation.
The original coordinate integral is used only as a separate check.
For general $H>0$, direct expansion of the rational coefficients gives
\begin{align}
 \mathcal P_{\mathrm{UV}}(H,1+\zeta)
 &=-\frac{256(H+1)(H+5)}{15H^2}\zeta+O(\zeta^2),\\
 \mathcal F_{\mathrm{UV}}(H,1+\zeta)
 &=-\frac{256(H+1)}{15H^3}
 \bigl[2H(H+5)\ln H-H^2-5H-15\bigr]\zeta
 +O(\zeta^2).
\end{align}
These follow by expanding $\ln(1+\zeta)$ and
$\ln[1-\zeta^2/(H(1+\zeta))]$ before multiplying the rational Laurent
series. The coefficients of $\zeta^{-4},\ldots,\zeta^{-1},\zeta^0$
vanish identically in both sums. Only finitely many Taylor coefficients
are required to prove these cancellations, since the highest rational
origin pole is fourth order.
Thus the apparent origin poles cancel for every $H>0$.
\subsubsection{Constant letters for the UV contribution}
The moving letters of the integrated $UV$ contribution can be removed
without using the total cross section.  Introduce
\begin{equation}
 q=\frac{\sqrt{H+4}-\sqrt H}{\sqrt{H+4}+\sqrt H},\qquad
 H=\frac{(1-q)^2}{q},\qquad 0<q<1.
\end{equation}
Write $z=v-1$.  Its upper endpoint is $H$, and the two logarithms of the
inner Laurent coefficient become
\begin{align}
 \ln v&=G(-1;z),\\
 \ln(1-k)&=G(q-1;z)+G(q^{-1}-1;z)-G(-1;z),
 \qquad k=\frac{z^2}{H(1+z)}.
\end{align}
This follows by factoring $1-k$ and matching its value one at $z=0$.
All factors are positive on $0\leq z\leq H$: the negative letter $q-1$
is outside the interval, and $q^{-1}-1>H$.  Rational partial fractions
and integration by parts therefore produce GPL words with letters
$0,-1,q-1,q^{-1}-1$ and rational coefficients in $q,z$.

To replace their moving endpoint and letters by constants, differentiate
the whole word with respect to $q$.  For a word $w=(a_1,\ldots,a_n)$,
let $w\setminus j$ mean deletion of its $j$th entry.  The total differential
is
\begin{align}
 \dd G(w;H)={}&G(w\setminus1;H)\,\dd\ln(H-a_1)
 -G(w\setminus n;H)\,\dd\ln(-a_n)\\
 &+\sum_{j=1}^{n-1}
 [G(w\setminus(j+1);H)-G(w\setminus j;H)]
                      \,\dd\ln(a_j-a_{j+1}).
\end{align}
The last endpoint term is absent when $a_n=0$.  Adjacent equal letters
are treated by continuity, with their paired difference equal to zero.
Factoring each logarithmic differential gives only constant simple poles.
In this case the complete alphabet is
\begin{equation}
 \left\{0,1,-1,\frac{1+\ii\sqrt3}{2},
                    \frac{1-\ii\sqrt3}{2}\right\}.
\end{equation}
For example, $H+1=(q^2-q+1)/q$ explains the complex-conjugate pair.
Its appearance does not change the real physical boundary value.

The recursive integration uses basepoint $q=1$ approached from below.
For every nonzero moving letter, $H/a_j\to0$ there.  Thus every original
word containing a nonzero letter vanishes, including powers of logarithms
multiplying that vanishing limit.  Pure-zero words instead have
\begin{equation}
 G(\underbrace{0,\ldots,0}_{n};H)=\frac{\ln^nH}{n!},\qquad
 \ln H=2\ln(1-q)-\ln q.
\end{equation}
These relations fix all integration constants.  The recursion checks that
pure powers of $\ln(1-q)$ cancel for every nonzero-letter original word.

The lower endpoint of the angular primitive supplies an additional finite
constant, even though the combined integral is convergent.  With
\begin{align}
 P(q)&=3q^8+18q^7+12q^6-1874q^5+1666q^4
                         -1874q^3+12q^2+18q+3,\\
 Q(q)&=-25q^8+238q^7-436q^6+2274q^5-4294q^4
                         +2274q^3-436q^2+238q-25,
\end{align}
the constants to subtract from the upper primitive are
\begin{align}
 b_{UV,-1}&=-\frac{(q^2-q+1)P(q)}{9q^3(q-1)^4},\\
 b_{UV,0}&=-\frac{q^2-q+1}{18q^3(q-1)^4}
                         [4P(q)\ln H+Q(q)].
\end{align}
The subscripts label the pole and finite Laurent coefficients, respectively.
These constants follow from the Taylor coefficients of each GPL multiplied
by its rational prefactor at $z=0$.  Setting every GPL to zero before
multiplying by the prefactor would lose them.

The resulting pole and finite coefficients contain respectively seven and
41 constant-letter words of weight at most three, before multiplication
by $1/(8sH)$. The script \nolinkurl{compton_soft_uv_constant_basis.py}
generates their explicit coefficient lists,
\nolinkurl{soft_uv_constant_basis_pole.json} and
\nolinkurl{soft_uv_constant_basis_finite.json}, in
\nolinkurl{build/intermediate/}.
The reduction script uses only the independent inner Laurent coefficients.
The primitive follows the exact rational integration recurrence above.
The primitive is checked by differentiating each independent kernel
term before summing. If its coefficient dictionary is $Q_w(y)$,
the derivative follows from
\begin{equation}
 \frac{\mathrm d}{\mathrm dy}\bigl[Q_w(y)G(w;y)\bigr]
 =Q'_w(y)G(w;y)+\frac{Q_w(y)}{y-a}G(\widehat w;y),
 \qquad w=(a,\widehat w).
\end{equation}
For the empty word only the first term occurs. Subtracting the input
kernel leaves a rational coefficient for each GPL word.
Every coefficient vanishes exactly, for both pole and finite parts.
The calculation holds the external parameter $q$ and $\ln H$ fixed
while differentiating in $y=v-1$.
Canceling rational factors before differentiation makes this test finite
in size. Linearity then proves the derivative identity for their sum.
The reproducible checks are in
\texttt{compton\_uv\_primitive\_certificate.py}.
Evaluation against the earlier moving-letter expression provides
a separate numerical check of the conversion.


\subsection{The UW endpoint matching}
Use the original Dalitz coordinates defined above, and set
$a=H-x$. The positive angular variable satisfies $1<v<1+a$, so
$t=(v-1)/a$ places its two endpoints at zero and one. The complete map is
\begin{align}
 x&=H-a,\qquad v=1+at,\qquad
 z=\frac{a(H-a)t^2}{1+at},\\
 y&=\frac{a[1+Ht^2+at(1-t)]}{1+at},\\
 \left|\frac{\partial(x,z)}{\partial(a,t)}\right|
 &=\frac{a(H-a)t(2+at)}{(1+at)^2}>0.
\end{align}
The new domain is exactly $0<a<H$, $0<t<1$. Both the sign of the
Jacobian and the range of $t$ follow from the positive solution for $v$.
In these rectangular coordinates the analytic
regulator is $[a(H-a)f(a,t)]^\eta$, with
\begin{equation}
 f(a,t)=\frac{1+Ht^2+at(1-t)}{1+at},\qquad
 0<a<H,\quad 0<t<1.
\end{equation}
The useful alternative form $f=1+(H-a)t^2/(1+at)$ shows that
$1\leq f\leq1+H$ on the rectangle. Its logarithm and endpoint derivatives
are real and bounded for fixed $H>0$. For a fixed interior $t$, start
again at $\operatorname{Re}\eta>1$ and continue analytically only after
subtracting the endpoint Taylor terms.
Write the unregulated difference as
$F(a,t)=L(a,t)\ln(1+at)-T(a,t)$, where $L$ and $T$ are the rational
coefficient and Taylor subtraction obtained after including the Jacobian.
At the upper endpoint put $b=H-a$ and expand
\begin{equation}
 L=\frac{L_2}{b^2}+\frac{L_1}{b}+O(1),\qquad
 T=\frac{T_2}{b^2}+\frac{T_1}{b}+O(1).
\end{equation}
These coefficients are fixed algebraically, for example
$L_2=[b^2L]_{b=0}$ and $L_1=[\partial_b(b^2L)]_{b=0}$.
Because $\ln[1+(H-b)t]=\ell-bt/(1+Ht)+O(b^2)$, the complete singular
coefficients are
\begin{align}
 g_2&=L_2\ell-T_2,\qquad \ell=\ln(1+Ht),\\
 g_1&=L_1\ell-\frac{L_2t}{1+Ht}-T_1.
\end{align}
In particular, the derivative of the logarithm contributes to the simple
pole and cannot be omitted. Thus $F=g_2/b^2+g_1/b+O(1)$.
Since $f(H,t)=1$, direct expansion gives
\begin{equation}
 a^\eta f(a,t)^\eta
 =H^\eta\left[1+\eta b\left(-\frac1H+\frac{t^2}{1+Ht}\right)
 +O(b^2)\right].
\end{equation}
The ordinary finite part used for comparison is defined explicitly by
\begin{equation}
 \operatorname{FP}\int_0^H F(a,t)\,\dd a
 =\lim_{\delta\to0^+}\left[
 \int_0^{H-\delta}F(a,t)\,\dd a-\frac{g_2}{\delta}
 +g_1\ln\delta\right].
\end{equation}
After the two singular terms are subtracted, the remainder is bounded
at $b=0$. Its integral is analytic in $\eta$ near zero. The removed terms
are integrated before expansion using
\begin{equation}
 \int_0^H b^{\eta-2}\,\dd b=\frac{H^{\eta-1}}{\eta-1},
 \qquad
 \int_0^H b^{\eta-1}\,\dd b=\frac{H^\eta}{\eta},
\end{equation}
first in their convergence half-planes and then by continuation.
The double-pole kernel therefore does not by itself produce a pole at
$\eta=0$. Its term proportional to $\eta b$, however, multiplies the
second integral and supplies a finite constant. The simple pole produces
$g_1/\eta$. Relative to the ordinary finite
part of the $a$ integral, the finite matching contribution is
\begin{equation}
 g_1\ln H+g_2\left(-\frac1H+\frac{t^2}{1+Ht}\right).
\end{equation}
The second term arises because the factor $\eta b$ multiplying
$b^{-2+\eta}$ leaves a finite limit after integration. It must not be
lost by setting $\eta=0$ before integrating. The finite part of the
unregulated integral already includes $-g_2/H$ and $g_1\ln H$.
The matching term is additional to those terms. At $a=0$ the possible
simple-pole coefficient cancels in the complete angular difference.
Explicitly, $\ln(1+at)=at+O(a^2)$ implies
\begin{equation}
 \lim_{a\to0}aF(a,t)
 =t\lim_{a\to0}a^2L(a,t)-\lim_{a\to0}aT(a,t)=0.
\end{equation}
The last equality is an exact identity of the rational coefficients.
Therefore there is no matching term from that endpoint.

The cancellation at the other boundary can also be seen before numerical
evaluation. Expanding $\ln(1+Ht)$ and the explicit rational endpoint
coefficients gives, for every fixed $H>0$,
\begin{align}
 g_2(H,t)&=\frac{64H^2(H+1)}{15}t+O(t^2),\\
 g_1(H,t)&=-\frac{64H(H+1)}5t+O(t^2).
\end{align}
All the coefficients of $t^{-4},t^{-3},t^{-2},t^{-1},t^0$ vanish in
these complete combinations. Thus the functions entering the pole and
finite endpoint matching are integrable at the lower angular boundary.
This algebraic cancellation, rather than their numerical values at a
few points, justifies treating that boundary as removable.

The remaining $t$ integral of $g_1$ is rational times $\ln(1+Ht)$.
The same finite-part polylogarithmic reduction gives seven terms.
\paragraph{Logarithmic lower endpoints.}
For the finite coefficient, the inner integral also contains $\ln t$
and weight-two functions. A word ending in zero can be rewritten by
shuffle identities. In particular,
\begin{equation}
 G(a,0;t)=\ln t\,G(a;t)-G(0,a;t).
\end{equation}
For a longer word the same identity recursively removes its trailing
zeros. The result is a polynomial in $\ln t$ whose coefficients are
regular-origin polylogarithms. In the repeated-pole primitive, the
finite-part constant is the coefficient of $t^{n-1}(\ln t)^0$.
Terms containing $\ln t$ are therefore included in the reduction rather
than incorrectly treated as ordinary Taylor coefficients.
For example, the shuffle identity gives the exact check
\begin{equation}
 \int_0^1\frac{\ln t}{1+t}\,\dd t
 =G(-1,0;1)=-G(0,-1;1)=\operatorname{Li}_2(-1)
 =-\frac{\pi^2}{12}.
\end{equation}
Differentiation proves the primitive, and its lower endpoint vanishes.

For the finite inner result, use the following names for its eight
logarithmic and dilogarithmic functions:
\begin{align}
 \ell&=\ln(1+Ht),\qquad l_H=\ln H,\qquad
 l_2=\ln(2+Ht),\\
 l_P&=\ln(1+Ht-Ht^2),\qquad l_t=\ln t,\\
 d_1&=\operatorname{Li}_2(-Ht),\qquad
 d_2=\operatorname{Li}_2(-1-Ht)+\frac{\pi^2}{12},\\
 d_3&=\operatorname{Li}_2\!\left(-\frac{(1-t)(1+Ht)}t\right)
       -\operatorname{Li}_2\!\left(-\frac{1-t}t\right).
\end{align}
All logarithm arguments are positive for $H>0$ and $0<t<1$.
The dilogarithm arguments are nonpositive, so their real branches
follow from the original physical interval. Expanding the boundary
formulas in \eqref{eq:comp-real-UW-inner-closed} in these functions
gives thirteen monomials. Write $R_f(H,t)$ for the rational coefficient
of a monomial $f$, including $R_1$ for the constant monomial.

In the complete $UW$ finite sum, the coefficient of
$\ell\ln t$, with $\ell=\ln(1+Ht)$, is exactly the negative of the
coefficient of the dilogarithm difference $d_3$.
This is an identity of the rational coefficients for arbitrary $H$ and
$t$, not a numerical cancellation. It groups these terms into
$d_3-\ell\ln t$. Writing $a=-1/H$ and
$t_\pm=(1\pm\sqrt{1+4/H})/2$, the group is
\begin{equation}
 d_3-\ell\ln t
 =-\sum_{r=t_\pm}\left[G(1,r;t)+G(a,r;t)-G(0,r;t)\right]
 -G(0,a;t).
\end{equation}
Every word on the right has a regular lower endpoint. This identity
also explains why the apparent $\ln t$ terms disappear from the
complete lower-endpoint expansion. The general finite-part rules remain
useful when these terms are integrated separately.

\paragraph{The finite UW coordinate integral.}
Write the inner finite result as a finite sum
$\sum_j R_j(t)M_j(t)$, where the $R_j$ are its explicit rational
coefficients and each $M_j$ is a monomial in the eight logarithmic or
dilogarithmic functions. Translating these functions to ordinary
polylogarithms and applying shuffle products gives finite coefficients
$c_{j\boldsymbol a}$ such that
\begin{equation}
 M_j(t)=\sum_{\boldsymbol a}c_{j\boldsymbol a}
 G(\boldsymbol a;t).
\end{equation}
The final coordinate integral is therefore the finite endpoint expression
\begin{equation}
 \mathcal F_{UW}(H)=\lim_{Z\to1^-}
 \sum_{j,\boldsymbol a}c_{j\boldsymbol a}
 \mathfrak T[R_j;\boldsymbol a](Z).
 \label{eq:comp-real-UW-final-primitive}
\end{equation}
All rational integrations in this formula are eliminated by the terminating
partial-fraction rules, including the origin constants described above.
The functions with a first letter equal to one are combined before the
endpoint limit. Their separate logarithmic divergences need not cancel
within a single term.

For evaluation, the complete inner finite function is analytic at $t=1$.
Indeed, its logarithm arguments are nonzero there and the two arguments
in the dilogarithm difference both tend to zero. The apparent rational
poles cancel in the summed expression. With $d=1-t$, the remaining tail
is consequently obtained algebraically,
\begin{equation}
 \mathcal F_{UW}^{\mathrm{inner}}(H,1-d)=\sum_{n\geq0}c_n d^n,
 \qquad
 \int_{1-\delta}^1\mathcal F_{UW}^{\mathrm{inner}}(H,t)\,\dd t
 =\sum_{n\geq0}\frac{c_n\delta^{n+1}}{n+1}.
\end{equation}
The coefficients are generated by rational power-series division and the
Taylor series of the displayed logarithms and dilogarithms. No phase-space
quadrature enters this endpoint evaluation. Individual terms near $t=0$ require powers of $t$ multiplied by
powers of $\ln t$. The combined function is analytic there, as follows
from the cancellation of $\ln t$ established above.
The generated finite expression requires 52 nonempty polylogarithm words.
\paragraph{The order of the limiting operations.}
All endpoint statements here concern a fixed $H>0$. They do not justify
interchanging the photon-mass limit or the regulator expansion with the
threshold limit $H\to0$. The analytic regulator is first introduced in
a common convergence half-plane. Only after the endpoint Taylor terms
have been isolated are their Mellin integrals continued to $\eta=0$.
The bounded remainders can then be expanded under the integral.

There is also a uniform reason for the removable angular origins. The
hyperbolic function entering each angular integral is analytic at
$q=0$, where
\begin{equation}
 q=\cosh(\ln v)-1=\frac{(v-1)^2}{2v}.
\end{equation}
Together with the Taylor expansion of its rational prefactors, its
constant and linear terms give precisely the terms removed in $J-T$. For the $UV$ map, $q=\zeta^2/[2(1+\zeta)]$ is independent
of the inner energy coordinate. For the $UW$ map,
$q=a^2t^2/[2(1+at)]$ is uniformly small as $t\to0$ for $0\leq a\leq H$.
Taylor's remainder is therefore uniform after the explicit soft endpoint
powers have been factored out. Those powers are at most second order,
and are controlled by $\operatorname{Re}\eta>1$ before continuation.
The local Mellin subtractions above preserve this uniform angular
behavior because their coefficient functions and the regulator shape
factors are analytic in the same angular neighborhood.
This gives the reason for taking the Laurent coefficients before the
last angular integration. A pointwise expansion of the un-subtracted
endpoint integrand would not give that justification.

At the upper $UW$ boundary, the logarithm arguments are positive and
the two $d_3$ dilogarithm arguments tend to zero. At nonzero electron mass, the original hard angular denominators have
no pole on the integration sphere for positive photon energies. The
soft endpoint factors have already been isolated explicitly. Thus the
complete soft-subtracted energy integral has an analytic neighborhood
of $t=1$, even when a particular angular partial-fraction representation
becomes singular there.
The largest apparent rational pole in the separated coefficient list is
fifth order. To determine the tail through order $d^n$, the atom series
are therefore expanded through at least order $d^{n+5}$ before summing.
All negative powers are canceled before the endpoint value is used.
The resulting convergent power series may be integrated term by term
inside its disk of convergence. These arguments establish the endpoint
operations. The numerical verification chapter
tests their implementation independently of these analytic arguments.

\subsubsection{Exact removal of the finite UW endpoint poles}
\label{sec:comp-UW-upper-certificate}
The cancellation of the apparent poles at $t=1$ can be checked for
arbitrary $H>0$, without evaluating the remaining integral numerically.
Write $d=1-t$, $E=\ln(1+H)$, and $C=(H+1)/H$ in this subsection.
Here $E$ is a logarithm, not an electron energy.
Denote the rational coefficient multiplying an atom $f$ in the
thirteen-term inner result by $R_f$.  Only six terms can have a pole
at this endpoint:
\begin{equation}
 R_1+R_{d_3}d_3+R_{l_P}l_P
       +R_\ell\ell+R_{\ell l_t}\ell l_t+R_{\ell l_P}\ell l_P.
\end{equation}
Their rational pole orders are respectively $4,3,5,2,3,3$.
The remaining seven terms have rational denominators nonzero at $t=1$
and analytic atoms there.  Thus they cannot contribute a negative power
of $d$.

The required Taylor coefficients follow from
\begin{align}
 \ell&=E-\sum_{n=1}^4\frac1n
                   \left(\frac{H}{1+H}\right)^n d^n+O(d^5),\\
 l_t&=-\sum_{n=1}^4\frac{d^n}{n}+O(d^5),\\
 l_P&=\sum_{n=1}^4\frac{(-1)^{n+1}}n
                         (Hd-Hd^2)^n+O(d^5),\\
 d_3&=\sum_{n=1}^4\frac{z_1^n-z_2^n}{n^2}+O(d^5),
 \qquad z_1=-\frac{d(1+H-Hd)}{1-d},\quad z_2=-\frac d{1-d}.
\end{align}
Terms of degree above four in the last two finite sums are simply omitted
when constructing the Taylor polynomial.  The two dilogarithm arguments
vanish at $d=0$, so their ordinary convergent series applies.  All these
statements are at fixed $H>0$ in a sufficiently small neighborhood of
$d=0$.  There is no continuation through a dilogarithm cut.

For completeness, the rational expansion requires only polynomial
arithmetic.  If $R_f(1-d)=(-d)^{-m}N(d)/D(d)$ with $D(0)\ne0$, write
$N=\sum N_jd^j$, $D=\sum D_jd^j$.  Its regular quotient coefficients obey
\begin{equation}
 r_0=\frac{N_0}{D_0},\qquad
 r_n=\frac{N_n-\sum_{j=1}^nD_jr_{n-j}}{D_0}.
\end{equation}
Multiplying this finite quotient series by the atom series determines
every negative power.  The largest pole is fifth order, so degree four
is sufficient.  No infinite-series truncation error enters a coefficient
of a negative power.

Let $\operatorname{PP}_{d=0}$ denote the sum of strictly negative powers
in a Laurent expansion.  Substituting the actual rational coefficients
of the inner result gives the particularly simple pair
\begin{align}
 \operatorname{PP}_{d=0}(R_1+R_{l_P}l_P)
 &=C\left[-\frac{16}{d^2}+\frac{4(H-10)}d\right],\\
 \operatorname{PP}_{d=0}(R_{d_3}d_3)
 &=C\left[\frac{16}{d^2}-\frac{4(H-10)}d\right].
\end{align}
The apparent fourth- and third-order poles have already canceled within
the first line.  The fifth-order coefficient vanishes because $l_P(0)=0$.
For the other three terms, divide their principal parts by $8C$.
The coefficients of $d^{-2}$ and $d^{-1}$ are
\begin{center}
\begin{tabular}{lll}
\toprule
Term & $[d^{-2}]/(8C)$ & $[d^{-1}]/(8C)$\\
\midrule
$R_\ell\ell$ & $2E(H+1)/H$ & $-E(H-4-7/H)-2$\\
$R_{\ell l_t}\ell l_t$ & $-2E/H$ & $-7E/H+2/(H+1)$\\
$R_{\ell l_P}\ell l_P$ & $-2E$ & $E(H-4)+2H/(H+1)$\\
\bottomrule
\end{tabular}
\end{center}
Each column sums to zero as a rational identity in $H,E$.  Their
third-order coefficients vanish because $l_t(0)=l_P(0)=0$.
Combining the two cancellations proves
\begin{equation}
 [d^{-j}]F_{UW}(H,1-d)=0,\qquad j=1,\ldots,5,\qquad H>0.
\end{equation}
Here $F_{UW}(H,t)$ denotes the finite inner integrand before the final
$t$ integration.  Since each atom is analytic at the endpoint and the
rational factors have finite pole order, the disappearance of their
complete principal part proves that this endpoint is removable.
It is an identity for the finite integrand, not only for its infrared pole.

The independent script
\nolinkurl{compton_soft_uv_uw_endpoint_certificate.py} reads the thirteen
rational coefficients, computes the six contributing expansions, and
checks every coefficient of $E$ separately using exact arithmetic.
\nolinkurl{uw_endpoint_certificate.json} records the individual terms
as well as the zero sums.  This certificate establishes endpoint
regularity.  It does not provide a bound on the entire omitted numerical
Taylor tail used when evaluating the final special functions.

\subsubsection{Exact removal of the finite UW origin poles}
\label{sec:comp-UW-lower-certificate}
The lower endpoint also admits an exact certificate for every fixed
$H>0$.  The only explicit $\ln t$ in the thirteen-term expression comes
from $d_3$ and the product $\ell\ln t$.  Their rational coefficients
satisfy
\begin{equation}
 R_{d_3}+R_{\ell l_t}=0.
\end{equation}
It is therefore essential to combine these two terms first.  Define
\begin{equation}
 \mathcal D(t)=d_3(t)-\ell(t)\ln t,\qquad \ell(t)=\ln(1+Ht).
\end{equation}
The logarithm cancels exactly, rather than merely becoming integrable.

Its boundary constant follows from the real inversion identity
\begin{equation}
 \operatorname{Li}_2(-z)=-\frac{\pi^2}{6}-\frac12\ln^2z
                              -\operatorname{Li}_2(-1/z),\qquad z>0.
\end{equation}
Apply it to $z=(1-t)(1+Ht)/t$ and $z=(1-t)/t$.  Both are positive on
$0<t<1$, so the same negative-real-axis boundary value is used in both
terms.  The exact result is
\begin{align}
 \mathcal D(t)={}&-\ell\ln(1-t)-\frac12\ell^2
 -\operatorname{Li}_2\!\left(-\frac{t}{(1-t)(1+Ht)}\right)
 +\operatorname{Li}_2\!\left(-\frac{t}{1-t}\right).
\end{align}
Every term on this line is analytic at $t=0$ and vanishes there.
This proves $\mathcal D(0)=0$ independently of any numerical evaluation.
Differentiating either this expression or the original pair gives
\begin{equation}
 \mathcal D'(t)=\frac{l_P}{1-t}-\frac{Hl_P}{1+Ht}
                         +\frac{l_P-\ell}{t},\qquad
 l_P=\ln(1+Ht-Ht^2).
\end{equation}
The apparent $1/t$ is removable because $l_P(0)=\ell(0)=0$.
The resulting series starts as
\begin{equation}
 \mathcal D(t)=-\frac{H^2}{2}t^2
 +\frac{H^2(H+1)}2t^3
 -\left(\frac{11H^4}{24}+\frac{2H^3}{3}\right)t^4+O(t^5).
\end{equation}
In particular, $d_3=Ht\ln t+O(t^2[1+|\ln t|])$ has a zero constant,
fixed by the inversion identity, not chosen to enforce cancellation.

After this combination the entire kernel contains only atoms analytic
at the origin.  Its rational factors have poles of order at most five.
Using the finite quotient recurrence of the preceding subsection,
expand the atom series through degree six.  This determines all negative
powers, the constant, and the coefficient of $t$ exactly.
The cancellation can be displayed without a long list of individual
rational coefficients.  Define the three groups
\begin{align}
 \mathcal B_0&=R_1+R_{l_2}l_2+R_{l_P}l_P+R_\ell\ell,\\
 \mathcal B_1&=R_{d_3}\mathcal D+R_{\ell l_P}\ell l_P+R_{\ell^2}\ell^2,\\
 \mathcal B_H&=\ln H\,[R_{l_H}+R_{\ell l_H}\ell].
\end{align}
Their expansions obtained from the explicit coefficient table are
\begin{align}
 \mathcal B_0={}&(H+1)\left[-\frac{128}{H^2t^2}+
 \frac{64}{Ht}+\frac{4(3H^2-67H-72)}{3H}\right]+O(t),\\
 \mathcal B_1={}&(H+1)\left[\frac{128}{H^2t^2}-
 \frac{64}{Ht}-\frac{4(3H^2-67H-72)}{3H}\right]+O(t),\\
 \mathcal B_H={}&O(t).
\end{align}
Their fifth-, fourth-, and third-order poles have already canceled
inside these groups.  The remaining three atom terms are individually
$O(t)$ because their rational coefficients are regular and their atoms
vanish.  Thus all negative powers and the constant cancel in the full
finite kernel.  The first surviving coefficient is
\begin{equation}
 F_{UW}(H,t)=-\frac{32H(H+1)}{15}(12\ln H-1)t+O(t^2),
                    \qquad t\to0^+,\quad H>0\text{ fixed}.
\end{equation}
There is no residual $\ln t$ in this combined expression.  The origin is
an ordinary analytic endpoint, so the finite lower boundary of the GPL
primitive equals the lower limit of the convergent integral.

The script \nolinkurl{compton_soft_uv_uw_origin_certificate.py} first
checks the exact rational coefficient identity removing $\ln t$.
It then treats $\ln H$ and $\ln2$ as independent symbols and checks
the coefficients of $t^{-5},\ldots,t^{-1},t^0$ separately in exact
arithmetic.  Their values and the nonzero coefficient of $t$ are stored
in \nolinkurl{uw_origin_certificate.json}.  Together with the upper
certificate, this establishes both finite endpoints for arbitrary
$H>0$.  A rigorous error bound for the numerical special-function
series is a separate issue.

\subsection{An exact certificate for the integrated pole sum}
This additional algebraic check tests the separately implemented sector
primitives for every $H>0$.  The physical soft factorization was established
earlier.  No numerical comparison or published total enters this check.
The $UV$ pole density already uses $1<v<H+1$.  In the $UW$ pole density,
set $v=1+Ht$, including $\dd t=\dd v/H$.  The physical infrared angular
integral has the same range after the substitution
\begin{align}
 T&=-t_{\mathrm{elastic}}=\frac{(v-1)^2}{v},
 &y_B&=H-T,\\
 b_e&=\frac{v^2+1}{v^2-1}\ln v,
 &\left|\frac{\dd c}{\dd v}\right|
 &=\frac{2(H+1)}{H^2}\left(1-\frac1{v^2}\right).
\end{align}
Here $y_B=1-u_{\mathrm{elastic}}$ distinguishes the elastic invariant
from the real-emission Dalitz variable. Let
$\mathcal B(v)=\overline{|\mathcal M_0|^2}/e_{\mathrm T}^4$ be the
Born square with its coupling factor removed, evaluated at $x=H$ and
$y=y_B$. The target density for $8H\mathcal C_{\mathrm{IR}}$
is consequently
\begin{equation}
 P(v)=\frac{16(H+1)}H\mathcal B(v)
 \left[\frac{v^2+1}{v^2}\ln v-1+\frac1{v^2}\right].
\end{equation}
Subtract this density from the sum of the $UV$ density and twice the
transformed $UW$ density.  The result is $L(v)\ln v+Q(v)$, with rational
$L$ and $Q$ fixed directly by the sector coefficient tables.  Integration by parts separates rational total derivatives from
simple-pole remainders. This operation, called rational Hermite reduction,
provides a short certificate for its integral.
Define
\begin{align}
 D(v)&=v^2-(H+2)v+1,\\
 p&=H^3-5H^2-28H-24,\qquad d=H^2+6H+4,\\
 j_0&=H^5+4H^4-12H^3+48H^2-32H-384,\\
 j_1&=2H^4+16H^3-40H^2+64H-192,\\
 j_2&=3H^4-16H^3-184H^2-736H-480.
\end{align}
Two rational functions suffice:
\begin{align}
 \mathcal R(v)&=\frac{H+1}{H^3}\left\{
 \frac{128H^2}{(v-1)^2}+16Hv^2
 +\frac{128H(H+2)(v-1)}{D(v)}
 +\frac{32H(3H-4)}{v-1}\right.\\
 &\hspace{14mm}\left.
 -\frac{16H}{v^2}-32dv
 +\frac{192(H+2)^2}{(v+1)^2}
 -\frac{128(H+2)^2}{(v+1)^3}
 +\frac{32p}{v+1}+\frac{32d}{v}\right\},\\
 \mathcal J(v)&=\frac{H+1}{3H^3}\left\{
 \frac{H^3(H+2)(H+6)(H^2-4H+8)[(H+2)v-1]}{D(v)^2}
 +\frac{H^2(j_0+j_1v)}{D(v)}\right.\\
 &\hspace{12mm}
 -\frac{384H^2}{v-1}-72Hv^2-\frac{72H}{v^2}
 -(6H^4-368H^2-1472H-960)v\\
 &\hspace{12mm}\left.
 -\frac{1152(H+2)^2}{v+1}
 +\frac{1152(H+2)^2}{(v+1)^2}-\frac{2j_2}{v}\right\}.
\end{align}
The four constant coefficients needed below are
\begin{align}
 C&=\frac{32(H+1)(H^2-4H-6)}{H^2},\\
 c_1&=\frac{2(H+1)}{3H^3}
 (3H^5-24H^4-112H^3+608H^2+2368H+1920),\\
 c_+&=\frac{256(H+1)p}{3H^3},\qquad
 c_D=-\frac{2(H+1)(H^3-8H^2-8H+48)}H.
\end{align}
Direct differentiation and rational cancellation give the two identities
\begin{align}
 L&=\mathcal R'+\frac Cv,\\
 Q-\frac{\mathcal R}{v}
 &=\mathcal J'+\frac{c_1}{v}+\frac{c_+}{v+1}
                      +c_D\frac{D'}D.
\end{align}
Thus an elementary primitive of the full difference is
\begin{equation}
 \Phi(v)=\mathcal R(v)\ln v+\frac C2\ln^2v+\mathcal J(v)
             +c_1\ln v+c_+\ln(v+1)+c_D\ln[-D(v)].
\end{equation}
There is no unresolved branch in this expression.  On $1\leq v\leq H+1$
the convex quadratic $D$ is negative, since its endpoint values both
 equal $-H$.  All displayed logarithm arguments are positive.
In particular, the $c_D$ logarithm has identical endpoint values.

The apparent pole at $v=1$ must be combined before taking its limit.
If the coefficients of $(v-1)^{-2}$ and $(v-1)^{-1}$ in $\mathcal R$
are $r_{-2}$ and $r_{-1}$, the constant from
$\mathcal R\ln v$ is $r_{-1}-r_{-2}/2$.  Adding the constant from
$\mathcal J$ gives
\begin{align}
 j_-&=-\frac{H+1}{3H^3}
       (H^5+24H^4-28H^3-624H^2-1936H-768),\\
 j_+&=\frac{1}{3H^3(H+1)}
 \left(H^9+6H^8-18H^7-36H^6+351H^5\right.\\
 &\hspace{35mm}\left.+2912H^4+6572H^3+6800H^2+3472H+768\right),\\
 r_+&=-\frac{16(H^5+20H^4+78H^3+114H^2+72H+16)}{H^3(H+1)}.
\end{align}
Here $j_+=\mathcal J(H+1)$ and $r_+=\mathcal R(H+1)$.
With $L_s=\ln(H+1)$ and $L_2=\ln[(H+2)/2]$, the integrated difference is
\begin{equation}
 E(H)=j_+-j_-+(r_++c_1)L_s+\frac C2L_s^2+c_+L_2.
\end{equation}
The remaining rational, single-logarithm, and Bose-restoration pole
expressions contain only the additional logarithms
\begin{equation}
 \ln\frac1{H+1}=-L_s,\qquad
 \ln\left(1+\frac{H^2}{4(H+1)}\right)=2L_2-L_s.
\end{equation}
Substitution and polynomial collection in the two independent symbols
$L_s,L_2$ gives zero for every rational coefficient in
$E+A_{\mathrm{rat}}+A_{\log}+A_D$.  Therefore
\begin{equation}
 A_{\mathrm{rat}}+A_{\log}+A_{UV}+2A_{UW}+A_D
              =8H\mathcal C_{\mathrm{IR}},\qquad H>0.
\end{equation}
The script \nolinkurl{compton_soft_uv_pole_certificate.py} verifies both
rational derivative identities, the endpoint constants, and the final
zero polynomial using exact arithmetic.  It also verifies that the
combined $UV$ and $UW$ pole densities start at order $v-1$, so their
lower endpoints are ordinary convergent integrals.  The machine-readable
certificate is \nolinkurl{soft_pole_certificate.json} in this supplement.


\section{Assembly in a common infrared convention}
All quantities in this section use $m=1$, so that $s$ means
$s_{\mathrm{phys}}/m^2$ and $\lambda$ means
$\lambda_{\mathrm{phys}}/m$. Put $H=s-1>0$. Define
\begin{equation}
 A(s)=\frac1s\int_{-1}^1\mathcal B(c)[b_e(c)-1]\,\dd c,
 \qquad \mathcal C_K(s)=\int_{-1}^1\mathcal B(c)K(s,c)\,\dd c.
\end{equation}
These are the integrated eikonal coefficient and the finite soft
angular integral derived above. Let $C_V$ be the complete virtual
contribution, including both external-photon orderings and the
on-shell subtractions, in units of $\alpha_{\mathrm T}^3/m^2$.
Its photon-mass term is $A(s)\ln\lambda$.

The relation to the local subtraction follows by angular integration.
For a future-directed timelike vector $v$, direct integration in polar
coordinates gives $\int\dd\Omega_1/(v\cdot n_1)^2=4\pi/v^2$.
Combining the two denominators in the cross term with a Feynman
parameter $\xi$ therefore gives
\begin{equation}
 \int\dd\Omega_1\,\frac{p\cdot p'_B}
 {(p\cdot n_1)(p'_B\cdot n_1)}
 =4\pi z_e\int_0^1\frac{\dd\xi}
 {1+2(z_e-1)\xi(1-\xi)}=4\pi b_e.
\end{equation}
In the last step, the substitution $2\xi-1$ gives the parameter integral
\begin{equation}
 \frac{2\operatorname{artanh}\sqrt{(z_e-1)/(z_e+1)}}{\sqrt{z_e^2-1}}
 =\frac{\operatorname{arcosh}z_e}{\sqrt{z_e^2-1}},
\end{equation}
with continuous value one at $z_e=1$.
The two diagonal terms each give $4\pi$, so
$\int\dd\Omega_1(-j_n^2)=8\pi(b_e-1)$. Using the factor
$q/\sqrt{s}=(s-1)/(2s)$ in $\mathcal A_{\mathrm{loc}}$ then yields
\begin{equation}
 \frac1{4\pi(s-1)}\int_{-1}^1\dd c_2\int\dd\Omega_1\,
 \mathcal A_{\mathrm{loc}}
 =\frac1s\int_{-1}^1\mathcal B(c_2)[b_e(c_2)-1]\,\dd c_2=A(s).
\end{equation}
Thus the local soft subtraction integrates to the same infrared
coefficient $A(1+X)$ used in Chapter 1.

The real integral regulated by $(xy)^\eta$ has the Laurent expansion
\begin{equation}
 C_R(\eta)=\frac{A(s)}\eta+F_R(s)+O(\eta).
 \label{eq:comp-assembly-real-laurent}
\end{equation}
This auxiliary regulator is used to evaluate the energy integrals.
Its finite part is converted to the same photon-mass convention as
the virtual result before either contribution is combined.
Near either soft corner the outgoing photon energy in the center-of-mass
frame is $\omega=e/(2\sqrt s)$, and $xy=He+O(e^2)$.
Consequently a real-energy cut $\omega>\varepsilon$ gives
\begin{equation}
 C_R^{\mathrm{hard}}(\varepsilon)
 =F_R-A\ln(2\sqrt s H\varepsilon)+o(1).
\end{equation}
The soft integral with photon mass $\lambda$ gives
$A\ln(2\varepsilon/\lambda)-\mathcal C_K/s$. Adding it first to the hard
integral removes $\varepsilon$. Adding the virtual term then gives
\begin{equation}
 C_{\mathrm{total}}
 =C_V(s,\lambda)+F_R(s)
   +A(s)\ln\frac1{\sqrt s H\lambda}-\frac{\mathcal C_K(s)}s.
 \label{eq:comp-assembly-common-regulator}
\end{equation}
The $\ln\lambda$ terms cancel algebraically. This argument fixes
both the coefficient of the regulator logarithm and the finite
scale factor $\sqrt s H$.

For the energy-sector decomposition used here, let
$F_{\mathrm{rat}}^z,F_{\log}^z,F_W^z,F_D^z$ denote finite parts
of the integrated rational, single-logarithm, $L_W$, and restoring
term $D$ with $z>\zeta$. Power and logarithmic divergences are
removed before taking $\zeta\to0$. Write $\Delta_j$ for the finite
conversion from this definition to $(xy)^\eta$. The $UV$ and $UW$
finite parts, $F_{UV}^{\eta}$ and $F_{UW}^{\eta}$, have already
been computed directly with the latter regulator. The required
assembly is
\begin{equation}
 \begin{aligned}
 F_R=\frac1{8sH}\big[&F_{\mathrm{rat}}^z+\Delta_{\mathrm{rat}}
       +F_{\log}^z+\Delta_{\log}+F_W^z\\
       &+F_D^z+\Delta_D+F_{UV}^{\eta}+2F_{UW}^{\eta}\big].
 \end{aligned}
 \label{eq:comp-assembly-real-sectors}
\end{equation}
Here $\Delta_W=0$. The term $D$ makes the replacement of the
subtracted $VW$ integral by a second $UW$ integral exact.
The factor $1/(8sH)$ includes the incident flux, the phase-space
normalization, and the identical-photon factor already established
in the real-emission derivation. It is applied once, after summing
the unnormalized energy integrals.

Every energy integral in \eqref{eq:comp-assembly-real-sectors}
is represented by a finite GPL primitive evaluated at its endpoints.
The endpoint conversions are logarithms and dilogarithms of the
explicit rational soft residues. Consequently the assembled answer
has no unevaluated phase-space integral:
\begin{equation}
 \boxed{\sigma_{\mathrm{inc}}^{\mathrm{NLO}}
 =\sigma_0+\frac{\alpha_{\mathrm T}^3}{m^2}
 C_{\mathrm{total}}(s_{\mathrm{phys}}/m^2)
 +\frac1{m^2}O(\alpha_{\mathrm T}^4).}
 \label{eq:comp-assembly-sigma}
\end{equation}
The remainder is at fixed dimensionless kinematics. The preceding
sector formulas and finite primitive recursions specify
$C_{\mathrm{total}}$ constructively, without using a published total
as an integration condition. Evaluating the GPLs by convergent
series is an evaluation of these analytic special functions, not
a numerical integration of the scattering phase space.


\subsection{Scope of the analytic reduction}
The complete virtual amplitude has been evaluated in closed scalar
functions. The vertex and self-energy angular integrals are given
explicitly in \eqref{eq:comp-vt-integrated}. The spanning and finite
soft angular integrals have finite multiple-polylogarithm endpoint
expressions, with explicit boundary constants. The coefficient
of the infrared logarithm cancels by the rational identity
\eqref{eq:comp-supp-box-IR-identity}. No published Compton
total-cross-section coefficient is used in these reductions.

The real-energy integrations reduce to finite GPL primitives, with
explicit endpoint constants and a common infrared definition.
Their assembly with the virtual and soft terms is
\eqref{eq:comp-assembly-sigma}. The separate finite pieces depend on
the subtraction convention, while their sum does not.
Independent numerical integration checks the result without
replacing the analytic primitives. The final verification chapter states the scope and estimated
numerical accuracy of these checks.
